% Options for packages loaded elsewhere
\PassOptionsToPackage{unicode}{hyperref}
\PassOptionsToPackage{hyphens}{url}
\documentclass[
  ukrainian,
]{article}
\usepackage{xcolor}
\usepackage[margin=2.5cm]{geometry}
\usepackage{amsmath,amssymb}
\setcounter{secnumdepth}{-\maxdimen} % remove section numbering
\usepackage{iftex}
\ifPDFTeX
  \usepackage[T1]{fontenc}
  \usepackage[utf8]{inputenc}
  \usepackage{textcomp} % provide euro and other symbols
\else % if luatex or xetex
  \usepackage{unicode-math} % this also loads fontspec
  \defaultfontfeatures{Scale=MatchLowercase}
  \defaultfontfeatures[\rmfamily]{Ligatures=TeX,Scale=1}
\fi
\usepackage{lmodern}
\ifPDFTeX\else
\setmainfont{DejaVu Serif}
\setsansfont{DejaVu Sans}
\setmonofont{DejaVu Sans Mono}
\setmathfont{STIX Two Math}
\fi
\ifPDFTeX\else
  % xetex/luatex font selection
\fi
% Use upquote if available, for straight quotes in verbatim environments
\IfFileExists{upquote.sty}{\usepackage{upquote}}{}
\IfFileExists{microtype.sty}{% use microtype if available
  \usepackage[]{microtype}
  \UseMicrotypeSet[protrusion]{basicmath} % disable protrusion for tt fonts
}{}
\makeatletter
\@ifundefined{KOMAClassName}{% if non-KOMA class
  \IfFileExists{parskip.sty}{%
    \usepackage{parskip}
  }{% else
    \setlength{\parindent}{0pt}
    \setlength{\parskip}{6pt plus 2pt minus 1pt}}
}{% if KOMA class
  \KOMAoptions{parskip=half}}
\makeatother
\usepackage{longtable,booktabs,array}
\newcounter{none} % for unnumbered tables
\usepackage{calc} % for calculating minipage widths
% Correct order of tables after \paragraph or \subparagraph
\usepackage{etoolbox}
\makeatletter
\patchcmd\longtable{\par}{\if@noskipsec\mbox{}\fi\par}{}{}
\makeatother
% Allow footnotes in longtable head/foot
\IfFileExists{footnotehyper.sty}{\usepackage{footnotehyper}}{\usepackage{footnote}}
\makesavenoteenv{longtable}
\usepackage{graphicx}
\graphicspath{{/Users/kalinicenkovitalijmikolajovic/.codex/worktrees/5c15/Script_for_translate_book/Output_Final/}}
\setkeys{Gin}{width=0.82\linewidth,height=0.82\textheight,keepaspectratio}
\makeatletter
\newsavebox\pandoc@box
\newcommand*\pandocbounded[1]{% scales image to fit in text height/width
  \sbox\pandoc@box{#1}%
  \Gscale@div\@tempa{\textheight}{\dimexpr\ht\pandoc@box+\dp\pandoc@box\relax}%
  \Gscale@div\@tempb{\linewidth}{\wd\pandoc@box}%
  \ifdim\@tempb\p@<\@tempa\p@\let\@tempa\@tempb\fi% select the smaller of both
  \ifdim\@tempa\p@<\p@\scalebox{\@tempa}{\usebox\pandoc@box}%
  \else\usebox{\pandoc@box}%
  \fi%
}
% Set default figure placement to htbp
\def\fps@figure{htbp}
\makeatother
\ifLuaTeX
\usepackage[bidi=basic,shorthands=off,provide=*]{babel}
\else
\usepackage[bidi=default,shorthands=off,provide=*]{babel}
\fi
\ifLuaTeX
  \usepackage{selnolig} % disable illegal ligatures
\fi
\setlength{\emergencystretch}{3em} % prevent overfull lines
\providecommand{\tightlist}{%
  \setlength{\itemsep}{0pt}\setlength{\parskip}{0pt}}
\usepackage{bookmark}
\IfFileExists{xurl.sty}{\usepackage{xurl}}{} % add URL line breaks if available
\urlstyle{same}
\hypersetup{
  pdflang={uk},
  hidelinks,
  pdfcreator={LaTeX via pandoc}}

\author{}
\date{}

\begin{document}

ЧАСТКОВО ДИФЕРЕНЦІАЛ РІВНЯННЯ для Вчені та Інженери

Стенлі Дж. Фарлоу, професор математики, Університет Мену

\section{С. Фарлоу}\label{ux441.-ux444ux430ux440ux43bux43eux443}

\section{РІВНЯННЯ З ПОХІДНИМИ ПОХІДНИМИ ДЛЯ НАУКОВЦІВ І
ІНЖЕНЕРІВ}\label{ux440ux456ux432ux43dux44fux43dux43dux44f-ux437-ux43fux43eux445ux456ux434ux43dux438ux43cux438-ux43fux43eux445ux456ux434ux43dux438ux43cux438-ux434ux43bux44f-ux43dux430ux443ux43aux43eux432ux446ux456ux432-ux456-ux456ux43dux436ux435ux43dux435ux440ux456ux432}

Перекладено з англійської \textbf{A.I.} Плісом

Редактор

С. І. Похожаєва

5\textbar{} 66K 22.161.6 φ24 υ ΥΠΚ 517.2

\section{Фарлоу С.}\label{ux444ux430ux440ux43bux43eux443-ux441.}

F 24 Рівняння з частими похідними для науковців та інженерів. Москва,
видавництво Мир, 1985, 384 с.

King's American Mathematician, який є підручником з теорії
диференціальних рівнянь у похідних похідних. Він характеризується
компактністю, ясністю та ясністю подачі, а також неформальним підходом
до подачі матеріалу. Вона містить багато ілюстрацій, графіків і діаграм,
але замість суворих доказів часто даються роздуми, засновані на інтуїції
або аналогіях.

Для інженерів і нематематиків --- біологів, диміків, а також студентів
університету.
\(\Phi = \frac{1702050000-380}{041(01)-85} 26-85, \text{ q. 1}\) BBK
22.161.6

Редакційна колегія журналу Literature on Mathematical Sciences

© переклад російською, Мир, 1985

© John Wiley \& Sons, Inc.~1982. Всі права захищені. Авторизований
переклад з англомовного видання, опублікований John Wiley \& Sons, Inc.

\section{З РЕДАКТОРА
ПЕРЕКЛАДУ}\label{ux437-ux440ux435ux434ux430ux43aux442ux43eux440ux430-ux43fux435ux440ux435ux43aux43bux430ux434ux443}

Книга, запропонована для читачів, слід розглядати як підручник для
курсів: «Диференціальні рівняння в похідних похідних» та «Рівняння
математичної фізики» для студентів технічних університетів. Вона також
буде корисною для інженерів різних спеціальностей. У книзі представлено
традиційні методи розв'язання диференціальних рівнянь у похідних
похідних, включаючи деякі чисельні методи. Водночас особлива увага
приділяється методу розділення змінних і методу інтегральних
перетворень, і цю частину книги можна використовувати на практичних
заняттях в університетах

Орієнтація книги на широку «нематематичну» аудиторію визначила стиль
викладання розглянутих питань: інтуїтивний підхід і велика увага до
фізичного значення не лише самих рівнянь, а й меж і початкових умов для
різних задач. Вважається, що читачі знайомі лише з основами
диференціального та інтегрального числення та елементами теорії
звичайних диференціальних рівнянь.

Фахівці можуть не задовольнитися інтуїтивним описом певних понять і
методів автором, але більш суворий виклад неминуче призведе до
ускладнення книги і звузить коло її читачів, що, очевидно, не було
наміром автора.

Більшим запереченням може бути відсутність загальних результатів щодо
задач для диференціальних рівнянь у похідних з змінними коефіцієнтами,
включно з рівняннями вищого порядку. Однак читач має чітко розуміти, що
ця книга є вступним курсом з теорії диференціальних рівнянь у часткових
похідних. Питання лінійної теорії диференціальних рівнянь, які автор не
включив, можна знайти в літературі, наведеній наприкінці книги.

Матеріал книги подається автором з великою методологічною майстерністю.
Однак не все було однаково успішним, і винний читач знайде це,
наприклад, у вивченні принципу Дюамеля, сферичних заповітів та деяких
інших уривків у книзі.

Без сумніву, запропонована книга С. Фарлоу знайде читача не лише серед
студентів і інженерів, а й серед викладачів, які можуть багато чого
навчитися з конкретного способу презентації та розробки спеціальних
математичних курсів.

\section{ПЕРЕДМОВА}\label{ux43fux435ux440ux435ux434ux43cux43eux432ux430}

За останні кілька років спостерігається значне зростання кількості
учнів, які починають вивчати диференціальні рівняння з похідними у своїх
одинадцятих класах. Крім того, багато з них спеціалізуються на сферах,
де інтуїція важливіша за математичну строгість. Тому, пишучи цю книгу, я
прагнув створити текст, який стимулює інтуїтивне мислення читачів, не
втрачаючи при цьому значної математичної строгості. Можливо, на одному
полюсі, подати предмет на високому рівні «епсилон-дельта», але тоді, як
це зазвичай буває, студенти не знатимуть, що робити з отриманою
інформацією. Інша крайність --- повне ігнорування математичних тонкощів,
але тоді ні учні, ні вчителі не знатимуть, як діяти. Я намагався знайти
розумний компроміс між цими крайнощами.

Ця книга виникла з курсу лекцій, які я читав протягом останніх п'яти
років. Розташування матеріалу не є загальноприйнятим --- вони майже
незалежні один від одного

Друг лекцій.

Основна увага приділяється двом найважливішим аналітичним методам:
розділенню змінних та інтегральним перетворенням. Також представлені
деякі нестандартні розділи, такі як метод Монте-Карло, варіаційне
числення, теорія керування, теорія потенціалів, інтегральні рівняння. Це
пов'язано з тим, що більшості студентів доведеться стикатися з цими
питаннями під час навчання, і якщо вони не ознайомляться з ними зараз,
навряд чи коли-небудь систематично їх вивчатимуть.

Цю книгу можна використовувати як навчальний посібник для одно- або
двосеместрового курсу як для початківців, так і для старших студентів.
Читачі повинні знати лише диференціальне та інтегральне числення та
звичайні диференціальні рівняння. Матеріал більшості лекцій подається в
одному-двох уроках. Типовий односеместровий курс може включати лекції
1-13, 15-17, 19-20, 22-23, 25-27, 30-32, 37-39. Усі 47 лекцій легко
читати протягом двох семестрів. Водночас часу буде достатньо для
вирішення проблем.

Автор вдячний Вайлі за пропозицію написати книгу, а також рецензентам
--- професорам Крісу Рорресу та М. Куршиду Алі, які зробили низку
корисних коментарів до тексту книги. Будь-які пропозиції щодо подальшого
покращення цього підручника будуть вдячні як від учнів, так і від
викладачів. Я також вдячний Дороті, Сюзан, Александер і Дезі Фарлоу.

Стенлі Дж. Фарлоу

\section{ВСТУП ДО ТЕОРІЇ ДИФЕРЕНЦІАЛЬНИХ РІВНЯНЬ У ПОХІДНИХ
ПОХІДНИХ}\label{ux432ux441ux442ux443ux43f-ux434ux43e-ux442ux435ux43eux440ux456ux457-ux434ux438ux444ux435ux440ux435ux43dux446ux456ux430ux43bux44cux43dux438ux445-ux440ux456ux432ux43dux44fux43dux44c-ux443-ux43fux43eux445ux456ux434ux43dux438ux445-ux43fux43eux445ux456ux434ux43dux438ux445}

МЕТА ЛЕКЦІЇ: Показати, що таке диференціальні рівняння в похідних
похідних, де і як вони виникають і як їх розв'язують. Коротко розгляньте
методи класифікації рівнянь і надайте список базових понять, які будуть
детально вивчені на наступних лекціях.

Більшість фізичних явищ у таких сферах, як гідродинаміка, електрика та
магнетизм, механіка, оптика та теплопередача, загалом можна описати за
допомогою рівнянь у похідних похідних (PNP). Більшість рівнянь у
математичній фізиці є диференціальними рівняннями з похідними похідними.
Правда, з деякими спрощеними припущеннями ці рівняння зводяться до
звичайних диференціальних рівнянь, але повний опис таких систем неминуче
призводить до використання диференціальних рівнянь у похідних похідних.

\section{Що таке часткові диференціальні
рівняння?}\label{ux449ux43e-ux442ux430ux43aux435-ux447ux430ux441ux442ux43aux43eux432ux456-ux434ux438ux444ux435ux440ux435ux43dux446ux456ux430ux43bux44cux43dux456-ux440ux456ux432ux43dux44fux43dux43dux44f}

Рівняння з частковими похідними --- це рівняння, яке містить
диференціали у похідних похідних. На відміну від звичайних
диференціальних рівнянь (TAC), де невідома функція залежить лише від
однієї змінної, у рівняннях часткових похідних вона залежить від кількох
змінних (наприклад, температура \(u(x, t)\) залежить від координати x і
часу t).

Тепер наведемо деякі з найважливіших диференціальних рівнянь у похідних
похідних. Для спрощення запису ми використаємо наступні нотації

\[
u_t \equiv \frac{\partial u}{\partial t}
\]

, \(u_x \equiv \frac{\partial u}{\partial x}\) ,
\(u_{xx} \equiv \frac{\partial^2 u}{\partial x^2}\) , \ldots{}

\section{Деякі диференціальні рівняння в похідних
похідних}\label{ux434ux435ux44fux43aux456-ux434ux438ux444ux435ux440ux435ux43dux446ux456ux430ux43bux44cux43dux456-ux440ux456ux432ux43dux44fux43dux43dux44f-ux432-ux43fux43eux445ux456ux434ux43dux438ux445-ux43fux43eux445ux456ux434ux43dux438ux445}

\[
u_t = u_{xx}
\]

(одновимірне рівняння теплопровідності), \(u_t = u_{xx} + u_{yy}\)
(двовимірне рівняння теплопровідності),
\(u_{rr} + \frac{1}{r}u_r + \frac{1}{r^2}u_{00} = 0\) (рівняння Лапласа
за полярними координатами), \(u_{tt} = u_{xx} + u_{yy} + u_{zz}\)
(тривимірне хвильове рівняння), \(u_{tt} = u_{xx} + zu_t + \beta u\)
(телеграфне рівняння).

\section{Примітки щодо наведених
прикладів}\label{ux43fux440ux438ux43cux456ux442ux43aux438-ux449ux43eux434ux43e-ux43dux430ux432ux435ux434ux435ux43dux438ux445-ux43fux440ux438ux43aux43bux430ux434ux456ux432}

У всіх наведених вище прикладах функція невідома і залежить від більше
ніж однієї змінної. Змінна і (яку ми диференціюємо) називається залежною
змінною. Змінні, за якими відбувається диференціювання, називаються
незалежними змінними. Наприклад, у рівнянні.

\[
u_t = u_{xx}
\]

Залежна змінна \(u(x, t)\) є функцією двох незалежних змінних x і t; та
в рівнянні

\[
u_t = u_{rr} + \frac{1}{r} u_r + \frac{1}{r^2} u_{00}
\]

змінна \(u(r, \theta, t)\) залежить від \(r, \theta, t\) .

\section{Чому необхідно вивчати диференціальні рівняння у похідних
похідних?}\label{ux447ux43eux43cux443-ux43dux435ux43eux431ux445ux456ux434ux43dux43e-ux432ux438ux432ux447ux430ux442ux438-ux434ux438ux444ux435ux440ux435ux43dux446ux456ux430ux43bux44cux43dux456-ux440ux456ux432ux43dux44fux43dux43dux44f-ux443-ux43fux43eux445ux456ux434ux43dux438ux445-ux43fux43eux445ux456ux434ux43dux438ux445}

Більшість фізичних законів природи можна сформулювати через рівняння з
похідними похідними. Прикладами є рівняння Максвелла, закон
теплопередачі Ньютона, рівняння Нав'є-Стокса, рівняння руху Ньютона та
рівняння Шредінгера в кватичній механіці. У всіх цих рівняннях фізичні
явища описуються мовою просторових і часових похідних. Похідні
з'являються в рівняннях, оскільки вони описують найважливіші фізичні
величини (такі як швидкість, прискорення, сила, тертя, потік, струм
тощо). Таким чином, виникають диференціальні рівняння в частинних
похідних, які містять невідому функцію, яку потрібно визначити.

Мета цієї книги --- показати читачеві:

\begin{enumerate}
\def\labelenumi{\arabic{enumi}.}
\item
  Як сформулювати фізичну задачу у формі диференціального рівняння у
  похідних похідних (побудова математичної моделі).
\item
  Як розв'язати диференціальне рівняння в похідних (з урахуванням
  початкових і граничних умов).
\end{enumerate}

Перш ніж перейти до побудови математичних моделей, давайте коротко
розглянемо методи розв'язання рівнянь у похідних похідних.

\section{Як розв'язати диференціальні рівняння в похідних
похідних?}\label{ux44fux43a-ux440ux43eux437ux432ux44fux437ux430ux442ux438-ux434ux438ux444ux435ux440ux435ux43dux446ux456ux430ux43bux44cux43dux456-ux440ux456ux432ux43dux44fux43dux43dux44f-ux432-ux43fux43eux445ux456ux434ux43dux438ux445-ux43fux43eux445ux456ux434ux43dux438ux445}

Гарне питання. Виявляється, існує цілий арсенал методів, придатних для
практичного використання. Найважливішими є ті, у яких диференціальні
рівняння з похідними похідними зводяться до звичайних диференціальних
рівнянь. Перелічимо десять методів розв'язання диференціальних рівнянь у
похідних похідних:

\begin{enumerate}
\def\labelenumi{\arabic{enumi}.}
\tightlist
\item
  Метод розділення змінних. Рівняння з факторами, отриманими з n
  незалежними змінними, зводиться до n звичайних-
\end{enumerate}

диференціальнірівняння 1).

\begin{enumerate}
\def\labelenumi{\arabic{enumi}.}
\setcounter{enumi}{1}
\item
  Метод інтегральних перетворень. Рівняння в похідних з n незалежними
  змінними зводиться до рівняння з похідними у (n-1) незалежних змінних;
  Відповідно, рівняння з частковими похідними з двома незалежними
  змінними можна звести до звичайного диференціального рівняння \(^2\)
  ).
\item
  Метод перетворення координат. Початкове диференціальне рівняння в
  похідних похідних зводиться до звичайного рівняння з частковими
  похідними або до іншого, простішого рівняння у похідних похідних
  шляхом відповідного перетворення координат (наприклад, обертання осей
  координат тощо).
\item
  Перетворення залежної змінної. Початкове диференціальне рівняння в
  похідних похідних перетворюється у таке рівняння з похідними похідними
  для іншої невідомої функції, яка
\end{enumerate}

легше розв'язати, ніж початкові 3).

\begin{enumerate}
\def\labelenumi{\arabic{enumi}.}
\setcounter{enumi}{4}
\tightlist
\item
  Чисельні методи. Початкове диференціальне рівняння в похідних похідних
  зводиться до системи різницей, яку розв'язують методом ітерацій на
  комп'ютері. У багатьох випадках це єдиний спосіб розв'язати
  диференціальне рівняння з похідними похідними. Окрім різницьких
  методів для розв'язання рівнянь з частковими про-
\end{enumerate}

\begin{enumerate}
\def\labelenumi{\alph{enumi})}
\setcounter{enumi}{4}
\tightlist
\item
  Ефективніше використовувати загальні перетворення незалежних і
  залежних змінних, включаючи похідні. Ед.
\end{enumerate}

4) Це стосується повного розділення змінних. При частковому розділенні
змінних одне диференціальне рівняння в похідних похідних зводиться до
кількох рівнянь у часткових похідних з меншою кількістю незалежних
змінних. --- Ед. Ед.

2) Це відбувається у випадку одномірного інтегрального перетворення. У
випадку k-вимірного інтегрального перетворення LCP з n незалежними
змінними зводиться до рівняння часткових з n---k незалежними змінними.
--- Ед. Ед.

Існують також інші чисельні методи, зокрема ті, що базуються на
наближенні розв'язків за поліноміальними поверхнями (наближення за
сплайнами).

\begin{enumerate}
\def\labelenumi{\arabic{enumi}.}
\setcounter{enumi}{5}
\tightlist
\item
  Метод теорії збурень. Початкова нелінійна задача зводиться до
  послідовності лінійних задач, апрокенн-
\end{enumerate}

нелінійна задача 1).

\begin{itemize}
\tightlist
\item
  \begin{enumerate}
  \def\labelenumi{\arabic{enumi}.}
  \setcounter{enumi}{6}
  \tightlist
  \item
    Метод функцій Гріна. Початкові та граничні умови замінюються
    системою найпростіших джерел, і задача розв'язується для кожного
    найпростішого джерела. Повне розв'язання початкової задачі
    отримується шляхом підсумовування розв'язків елементарних джерел.
  \end{enumerate}
\item
  \begin{enumerate}
  \def\labelenumi{\arabic{enumi}.}
  \setcounter{enumi}{7}
  \tightlist
  \item
    Метод інтегральних рівнянь. Рівняння з частковими похідними
    зводяться до інтегрального рівняння (рівняння, в якому невідома
    функція знаходиться під знаком інтеграла). Існує багато різних
    методів розв'язання інтегральних рівнянь.
  \end{enumerate}
\item
  \begin{enumerate}
  \def\labelenumi{\arabic{enumi}.}
  \setcounter{enumi}{8}
  \tightlist
  \item
    Варіаційні методи. Замість рівняння з частковими похідними
    диференціальними розв'язують задачу мінімізації. Виявляється, що
    функція, яка дає мінімум певному виразу (наприклад, загальній
    енергії системи), одночасно є розв'язком початкового рівняння
    похідниху похідних 2.
  \end{enumerate}
\item
  \begin{enumerate}
  \def\labelenumi{\arabic{enumi}.}
  \setcounter{enumi}{9}
  \tightlist
  \item
    Метод розкладу за власними функціями. Розв'язок рівняння часткових
    похідних шукається як послідовність власних функцій. Ці власні
    функції знаходяться як розв'язки так званої задачі власних значень,
    що відповідає початковій задачі для рівняння часткових похідних.
  \end{enumerate}
\end{itemize}

\section{Типи рівнянь у похідних
похідних}\label{ux442ux438ux43fux438-ux440ux456ux432ux43dux44fux43dux44c-ux443-ux43fux43eux445ux456ux434ux43dux438ux445-ux43fux43eux445ux456ux434ux43dux438ux445}

Рівняння в похідних похідних можна класифікувати за багатьма
характеристиками. Класифікація рівнянь важлива, оскільки, як виявилося,
кожен клас має свою загальну теорію та методи розв'язання рівнянь.

Тут ми наведемо шість основних методів класифікації

рівняння.

\begin{enumerate}
\def\labelenumi{\arabic{enumi}.}
\tightlist
\item
  Порядок рівняння. Порядок рівняння --- це найвищий порядок часткових
  похідних, включених до рівняння. Наприклад,
\end{enumerate}

Примітка. Ед.

1) Точніше, це метод лінеаризації нелінійних задач. Загальна теорія
збурень містить лінійні та нелінійні теорії збурень, кожна з яких, у
свою чергу, є регулярною та особливою. --- Ед. Ед.

Ці методи застосовні до спеціальних задач, що зумовлено існуванням
відповідних виразів --- функціоналів для цих задач.

Застосовувався раніше

\[
u_t = u_{xx}
\]

(рівняння другого порядку), \(u_t = u_x\) (рівняння першого порядку),
\(u_1 = uu_{xxx} + \sin x\) (рівняння трьох порядків).

\begin{enumerate}
\def\labelenumi{\arabic{enumi}.}
\setcounter{enumi}{1}
\tightlist
\item
  Кількість змінних. Кількість змінних --- це кількість незалежних
  змінних. Наприклад,
\end{enumerate}

\[
u_t = u_{xx}
\]

(рівняння з двома змінними \(x\) і \(t\) ),
\(u_t = u_{rr} + \frac{1}{r} u_r + \frac{1}{r^2} u_{\theta\theta}\)
(рівняння з трьома змінними \(r\) ,

\begin{enumerate}
\def\labelenumi{\arabic{enumi}.}
\setcounter{enumi}{2}
\tightlist
\item
  Лінійність. Рівняння з частковими похідними можуть бути лінійними та
  нелінійними. У лінійних рівняннях залежна змінна та всі її часткові
  похідні включаються лінійно, зокрема, вони не множаться одна на одну,
  не множаться в квадраті тощо). Точніше, лінійне рівняння другого
  порядку з двома незалежними змінними є рівнянням вигляду
\end{enumerate}

\[
(1.1) Au_{xx} + Bu_{xy} + Cu_{yy} + Du_x + Eu_y + Fu = G
\]

де A, B, C, D, E, F і G --- це константи або задані функції
пояснювальних змінних x і y.

\[
u_{tt} = e^{-t}u_{xx} + \sin t
\]

(лінійне рівняння), \(uu_{xx} + u_t = 0\) (нелінійне рівняння),
\(u_{xx} + yu_{yy} = 0\) (лінійне рівняння), \(xu_x + yu_y + u^2 = 0\)
(нелінійне рівняння).

\begin{itemize}
\tightlist
\item
  \begin{enumerate}
  \def\labelenumi{\arabic{enumi}.}
  \setcounter{enumi}{3}
  \tightlist
  \item
    Однорідність. Рівняння (1.1) називають однорідним, якщо права
    частина \(G(x, y)\) ідентична нулю для всіх x і y. Якщо \(G(x, y)\)
    не ідентичне нулю, то рівняння називають неоднорідним.
  \end{enumerate}
\item
  \begin{enumerate}
  \def\labelenumi{\arabic{enumi}.}
  \setcounter{enumi}{4}
  \tightlist
  \item
    Типи шансів. Якщо коефіцієнти A, B, C, D, E і F рівняння (1.1) є
    сталими, то це рівняння називається рівнянням зі сталими
    коефіцієнтами (інакше рівнянням зі змінними коефіцієнтами).
  \end{enumerate}
\item
  \begin{enumerate}
  \def\labelenumi{\arabic{enumi}.}
  \setcounter{enumi}{5}
  \tightlist
  \item
    Три основні типи лінійних рівнянь. Усі лінійні рівняння з частковими
    похідними другого порядку виду (1.1) належать до одного з трьох
    типів: а) параболічних, б) гіперболічних, в) еліптичних.
  \end{enumerate}
\end{itemize}

Параболічний тип. Рівняння параболічного типу описують процеси
теплопровідності та дифузії і визначаються умовою \(B^2 - 4AC = 0\) .

Гіперболічний тип. Рівняння гіперболічного типу описів

осцилюючі системи та хвильові рухи визначаються умовою \(B^2-4AC>0\) .

Еліпічний тип. Рівняння еліптичного типу описують усталені процеси і
визначаються умовним \(B^2-4AC<0\) . Приклади

\begin{enumerate}
\def\labelenumi{\alph{enumi})}
\tightlist
\item
\end{enumerate}

\[
u_i = u_{xx}
\]

, \(B^2 - 4AC = 0\) (параболічна), (гіперболічна), (гіперболічна),
(гіперболічна \(y > 0\) ), (гіперболічна), (еліптична), (еліптична),
параболічна на \(y = 0\) , гіперболічна на \(y < 0\) .

(У випадку змінних коефіцієнтів тип рівняння може змінюватися від точки
до точки.)

\section{ЗАУВАЖЕННЯ}\label{ux437ux430ux443ux432ux430ux436ux435ux43dux43dux44f}

\begin{itemize}
\tightlist
\item
  \begin{enumerate}
  \def\labelenumi{\arabic{enumi}.}
  \tightlist
  \item
    Загалом, значення \(B^2-4AC\) є функцією незалежних змінних.
    Відповідно, базовий тип рівняння може змінюватися в межах визначення
    рівняння, хоча це не є обов'язковим.
  \end{enumerate}
\item
  \begin{enumerate}
  \def\labelenumi{\arabic{enumi}.}
  \setcounter{enumi}{1}
  \tightlist
  \item
    У рівнянні (1.1) незалежними змінними є x і y. У багатьох задачах
    однією з двох незалежних змінних є час, і рівняння (1.1) можна
    записати через x і t.
  \end{enumerate}
\end{itemize}

\pandocbounded{\includegraphics[keepaspectratio]{images/_page_10_Figure_9.jpeg}}

PHC 1.1. Діаграма класифікації диференціальних рівнянь у похідних
похідних.

\begin{enumerate}
\def\labelenumi{\arabic{enumi}.}
\setcounter{enumi}{2}
\tightlist
\item
  Рис. 1.1. Представлено діаграму класифікації для диференціальних
  рівнянь у похідних похідних.
\end{enumerate}

\section{ЗАВДАННЯ}\label{ux437ux430ux432ux434ux430ux43dux43dux44f}

\begin{itemize}
\tightlist
\item
  \begin{enumerate}
  \def\labelenumi{\arabic{enumi}.}
  \tightlist
  \item
    Класифікуємо наступні рівняння відповідно до всіх характеристик,
    зазначених у днаграмі на рис. 1.1:
  \end{enumerate}
\item
  \begin{enumerate}
  \def\labelenumi{\alph{enumi})}
  \tightlist
  \item
    \(u_t = u_{xx} + 2u_x + u\) ,
  \end{enumerate}
\end{itemize}

\begin{enumerate}
\def\labelenumi{\arabic{enumi})}
\setcounter{enumi}{5}
\tightlist
\item
  \(u_t = u_{xx} + e^{-t}\) , B) \(u_{xx} + 3u_{xy} + u_{yy} = \sin x\)
  , \(u_{tt} = uu_{xxx} + e^{-t}\) .
\end{enumerate}

\begin{itemize}
\tightlist
\item
  \begin{enumerate}
  \def\labelenumi{\arabic{enumi}.}
  \setcounter{enumi}{1}
  \tightlist
  \item
    Скільки розв'язків \(u_1 = u_{xx}\) рівняння існує? Спробуйте знайти
    рішення, як \(u(x, t) = e^{a\lambda + bt}\) .
  \end{enumerate}
\item
  \begin{enumerate}
  \def\labelenumi{\arabic{enumi}.}
  \setcounter{enumi}{2}
  \tightlist
  \item
    Якщо функції \(u_1(x, t)\) і \(u_2(x, t)\) задовольняють рівняння
    (1.1), чи задовольняє сума цих функцій його? Доведи це.
  \end{enumerate}
\item
  \begin{enumerate}
  \def\labelenumi{\arabic{enumi}.}
  \setcounter{enumi}{3}
  \tightlist
  \item
    Ймовірно, найпростіше диференціальне рівняння в похідних ---
  \end{enumerate}
\end{itemize}

\[
\frac{\partial u(x, y)}{\partial x} = 0.
\]

Чи можете ви розв'язати це рівняння? (Знайдіть усі функції, що
задовольняють це рівняння.)

\begin{enumerate}
\def\labelenumi{\arabic{enumi}.}
\setcounter{enumi}{4}
\tightlist
\item
  Що можна сказати про рівняння
\end{enumerate}

\[
\frac{\partial^2 u(x, y)}{\partial x \partial y} = 0
\]

?

Чи можете ви знайти всі розв'язки цього рівняння? Скільки їх? Порівняйте
з кількістю розв'язків звичайного диференціального рівняння

\[
\frac{d^2y}{dx^2} = 0.
\]

\section{ЗАДАЧІ ТИПУ ДИФУЗІЇ (ПАРАБОЛІЧНІ
РІВНЯННЯ)}\label{ux437ux430ux434ux430ux447ux456-ux442ux438ux43fux443-ux434ux438ux444ux443ux437ux456ux457-ux43fux430ux440ux430ux431ux43eux43bux456ux447ux43dux456-ux440ux456ux432ux43dux44fux43dux43dux44f}

МЕТА ЛЕКЦІЇ: Показати, як параболічні рівняння використовуються для
розв'язання задач теплопровідності та дифузії. Наведено кілька прикладів
параболічних рівнянь, а також пояснюється фізичне значення різних членів
\((u_1, u_2, u_3, u_4, u_6)\) , включених у ці рівняння.

На конкретному прикладі розглядається основна ідея зміщеної задачі для
параболічного рівняння. Одна з головних цілей цієї лекції --- надати
читачеві інтуїтивне уявлення

Розуміння параболічних проблем.

Почнемо з розгляньмо просту фізичну задачу і те, як її можна описати
математичною моделлю з рівнянням у похідних похідних.

Потім ми ускладнимо задачу і покажемо, як нові фізичні умови призводять
до появи нових рівнянь у похідних похідних. У цій лекції диференціальні
рівняння не виводяться і не розв'язуються. Про це ми поговоримо пізніше.

\section{Простий експеримент з
теплопровідністю}\label{ux43fux440ux43eux441ux442ux438ux439-ux435ux43aux441ux43fux435ux440ux438ux43cux435ux43dux442-ux437-ux442ux435ux43fux43bux43eux43fux440ux43eux432ux456ux434ux43dux456ux441ux442ux44e}

Припустимо, що ми провели простий експеримент, який розбили на такі
кроки:

КРОК 1. Візьміть достатньо довгий (скажімо, \(L=2\,\mathrm{m}\) )
стрижень (наприклад, мідний) діаметром \(2\,\mathrm{cm}\) , бічна
поверхня (але не кінці) якої термічно ізольована. Замість стрижня можна
використовувати мідну трубку, покриту теплоізоляцією зовні та зсередини.

Іншими словами, тепло може проходити лише через кінці стрижня і не може
проходити через його бічну поверхню.

КРОК 2. Тепер помістимо цей стрижень у пристрій із фіксованою
температурою \(T_{\mathfrak{o}}\) (градуси Цельсія) на достатньо довгий
час, щоб температура всередині стрижня стала

Так само, як і температура пристрою. Для простоти припустимо, що
температура пристрою \(T_0 = 10^{\circ}\) C.

КРОК 3. Давайте витягнемо стрижень з нашого пристрою у певний момент
часу, що зручно взяти як нуль t = 0, і приєднаємо до нього дві термопари
з кінців. Завдання цих елементів --- підтримувати фіксовані температури
\(T_1\) і \(T_2\) на кінцях (наприклад, \(T_1=0^{\circ}\) C і
\(T_2=50^{\circ}\) C). Іншими словами, два термостати постійно
контролюють температуру на кінцях стрижня, і якщо температура
відхиляється від встановлених значень \(T_1\) і \(T_2\) , активуються
потужні нагрівальні або охолоджувальні елементи (пристрої) для
відповідного регулювання температури. Наш експеримент ілюструє RPS, 2 1,

\pandocbounded{\includegraphics[keepaspectratio]{images/_page_13_Figure_4.jpeg}}

Рис. 2.1. Схема експерименту: a --- це термопара, яка підтримує
встановлену температуру \(T_1\) на лівому кінці стрижня; \(\delta\) ---
це термопара, яка підтримує встановлену температуру \(T_2\) на правому
кінці стрижня.

КРОК 4. Тепер ми будемо слідкувати за температурним профілем у стрижні
за допомогою пристрою візуалізації. Чому ми хочемо провести такий
експеримент? Це інше питання, про яке ми поговоримо пізніше.

Це кінець опису нашого експерименту.

Головна мета цієї лекції --- показати, як цю фізичну проблему та її
модифікації можна змоделювати мовою параболічних рівнянь.

\section{Математична модель
теплопровідності}\label{ux43cux430ux442ux435ux43cux430ux442ux438ux447ux43dux430-ux43cux43eux434ux435ux43bux44c-ux442ux435ux43fux43bux43eux43fux440ux43eux432ux456ux434ux43dux43eux441ux442ux456}

Опис нашої фізичної проблеми вимагає трьох типів відносин.

\begin{itemize}
\tightlist
\item
  \begin{enumerate}
  \def\labelenumi{\arabic{enumi}.}
  \tightlist
  \item
    Рівняння з частковими похідними, які описують фізичне явище
    теплопровідності.
  \end{enumerate}
\item
  \begin{enumerate}
  \def\labelenumi{\arabic{enumi}.}
  \setcounter{enumi}{1}
  \tightlist
  \item
    Граничні умови, що описують процеси теплопередачі на межах.
  \end{enumerate}
\item
  \begin{enumerate}
  \def\labelenumi{\arabic{enumi}.}
  \setcounter{enumi}{2}
  \tightlist
  \item
    Початкові твердження, які описують стан системи на початку процесу.
  \end{enumerate}
\end{itemize}

\section{Рівняння
теплопровідності}\label{ux440ux456ux432ux43dux44fux43dux43dux44f-ux442ux435ux43fux43bux43eux43fux440ux43eux432ux456ux434ux43dux43eux441ux442ux456}

Базове одновимірне рівняння теплопровідності записується у вилці

(2.1)

\[
u_t = \alpha^2 u_{xx}, \quad 0 < x < L, \quad 0 < t < \infty.
\]

Це рівняння пов'язує значення \(u_t\) --- швидкість зміни температури з
часом (виміряна в градусах/с)

та величину \(u_{xx}\) --- це увігнутий профілю температури \(u(x, t)\)
(який є мірою різниці між температурою в заданій точці та температурою в
сусідніх точках).

\pandocbounded{\includegraphics[keepaspectratio]{images/_page_14_Figure_14.jpeg}}

Рис. 2.2. Стрілки показують зміну температури відповідно до рівняння
\(u_t = \alpha^2 u_{xx}\) .

Це рівняння буде виведено з базового рівняння збереження тепла в
наступних лекціях, але наразі ми розглянемо його заданим. Це рівняння
каже, що температура \(u\left(x,\ t\right)\) (у певній точці t у певній
точці x стрижня) зростає \((u_t>0)\) або зменшується \((u_t<0)\) залежно
від того, чи є друга похідна \(u_{xx}\) позитивною чи негативною. Фіг.
Рисунок 2.2 ілюструє зміну температури в різних точках на елементі.

Давайте подивимось, як значення \(u_{xx}\) можна інтерпретувати мовою
теплопровідності. Припустимо, ми наблизимо значення \(u_{xx}\) фінальних
постів

\[
u_{xx}(x, t) \cong \frac{1}{\Delta x^2} [u(x + \Delta x, t) - 2u(x, t) + u(x - \Delta x, t)].
\]

Це співвідношення можна переписати у вигляді

\[
u_{\lambda x}(x, t) \simeq \frac{2}{\Delta x^2} \left[ \frac{u(x + \Delta x, t) + u(x - \Delta x, t)}{2} - u(x, t) \right].
\]

Тепер ми можемо дати наступну інтерпретацію значення \(u_{xx}\) .

\begin{enumerate}
\def\labelenumi{\arabic{enumi}.}
\item
  Якщо температура \(u(x, t)\) менша за середнє значення температури у
  двох сусідніх точках, то \(u_{xx} > 0\) (тут загальний тепловий потік
  по осі x додатний).
\item
  Якщо температура \(u(x, t)\) дорівнює середньому значенню температур у
  двох сусідніх точках, то \(u_{xx} = 0\) (тут загальний тепловий потік
  по осі x дорівнює нулю).
\item
  Якщо температура \(u\left(x,\ t\right)\) вища за середню температуру у
  двох сусідніх точках, то \(u_{xx}<0\) (тут --- повний потік
\end{enumerate}

тепла вздовж осі x є від'ємною).

Усе це ілюструється на рис. 2.2. Іншими словами, якщо температура в
точці x більша за середню температуру у двох сусідніх точках
\(x - \Delta x\) і \(x + \Delta x\) , то температура в точці x
зменшиться. Відповідно, точна швидкість зниження температури, а саме
величина \(u_t\) , пропорційна цій різниці. Коефіцієнт пропорційності
\(\alpha^2\) визначається властивостями матеріалу. Ми розглянемо цю
константу детальніше у майбутніх лекціях.

\section{Граничні умови
(ГІ)}\label{ux433ux440ux430ux43dux438ux447ux43dux456-ux443ux43cux43eux432ux438-ux433ux456}

Усі фізичні проблеми характеризуються наявністю певних меж. Отже, щоб
рухатися вперед, ми повинні включити межі в математичну модель, щоб
адекватно описати фізичну проблему. У нашому експерименті граничні умови
(ГІ) дуже прості. Оскільки температура на кінцях x=0 і x=L завжди
залишається сталою і дорівнює \(T_1\) та \(T_2\) , можливо записати

(2.2)

\begin{verbatim}
\left\{ \begin{array}{l} u\left(0,\ t\right) = T_{t}, \\ u\left(L,\ t\right) = T_{2}, \end{array} \right. \quad 0 < t < \infty,
\end{verbatim}

\section{Початкові умови
(NU)}\label{ux43fux43eux447ux430ux442ux43aux43eux432ux456-ux443ux43cux43eux432ux438-nu}

Усі фізичні процеси мають починатися у певний момент часу (зазвичай
приймається як нуль t=0), тому ми повинні встановити фізичні умови саме
в цю точку. Відтоді, як ми почали контролювати температуру з моменту,
коли паличка мала постійну температуру \(T_{\mathrm{ }o}\) , ми

(2.3)

\[
u(x, 0) = T_0, \quad 0 \le x \le L
\]

Тепер ми створили математичну модель експерименту. Система співвідношень
(2.1), (2.2) та (2.3) називається задачею зміщення для рівняння
теплопровідності і зазвичай записується як

(2.4)

\[
\begin{aligned} & (\text{YUII}) \quad u_t = \alpha^2 u_{xx}, & 0 < x < L, \ 0 < t < \infty, \\ & \left( \Gamma \text{Y} \right) \quad \begin{cases} u \ (0, \ t) = T_1, \\ u \ (L, \ t) = T_2, \end{cases} & 0 < t < \infty, \\ & (\text{HY}) \quad u \ (x, \ 0) = T_0, \quad 0 \leqslant x \leqslant L. \end{aligned}
\]

Цікаво, що це зовсім не очевидно: чи існує одна функція \(u(x, t)\), яка
задовольняє умови (2.4), і чи описує ця функція температуру в стрижні?
Отже, наше безпосереднє завдання --- знайти цю функцію \(u(x, t)\) з
(2.4).

Перш ніж завершити цю лекцію, давайте розглянемо деякі варіації основної
проблеми. Почнемо з невеликих змін рівняння теплопровідності.

\section{Деякі рівняння дифузійного
типу}\label{ux434ux435ux44fux43aux456-ux440ux456ux432ux43dux44fux43dux43dux44f-ux434ux438ux444ux443ux437ux456ux439ux43dux43eux433ux43e-ux442ux438ux43fux443}

Теплопередача через бічну поверхню пропорційна різниці температур

Рівняння

\[
u_t = \alpha^2 u_{xx} - \beta (u - u_0)
\]

Описує теплопровідність елемента, враховуючи не лише дифузію теплових
\(\alpha^2 u_{xx}\) вздовж елемента, а й теплопередачу через бічну
поверхню елемента. Відтік тепла ( \(\beta > 0\) ) або його приплив (
\(\beta < 0\) ) пропорційний різниці між температурою u стрижня та
температурою навколишнього середовища \(u_0\) ( \(\beta\) --- це сталий
коефіцієнт пропорційності). Якщо коефіцієнт \(\beta\) великий порівняно
з \(\alpha^2\) , то теплопотік уздовж стрижня буде малим порівняно з

Протікання через бічну поверхню, і тому тепло буде проходити через неї
відповідно до закону \(u_t = -\beta (u - u_a)\) .

У хімії, де величина відіграє роль концентрації, рівняння

\[
u_1 = \alpha^2 u_{xx} - \beta (u - u_0)
\]

означає, що швидкість зміни \(u_i\) кількості речовини залежить як від
дифузії \(\alpha^2 u_{xx}\) (у напрямку осі x), так і від появи (
\(\beta < 0\) ) або розпаду ( \(\beta > 0\) ) речовини в хімічній
реакції, і пропорційна різниці між двома концентраціями u і
\(u_{\theta}\) .

\section{Внутрішнє джерело
тепла}\label{ux432ux43dux443ux442ux440ux456ux448ux43dux454-ux434ux436ux435ux440ux435ux43bux43e-ux442ux435ux43fux43bux430}

Неоднородне рівняння

\[
u_t = \alpha^2 u_{xx} + f(x, t)
\]

відповідає випадку, коли всередині елемента є джерело тепла (розташоване
вздовж усього елемента і діє в будь-який момент t). Наприклад, може бути
так, що дріт з електричним струмом проходить через стрижень і опір є
постійним джерелом тепла \(f(x,t)\) = K.

\section{Рівняння конвективної
дифузії}\label{ux440ux456ux432ux43dux44fux43dux43dux44f-ux43aux43eux43dux432ux435ux43aux442ux438ux432ux43dux43eux457-ux434ux438ux444ux443ux437ux456ux457}

Припустимо, що домішка поширюється вздовж потоку, що рухається зі
швидкістю v. Очевидно, що концентрація \(u\left(x,\,t\right)\) домішок
змінюється залежно від x (позитивний напрямок осі x обирається вздовж
потоку) і часу t. Швидкість зміни концентрації \(u_t\) визначається
рівнянням конвективної дифузії

\[
u_t = \alpha^2 u_{xx} - v u_x.
\]

Термін \(\alpha^2 u_{xx}\) описує внесок дифузії, а \(vu_x\) ---
конвективний компонент. Що важливіше, дифузія або конвекція залежить від
співвідношення \(\alpha^2\) до v. Димові частинки піднімаються
конвективно через рух гарячого повітря і одночасно розсіюють через
вихорові рухи повітря. Окрім цих модифікацій рівняння теплопровідності,
можливо змінювати граничні умови на кінцях стрижня відповідно до
фізичної ситуації. Деякі з цих змін ми обговоримо у третій лекції.

\section{ЗАУВАЖЕННЯ}\label{ux437ux430ux443ux432ux430ux436ux435ux43dux43dux44f-1}

Рівняння теплопровідності \(u_t = \alpha^2(x) u_{xx}\) з змінним
коефіцієнтом \(\alpha(x)\) відповідає задачі, в якій дифузія залежить
від x (тобто матеріал є гетерогенним). Наприклад, мідь і сталь---

\pandocbounded{\includegraphics[keepaspectratio]{images/_page_18_Picture_4.jpeg}}

PHC. 2.3.

Пластини контактують, як показано на рис. 2.3. Припустимо, що
температура \(u(0,t)=0^{\circ}\) C підтримується на лівому боці міді, а
температура \(u(L,t)=20^{\circ}\) C --- з правого боку сталі, тоді
рівняння, що описує теплопровідність у цій системі, виглядає так:

\[
u_t = \alpha^2(x) \, u_{_{\! A \, X}}, \quad 0 < x < L,
\]

де

\[
\alpha(x) = \begin{cases}
\alpha_1, & \text{коефіцієнт дифузії в міді}, \quad 0 < x < L/2, \\
\alpha_2, & \text{коефіцієнт дифузії в сталі}, \quad L/2 < x < L.
\end{cases}
\]

\section{ЗАВДАННЯ}\label{ux437ux430ux432ux434ux430ux43dux43dux44f-1}

\begin{enumerate}
\def\labelenumi{\arabic{enumi}.}
\tightlist
\item
  Якщо початкова температура в стрижні
\end{enumerate}

\[
u(x, 0) = \sin \pi x, \quad 0 \leqslant x \leqslant 1,
\]

а граничні умови такі:

\[
u(0, t) = 0, \qquad u(1, t) = 0.
\]

Тоді як поводитиметься температура в стрижні \(U(x, t)\) при \(t > 0\) ?
2. Припустимо, що в стрижні є постійне внутрішнє джерело тепла, так що
рівняння теплопровідності в стрижні записується як

\[
u_t = \alpha^2 u_{xx} + 1, \quad 0 < x < 1.
\]

Припустимо, що ми зафіксували температуру на межі з значеннями
\(u(0, t)\) = 0 і \(u(1, t)\) = 1. Якою буде стаціонарна температура у
стрижні? Іншими словами, який c-незалежний розподіл \(U(x, t)\)
перешкоджує \(u(x, t)\)?

\begin{enumerate}
\def\labelenumi{\arabic{enumi}.}
\setcounter{enumi}{2}
\tightlist
\item
  Припустимо, що металевий стрижень втрачає тепло через бічну поверхню
  згідно з рівнянням
\end{enumerate}

\[
u_1 = \alpha^2 u_{xx} - \beta u
\]

, \(0 < x < 1\) .

І припустимо, що ми підтримуємо кінці стрижня при температурах
\(U(0, T)\) = 1, \(U(1, T)\) = 1. Знайдіть стаціонарний розподіл
температури по стрижню. Побудуйте графік цього розподілу. У яких
моментах відбувається тепловий потік?

\begin{enumerate}
\def\labelenumi{\arabic{enumi}.}
\setcounter{enumi}{3}
\tightlist
\item
  Припустимо, що термоізольований металевий стрижень довжини L = 1 має
  початкову температуру \(\sin(3\pi x)\) , а фіксовану температуру --- 0
  і \(10\,^{\circ}\) C на лівому та правому кінцях.
\end{enumerate}

\section{Лекція 3}\label{ux43bux435ux43aux446ux456ux44f-3}

\section{ГРАНИЧНІ УМОВИ В ЗАДАЧАХ ДИФУЗІЙНОГО
ТИПУ}\label{ux433ux440ux430ux43dux438ux447ux43dux456-ux443ux43cux43eux432ux438-ux432-ux437ux430ux434ux430ux447ux430ux445-ux434ux438ux444ux443ux437ux456ux439ux43dux43eux433ux43e-ux442ux438ux43fux443}

МЕТА ЛЕКЦІЇ: Показати, як виникають різні граничні умови у задачах
теплопровідності та дифузії, а також ввести важливе поняття теплового
потоку.

Розглядаються три найважливіші типи граничних умов:

\begin{itemize}
\tightlist
\item
  \begin{enumerate}
  \def\labelenumi{\arabic{enumi}.}
  \tightlist
  \item
    u = \(g(t)\) (температура задається на межі),
  \end{enumerate}
\item
  \texttt{2.\ \textbackslash{}\ \textbackslash{}frac\{\textbackslash{}partial\ u\}\{\textbackslash{}partial\ n\}\ +\ \textbackslash{}lambda\ u\ =\ g\textbackslash{}left(t\textbackslash{}right)}
  (встановлення температури навколишнього середовища);
\end{itemize}

n --- вектор зовнішньої нормалі до межі),
\texttt{3.\ \textbackslash{}frac\{\textbackslash{}partial\ u\}\{\textbackslash{}partial\ n\}\ =\ g(t)}
(визначений тепловий потік через межу).

У задачах теплопровідності зазвичай використовують три основні типи
граничних умов. Ця лекція наводить приклади того, як різні фізичні умови
експерименту призводять до різних типів граничних умов.

\section{Граничні умови першого виду (температура задається на
межі)}\label{ux433ux440ux430ux43dux438ux447ux43dux456-ux443ux43cux43eux432ux438-ux43fux435ux440ux448ux43eux433ux43e-ux432ux438ux434ux443-ux442ux435ux43cux43fux435ux440ux430ux442ux443ux440ux430-ux437ux430ux434ux430ux454ux442ux44cux441ux44f-ux43dux430-ux43cux435ux436ux456}

Розглянемо тепловий потік в одномірному стрижні, показаний на рис. 3 1,
і припустимо, що температура кінців змінюється відповідно до законів
\(g_1(t)\) і \(g_2(t)\) .

\pandocbounded{\includegraphics[keepaspectratio]{images/_page_20_Picture_4.jpeg}}

Фіг. 31. Температура буде відновлена на межі.

Як згадувалося в попередній лекції, кожен пристрій, що підтримує
встановлену температуру на кінцях, складається з термостата та
нагрівального елемента для підтримки відповідних теплових умов. Проблеми
з граничними умовами першого виду дуже поширені. У деяких випадках
складність полягає у пошуку контролерів граничних умов, тобто температур
\(g_1(t)\) та \(g_2(t)\) меж, які призводять до певної зміни температури
всередині стрижня. Наприклад, у металургії часто необхідно обирати
граничні умови контролю, щоб температура металу всередині печі
змінювалася з часом, а температурний градієнт був невеликим.

\pandocbounded{\includegraphics[keepaspectratio]{images/_page_20_Picture_7.jpeg}}

Фіг. 3.2. Зниження температури на межі.

Граничні умови першого виду також виникають у багатовимірних задачах
теплопровідності. Як приклад можна навести цікаву задачу визначення
температурного поля всередині кругового диска з радіусом \emph{R} (див.
рис. 3.2), коли гранична температура задається у полярних координатах
формулою \(u(R, \theta, t) = \cos t \sin \theta\) .

Звісно, щоб розв'язати задачу, потрібно знати початкову температуру, але
в цьому випадку початкові умови перестають впливати на розв'язок через
дуже короткий час, і отримана температура всередині диска визначається
лише значеннями температури на межі.

\section{Граничні умови другого виду (задана температура навколишнього
середовища)}\label{ux433ux440ux430ux43dux438ux447ux43dux456-ux443ux43cux43eux432ux438-ux434ux440ux443ux433ux43eux433ux43e-ux432ux438ux434ux443-ux437ux430ux434ux430ux43dux430-ux442ux435ux43cux43fux435ux440ux430ux442ux443ux440ux430-ux43dux430ux432ux43aux43eux43bux438ux448ux43dux44cux43eux433ux43e-ux441ux435ux440ux435ux434ux43eux432ux438ux449ux430}

Припустимо, що ми знову розглядаємо термічно ізольований мідний
стержень, але тепер відмовляємося від заданого температурного режиму на
кінцях стрижня. Натомість давайте наблизимо кінці стрижня до контакту з
двома середовищами. Нехай температура одного з них \(g_1(t)\) , а іншого
\(g_2(t)\) . Іншими словами, припустимо, що лівий кінець стрижня
поміщений у контейнер із рідиною, температура якої змінюється відповідно
до закону \(g_1(t)\) . А правий кінець поміщається в іншу рідину з
температурою \(g_2(t)\) (рис. 3, 3)

\pandocbounded{\includegraphics[keepaspectratio]{images/_page_21_Picture_6.jpeg}}

Рис. 3.3. Конвективна теплопередача через межі, a --- це рідина при
температурі \(g_1(t)\) , \(\delta\) --- рідина при температурі
\(g_2(t)\) .

Задаючи граничні умови такого типу, ми не можемо вважати, що граничні
температури стрижня збігаються з температурами рідини. Але ми знаємо
закон Ньютона: якщо температура одного з кінців стрижня менша за
температуру відповідної рідини, тепло потече в стрижень зі швидкістю,
пропорційною різницям температур. Іншими словами, для одномірного
стрижня з межами x=0 і x=1 закон теплопередачі Ньютона формулюється так:

(3.1)

\begin{verbatim}
\{
\end{verbatim}

вихідний тепловий потік на \(x=0\) дорівнює
\(h\left[u\left(0,\,t\right)-g_{1}\left(t\right)\right]\) , вихідний
тепловий потік при \(x=L\) дорівнює
\(h\left[u\left(L,\,t\right)-g_{2}\left(t\right)\right]\) ,

де h --- коефіцієнт теплопередачі, який показує, скільки калорій
проходить через межу за одну секунду при різниці температур у один
градус. Вихідний потік дорівнює кількості кало-

Пройшов крізь палочний спис за одну секунду. Зверніть увагу, що вихідний
тепловий потік буде позитивним на кінці стрижня, де температура стрижня
вища за температуру навколишнього середовища. Рівняння (3.1) разом із
відомим законом теплопровідності Фур'є тепер можна використовувати для
отримання Р и с

\pandocbounded{\includegraphics[keepaspectratio]{images/_page_22_Figure_3.jpeg}}

. 3.4. Ілюстрація закону Фур'є: \(\frac{\partial u}{\partial n} > 0\)
означає, що тепло надходить у цю область;
\(\frac{\partial u}{\partial n} < 0\) означає, що тепло виходить із цієї
зони; n --- напрямок зовнішньої нормалі;
\(\frac{\partial u}{\partial n} -\) Зміна температури в напрямку N;
\(\frac{\partial u}{\partial n} = -\frac{\partial u}{\partial (-n)}\) є
похідною у напрямку внутрішньої нормальної --- це похідна, взята
\(\mathbf{e}\) протилежному знаку протилежного знака пормалі.

граничні умови. Закон Фур'є дає друге представлення (перше міститься у
(3.1)) для вихідного теплового потоку. Прирівнюючи ці два вирази до
потоку, ми отримуємо граничні умови, які шукаємо. Сформулюємо закон
Фур'є (встановлений експериментально):

(3.2) Тепловий потік, що проходить через межу площі, пропорційний
нормальній похідній температури у напрямку внутрішньої норми.

Цей закон стверджує, що якщо температура швидко підвищується у напрямку
зовнішньої температури до межі D (рисунок 3.4), то потік буде текти з
навколишнього середовища до області D.

У нашій одновимірній задачі закон Фур'є має вигляд:

(3.3)

\[
\begin{cases} \text{вытекающий поток тепла при } x=0 \text{ равен } k \frac{\partial u}{\partial x}, \\ \text{вытекающий поток тепла при } x=' \text{ равен } -k \frac{\partial u}{\partial x}, \end{cases}
\]

де k --- теплопровідність матеріалу, яка є мірою того, наскільки добре
матеріал проводить тепло. Погано провідні матеріали мають коефіцієнт
близький до нуля в одиницях GHS, тоді як мідь та алюміній мають значення
значно вище нуля.

\pandocbounded{\includegraphics[keepaspectratio]{images/_page_23_Figure_5.jpeg}}

Фіг. 3.5. Ще одна ілюстрація закону Чотирьох.

Закон Фур'є (3.3) фактично описує теплопровідність всюди всередині
стрижня, а не лише на його межі; Наприклад, (див. рис. 3.5).

(3.4) Потік, що проходить через точку

x

зліва направо \(=-k\frac{\partial n}{\partial x}\) Закон Фур'є (3.4)
стверджує: якщо \(u_x < 0\) , то в точці \(x_0\) тепло тече зліва
направо, якщо \(u_x > 0\) , то в точці \(x_0\) тепло тече справа наліво.
(Тепло завжди переходить від вищих температур до нижчих.)

Нарешті, якщо використовувати вирази (3.1) і (3.3) для теплового потоку,
то отримуємо бажані граничні умови для задачі, описаної на рис. 3.3, у
чисто математичній формі

\[
(\Gamma Y) \quad \begin{cases} \frac{\partial u\left(0,\,t\right)}{\partial x} = \frac{h}{k}\left[u\left(0,\,t\right) - g_{1}\left(t\right)\right], \\ \frac{\partial u\left(L,\,t\right)}{\partial x} = -\frac{h}{k}\left[u\left(L,\,t\right) - g_{2}\left(t\right)\right]. \end{cases} \quad 0 < t < \infty,
\]

Дуже часто константа h/k позначається як \(\lambda\) , а граничні умови
теплового потоку на межі записуються у вигляді

(3.5)

\[
u_{x}(0, t) = \lambda [u(0, t) - g_{1}(t)], u_{x}(L, t) = -\lambda [u(L, t) - g_{2}(t)].
\]

Для випадку вищих розмірностей отримують подібні граничні умови.
Наприклад, якщо межа кругового диска осяяна рухомою рідиною з
температурою \(g(\theta,t)\) , то граничні умови записуються як

\[
\frac{\partial u}{\partial r}(R, \theta, t) = -\frac{h}{k} [u(R, \theta, t) - g(\theta, t)].
\]

Тут \(\frac{\partial u}{\partial r}(R, \theta, t)\) позначає зовнішню
нормальну похідну (у позитивному напрямку r-осі) значення u, обчислену в
точці \((R, \theta)\) межі. Ми назвемо такі граничні умови лінійними
(оскільки вони лінійні в U і \(u_r\) ), але неоднорідними, оскільки
функція \(g(\theta, t)\) знаходиться праворуч.

\section{Граничні умови третього виду (враховуючи потік; зокрема,
враховується випадок термічно ізольованих
меж)}\label{ux433ux440ux430ux43dux438ux447ux43dux456-ux443ux43cux43eux432ux438-ux442ux440ux435ux442ux44cux43eux433ux43e-ux432ux438ux434ux443-ux432ux440ux430ux445ux43eux432ux443ux44eux447ux438-ux43fux43eux442ux456ux43a-ux437ux43eux43aux440ux435ux43cux430-ux432ux440ux430ux445ux43eux432ux443ux454ux442ux44cux441ux44f-ux432ux438ux43fux430ux434ux43eux43a-ux442ux435ux440ux43cux456ux447ux43dux43e-ux456ux437ux43eux43bux44cux43eux432ux430ux43dux438ux445-ux43cux435ux436}

Жоден потік не проходить через термоізольовані межі, тому нормальна
похідна (зовнішня чи внутрішня) має зрівнятися з нулем на межі, оскільки
потік пропорційний нормальній похідній. У випадку одномірного елемента з
ізольованими кінцями x=0 і x=L, граничні умови мають вигляд

\[
u_{\lambda}(0, t) = 0,
\]

\(u_{\lambda}(L, t) = 0.\) \(0 < t < \infty,\) У двовимірній області
ізоляція межі означає, що нормальна похідна температури на межі
перетворюється на нуль. Припустимо, наприклад, що круглий диск термічно
ізольований уздовж межі. Тоді граничну умову слід записати так:
\(u_r(R, \theta, t) = 0\) для всіх \(0 \le \theta < 2\pi\) і всіх
\(0 < t < \infty\) .

З іншого боку, якщо вказати кількість тепла, що проходить через межу
нашого диска, то граничні умови набувають вигляду

\[
u_r(R, \theta, t) = f(\theta, t),
\]

де \(f(0, t)\) --- це кількість тепла, що переходить у круглий диск із
зовнішнього джерела.

Тепер проілюструємо різні типи граничних умов.

\section{Типові граничні умови для задачі одномірної
теплопровідності}\label{ux442ux438ux43fux43eux432ux456-ux433ux440ux430ux43dux438ux447ux43dux456-ux443ux43cux43eux432ux438-ux434ux43bux44f-ux437ux430ux434ux430ux447ux456-ux43eux434ux43dux43eux43cux456ux440ux43dux43eux457-ux442ux435ux43fux43bux43eux43fux440ux43eux432ux456ux434ux43dux43eux441ux442ux456}

Припустимо, що у нас є мідний стрижень довжиною 200 см, бікова поверхня
якого термічно ізольована, а початкова температура становить 0 °C.
Припустимо також, що верхній кінець стрижня x=0 термічно ізольований, а
нижній кінець x=200 змивається рухомою водою, яка має постійну
температуру \(g_2(t)=20\) °C (див. рис. 3.6).

\pandocbounded{\includegraphics[keepaspectratio]{images/_page_25_Picture_6.jpeg}}

Фіг. 3.6. Змішане завдання.

Математична модель цієї задачі описується наступними чотирма
відношеннями:

(3.6)

\begin{verbatim}
\begin{aligned}
(\text{УЧП})\quad & u_t = \alpha^2 u_{xx}, && 0 < x < 200,\ 0 < t < \infty, \\
(\text{ГУ})\quad & u_x(0,t) = 0, && 0 < t < \infty, \\
& u_x(200,t) = -\frac{h}{k}\left[u(200,t) - 20\right], && 0 < t < \infty, \\
(\text{НУ})\quad & u(x,0) = 0^\circ\mathrm{C}, && 0 \leqslant x \leqslant 200.
\end{aligned}
\end{verbatim}

де \(\alpha^2 = 1,16\) см 2/с --- теплова дифузія міді;

k = 0,93 кал/см·с·°C --- теплопровідність міді;

h --- коефіцієнт теплопередачі. Визначення цього співвідношення ---
велика проблема. Він вимірюється швидкістю виходу тепла від другого
кінця стрижня до води. Це залежить від швидкості потоку води, стану
поверхні, яку слід змивати тощо.

Читач може спробувати провести експеримент, щоб визначити

ділення значення цього коефіцієнта.

\section{ЗАУВАЖЕННЯ}\label{ux437ux430ux443ux432ux430ux436ux435ux43dux43dux44f-2}

\begin{enumerate}
\def\labelenumi{\arabic{enumi}.}
\tightlist
\item
  Типова задача внутрішньої теплопровідності для квадрата показана на
  рис. 3.7.
\end{enumerate}

У цій задачі початкова температура \(u\left(x,\,y,\,t\right)\) при t = 0
задається всередині квадрата, і рівняння з похідними з частковими має
вигляд Р и с

\pandocbounded{\includegraphics[keepaspectratio]{images/_page_26_Figure_9.jpeg}}

. 3.7. Типові граничні умови для задачі дифузії в квадраті: a ---
кількість тепла, \(f_1(t)\) проходить через межу; b --- термоізольована
межа; a --- це навколишня температура Равії \(10^{\circ}\) C; z ---
температура підтримується рівною \(10^{\circ}\) C.

а граничні умови будуть виконані на \(0 < t < \infty\) і визначатимуть
температуру в будь-який момент. Якою б не була ця температура, вона
повинна відповідати граничним умовам, зазначеним на рис. 3.7.

\begin{enumerate}
\def\labelenumi{\arabic{enumi}.}
\setcounter{enumi}{1}
\tightlist
\item
  Зверніть увагу, що гранична умова
\end{enumerate}

\[
u_r(R, \theta, t) = -\frac{h}{k} [u(R, \theta, t) - g(\theta, t)]
\]

На колі не означає, що гранична температура дорівнює \(g(\theta,t)\) ,
але якщо коефіцієнт теплопередачі h великий, то ці умови фактично
еквівалентні вимогі \(g(\theta,t)\) температури на межі.

\section{ЗАВДАННЯ}\label{ux437ux430ux432ux434ux430ux43dux43dux44f-2}

\begin{enumerate}
\def\labelenumi{\arabic{enumi}.}
\item
  Намалюйте контури графіків для розв'язання змішаної задачі (3.6) для
  різних моментів часу. Чи відповідають ваші графи визначеним граничним
  умовам? Яка стаціонарна температура стрижня? Очевидно, ваша відповідь
  базується на інтуїції?
\item
  Як би ви інтерпретували наступне змішане завдання?
\end{enumerate}

\[
\begin{aligned}
(\text{УЧП})\quad & u_t = \alpha^2 u_{xx}, && 0 < x < 1,\ 0 < t < \infty, \\
(\text{ГУ})\quad & u(0,t) = 0,\quad u_x(1,t) = 1, && 0 < t < \infty, \\
(\text{НУ})\quad & u(x,0) = \sin(\pi x), && 0 \leqslant x \leqslant 1.
\end{aligned}
\]

Чи можете ви графічно показати рішення цієї проблеми в різні моменти
часу? Чи буде це рішення, як правило, стаціонарним? Це очевидно?

\begin{enumerate}
\def\labelenumi{\arabic{enumi}.}
\setcounter{enumi}{2}
\tightlist
\item
  Яку фізичну інтерпретацію ви можете дати проблемі
\end{enumerate}

\begin{verbatim}
\begin{aligned}
(\text{УЧП})\quad & u_t = \alpha^2 u_{xx}, && 0 < x < 1,\ 0 < t < \infty, \\
(\text{ГУ})\quad & u_x(0,t) = 0,\quad u_x(1,t) = 0, && 0 < t < \infty, \\
(\text{НУ})\quad & u(x,0) = \sin(\pi x), && 0 \leqslant x \leqslant 1.
\end{aligned}
\end{verbatim}

Чи можете ви графічно показати рішення цієї проблеми в різні моменти
часу?

Що можна сказати про стабільну температуру?

\begin{enumerate}
\def\labelenumi{\arabic{enumi}.}
\setcounter{enumi}{3}
\tightlist
\item
  Припустимо, що бічна поверхня металевого стрижня не ізольована і має
  початкову температуру 20°C, але температура одного кінця стає 50°C.
  Потім стрижень промивають рідиною з температурою 30°C. Як виглядає
  змішана задача, що відповідає цьому випадку?
\end{enumerate}

\section{ВИВЕДЕННЯ РІВНЯННЯ
ТЕПЛОПРОВІДНОСТІ}\label{ux432ux438ux432ux435ux434ux435ux43dux43dux44f-ux440ux456ux432ux43dux44fux43dux43dux44f-ux442ux435ux43fux43bux43eux43fux440ux43eux432ux456ux434ux43dux43eux441ux442ux456}

МЕТА ЛЕКЦІЇ: Показати, як можна отримати одномірне рівняння
теплопровідності

\[
u_t = \alpha^2 u_{xx} + f(x, t),
\]

спираючись лише на фундаментальний закон збереження кількості тепла.
Розглядається залежність швидкості теплопередачі від основних
термофізичних параметрів: теплопровідності, теплоємності та густини
матеріалу. Розглядаються різні модифікації базового рівняння
теплопровідності.

У кожній галузі знань існує набір базових принципів, правильність яких
очевидна. Усі інші твердження мають бути виведені з цих базових
принципів. Звісно, те, що очевидно одній людині, може здатися сумнівним
для іншої. Історія будь-якої науки полягає у постійному пошуку більш
фундаментальних принципів, що лежать в основі цієї науки. Основні
принципи мають бути настільки універсальними, щоб ніхто не сумнівався.

? --- Припущення A --- Припущення B --- Припущення C --- Припущення D
--- ? Рис. 4.1. Аксіоматичний метод.

Наприклад, одна людина може вважати, що всі факти певної науки можна
вивести з одного базового припущення, яке ми позначимо як B. На основі
припущення B вона або вона може довести теорему C, що призведе до
теореми D, а яка, у свою чергу, призведе до доведення багатьох інших
теорем (див. рисунок 4.1). Ось як виглядає прогрес у науці --- дослідник
може отримувати все більше нових результатів. Найвидатніші фізики,
хіміки та біологи або саме на цьому шляху.

З іншого боку, замість доведення нових теорем можна спробувати, якщо
можливо, знайти новий принцип (назвемо його A), який є більш
фундаментальним за B. Припущення A природно вважається більш
фундаментальним, ніж B.

Якщо B можна вивести з A. Отже, якщо можна знайти новий принцип A, тоді
межа наших знань зміщується до сфери більш фундаментальних ідей про
певну науку. У теорії теплопровідності основним принципом є закон
збереження енергії (теплова енергія). Усі інші твердження випливають із
цього базового принципу (рис. 4.2).

Закон збереження енергії

\[
u_t = \alpha^2 u_{xx} + f(x,t)
\]

Інші властивості теплопровідності Рис. 4.2. Енергозбереження: наріжний
камінь теплопровідності.

Звісно, ми могли б забути цю лекцію і використати рівняння
теплопровідності як фундаментальне поняття (дехто може так подумати,
якщо це очевидно), але тоді ми обдуримо більш серйозних студентів, адже
закон збереження вважається найфундаментальнішим принципом у науці.
Науковці, моделюючи свої задачі, часто записують відношення, що
виражають закони збереження, а потім переписують їх у вигляді рівнянь у
похідних похідних.

Повернемося до головної мети нашої лекції --- вивести рівняння
теплопровідності з рівняння збереження енергії.

\section{Виведення рівняння
теплопровідності}\label{ux432ux438ux432ux435ux434ux435ux43dux43dux44f-ux440ux456ux432ux43dux44fux43dux43dux44f-ux442ux435ux43fux43bux43eux43fux440ux43eux432ux456ux434ux43dux43eux441ux442ux456-1}

Розглянемо однорідний стрижень довжини L, щодо якого зробимо такі
припущення:

\begin{enumerate}
\def\labelenumi{\arabic{enumi}.}
\tightlist
\item
  Стрижень виготовлений з одного однорідного провідного матеріалу.
\end{enumerate}

\pandocbounded{\includegraphics[keepaspectratio]{images/_page_29_Figure_10.jpeg}}

Фіг. 4.3. Тонкий тепловий стержень.

\begin{enumerate}
\def\labelenumi{\arabic{enumi}.}
\setcounter{enumi}{1}
\item
  Бікова поверхня стрижня термічно ізольована (тепло може поширюватися
  лише вздовж осі x).
\item
  Стрижень тонка, тобто температура всіх точок у кожному поперечному
  перерізі стрижня є сталою. Якщо розглядати частину вудки на сегменті
  \([x, x+\Delta x]\) і використовувати
\end{enumerate}

закон збереження кількості тепла, тоді можна написати:

(4.1) Загальна зміна кількості тепла на сегменті \([x, x + \Delta x] =\)
дорівнює загальній кількості тепла, що пройшло через межі + Нульова
кількість тепла, утвореного всередині сегмента \([x, x + \Delta x]\) .

Оскільки загальна кількість тепла (у калоріях) всередині сегмента
\(\{x, x + \Delta x\}\) в будь-який момент часу t обчислюється формулою:

Загальна кількість тепла всередині сегмента
\([x, x+\Delta x] = \int\limits_{-\infty}^{x+\Delta x} coAu\) (s, t) ds,

де c --- питома теплоємність матеріалу (показує здатність матеріалу
накопичувати тепло),

Щільність матеріалу,

A --- це площа поперечного перерізу стрижня.

\textbf{va}con збереження енергії (4.1) можна надати у такій
математичній формі:

(4.2)

\[
\frac{d}{dt} \int_{x}^{x+\Delta x} c\rho Au(s, t) ds = c\rho A \int_{x}^{x+\Delta x} u_{t}(s, t) ds = eA \left[ u_{x}(x+\Delta x, t) - u_{x}(x, t) \right] + A \int_{x+\Delta x}^{x+\Delta x} f(s, t) ds.
\]

\textbf{де} k --- теплопровідність матеріалу (вказує на здатність
матеріалу проводити тепло), \(f(x, t)\) --- об'ємна потужність
зовнішнього джерела тепла (калорії/см·с).

Задача полягає в тому, щоб записати рівняння (4.2) у формі, яка не
містить інтегралів. Щоб розв'язати цю задачу, нагадаємо читачеві теорему
про середнє значення з ходу інтегральної нечисельності.

\section{Теорема про
значення}\label{ux442ux435ux43eux440ux435ux43cux430-ux43fux440ux43e-ux437ux43dux430ux447ux435ux43dux43dux44f}

Якщо функція \(f(x)\) є неперервною на сегменті {[}a, b{]}, то існує
принаймні одна точка \(\xi \in [a, b]\) така, що

\[
\int_{a}^{b} f(x) dx = f(\xi) (b-a).
\]

Застосовуючи цей результат до рівняння (4.2), ми отримуємо наступний
результат

\[
coAu_{t}(\xi, t) \Delta x = kA\{u_{x}(x + \Delta x, t) - u_{x}(x, t)\} + Af(\xi, t) \Delta x,
\]

x \textless{} \xi \textless{} x + \Delta x, або \(u_t\)(\xi,
t)\texttt{=\ \textbackslash{}frac\{k\}\{c\textbackslash{}rho\}\ \textbackslash{}left\textbackslash{}\{\ \textbackslash{}frac\{}u\_x\((x + \Delta x, t) -\)u\_x\texttt{(x,\ t)\}\{\textbackslash{}Delta\ x\}\ \textbackslash{}right\textbackslash{}\}\ +\ \textbackslash{}frac\{1\}\{c}\rho\}
f(\xi, t).

\begin{verbatim}
Спрямовуючи \Delta x до 0, отримуємо потрібне рівняння
\end{verbatim}

(4.3)

\[
u_t(x, t) = \alpha^2 u_{xx}(x, t) + F(x, t),
\]

де

\[
\alpha^{2}=\frac{k}{c\rho}
\]

\begin{verbatim}
— коефіцієнт теплової дифузії, F\left(x,\ t\right)=\frac{1}{c\rho}f\left(x,\ t\right) — густина джерел тепла.
\end{verbatim}

Ми могли б зупинитися на цьому. Однак ми також розглянемо випадок, коли
бічна поверхня стрижня не термічно ізольована. Припустимо, що кількість
теплового потоку через бічну поверхню стрижня в цьому випадку
пропорційна різниці між температурою стрижня \(u(x,t)\) і температурою
навколишнього середовища, яка залишається сталою і дорівнює нулю. У
цьому випадку закон збереження кількості тепла веде до рівняння

\[
(4.4) u_t = \alpha^2 u_{xx} - \beta u + F(x, t),
\]

\begin{verbatim}
де \beta — коефіцієнт пропорційності потоку через бічну поверхню.
\end{verbatim}

\section{ЗАУВАЖЕННЯ}\label{ux437ux430ux443ux432ux430ux436ux435ux43dux43dux44f-3}

\begin{enumerate}
\def\labelenumi{\arabic{enumi}.}
\tightlist
\item
  Як зазначалося вище, константа k називається теплопровідністю речовини
  і чисельно дорівнює кількості тепла (у калоріях), що проходить за
  секунду через пластину товщиною 1 см і площею поперечного перерізу 1
  см³ при різниці температур на протилежних гранях пластини рівною 1°C.
  Числові значення теплопровідності для різних матеріалів можна знайти в
  довідниках з фізики або хімії. Типові значення теплопровідності
  варіюються від 1 (для міді) до значень, близьких до нуля для хороших
  теплоізоляторів. Якщо стрижень зроблений із однорідного матеріалу, то
  значення k не залежить від x. Для деяких матеріалів значення k
  залежить від температури, і тому рівняння теплопровідності в цьому
  випадку стає нелінійним:
\end{enumerate}

\[
u_{t} = \frac{1}{c\rho} \frac{\partial}{\partial x} \{ k (u) | u_{x} \}.
\]

Однак найчастіше теплопровідність k дуже слабо впливає на температуру u,
і цю нелінійність можна ігнорувати.

\begin{enumerate}
\def\labelenumi{\arabic{enumi}.}
\setcounter{enumi}{1}
\item
  Значення c зазвичай називають питомою теплоємністю речовини. Він
  характеризує здатність речовини зберігати теплову енергію. Наприклад,
  запечена картопля має високу теплоємність, оскільки може зберігати
  велику кількість тепла на одиницю маси (саме тому вона зберігає високу
  температуру так довго). Чисельна теплоємність дорівнює кількості
  калорій тепла, які потрібно передати 1 граму речовини, щоб підвищити
  її температуру на 1°C. Для більшості розглянутих задач значення c
  можна вважати сталим, не записуючи ні x, ні i. Значення питомої
  теплоємності для конкретних матеріалів можна знайти у фізичних
  довідниках.
\item
  Одиниці вимірювання основних величин, включених у рівняння
\end{enumerate}

теплопровідність (у системі GHS) наступна:

i-температура (градуси Цельсія), \(u_t\) --- швидкість зміни температури
(°C/s), \(u_x\) --- нахил температурної кривої (°C/см), \(u_{xx}\) ---
опуклість температурної кривої (°C/см²),

\begin{verbatim}
c = щільність тепла (калор/г \cdot °C)
\end{verbatim}

\begin{verbatim}
k — теплопровідність (кал/см \cdot при \cdot °C), \rho — щільність (г/<sup>см3</sup>), \alpha^2 — це теплова дифузія (см<sup>²</sup>/с).
\end{verbatim}

\begin{verbatim}
4. Слід зазначити, що теплова дифузія матеріалу \alpha^2 = \frac{k}{\epsilon \rho} прямо пропорційна теплопровідності матеріалу і обернено пропорційна густині \rho та питому теплоємності \epsilon ; Ми сподіваємося, що цей зв'язок відповідає інтуїтивним ідеям читача.
\end{verbatim}

\section{ЗАВДАННЯ}\label{ux437ux430ux432ux434ux430ux43dux43dux44f-3}

\begin{verbatim}
1. Підставити розмірність кожної з u,\ u_t,\ \dots величин у рівняння
\end{verbatim}

\[
u_t = \alpha^2 u_{xx} - \beta u
\]

\begin{verbatim}
і переконайтеся, що всі терміни мають однаковий розмір {}^{\circ}\mathrm{C}/c .
\end{verbatim}

\begin{enumerate}
\def\labelenumi{\arabic{enumi}.}
\setcounter{enumi}{1}
\tightlist
\item
  Підставити розмірність кожної з величин у рівняння
\end{enumerate}

\[
u_t = \alpha^2 u_{xx} - v u_x,
\]

\begin{verbatim}
де \upsilon вимірюється в одиницях швидкості, і переконайтеся, що всі терми мають однакову розмірність.
\end{verbatim}

3 Отримаємо рівняння теплопровідності

\[
u_{t} = \frac{1}{c\rho} \frac{\partial}{\partial x} [h(x) u_{x}] + f(x, t)
\]

\begin{verbatim}
Для випадку, коли теплопровідність k\left( x\right) залежить від координати X.
\end{verbatim}

\begin{enumerate}
\def\labelenumi{\arabic{enumi}.}
\setcounter{enumi}{3}
\tightlist
\item
  Припустимо, що \(u(x, t)\) --- це значення концентрації речовини в
  потоку, що рухається зі швидкістю v. Припустимо, що концентрація
  змінюється внаслідок процесів дифузії та конвекції. Виведемо рівняння
\end{enumerate}

\[
u_1 = \alpha^2 u_{xx} - v u_x,
\]

\begin{verbatim}
Виходячи з того, що в будь-який момент часу в певному [x, x + \Delta x] сегменті матерії не з'являється і не зникає.
\end{verbatim}

ПРИМІТКА. Закон збереження матерії виглядає так:

\begin{verbatim}
Зміни маси речовини на сегменті [x, x + \Delta x] = = Зміна внаслідок дифузії через межі + Зміна внаслідок переносу через межі.
\end{verbatim}

Лекція 5

\section{ЗМІННЕ
РОЗДІЛЕННЯ}\label{ux437ux43cux456ux43dux43dux435-ux440ux43eux437ux434ux456ux43bux435ux43dux43dux44f}

МЕТА ЛЕКЦІЇ: Познайомити читача з потужним методом розділення змінних і
показати, як цей метод можна використати для розв'язання добре відомої
проблеми дифузії. Оскільки цей метод створює труднощі для деяких учнів
(головним чином через громіздкі алгебранічні перетворення), ми
супроводжуємо презентацію інтуїтивними міркуваннями.

Основна ідея методу полягає в тому, щоб розкласти початкову умову на
найпростіші компоненти, знайти відгук системи на кожен найпростіший
компонент і підсумувати всі відповіді. Таким чином можна знайти
відповідь на довільну початкову умову.

Існуючий метод поділу змінних певною мірою зменшує можливість такої
інтерпретації методу, але все ж не виключає його.

Метод розділення змінних є одним із найвідоміших методів розв'язання
зсувних задач і використовується, якщо:

\begin{itemize}
\tightlist
\item
  \begin{enumerate}
  \def\labelenumi{\arabic{enumi}.}
  \tightlist
  \item
    Рівняння є лінійним і однорідним (не обов'язково з постійними
    коефіцієнтами).
  \end{enumerate}
\item
  \begin{enumerate}
  \def\labelenumi{\arabic{enumi}.}
  \setcounter{enumi}{1}
  \tightlist
  \item
    Граничні умови визначаються як
  \end{enumerate}
\end{itemize}

\[
\alpha u_x(0, t) + \beta u(0, t) = 0, \qquad \gamma u_x(1, t) + \delta u(1, t) = 0.
\]

\begin{verbatim}
де \alpha, \beta, \gamma і \delta — константи (граничні умови, задані в такій формі, називаються лінійними однорідними граничними умовами).
\end{verbatim}

Метод був створений за часів Фур'є (зазвичай відомий як метод Фур'є) і,
здається, є найпопулярнішим (де це застосовно).

\pandocbounded{\includegraphics[keepaspectratio]{images/_page_34_Figure_8.jpeg}}

Рис. 5.1. Схематичне представлення задачі дифузії.

Замість того, щоб вивчати метод у загальному випадку, давайте спочатку
розглянемо конкретну задачу (загальний випадок обговоримо пізніше).
Розглянемо змішану задачу типу дифузії: щоб знайти розв'язок

(PPP)

\[
u_t = \alpha^2 u_{xx}
, 0 < x < 1 , 0 < t < \infty ,
\]

що задовольняє граничні умови

(GI)

\[
\begin{cases} u(0, t) = 0, & 0 < t < \infty, \\ u(1, t) = 0, & 0 < t < \infty, \end{cases}
\]

та початкова умова

(ГУ)

\[
u(x, 0) = \varphi(x), 0 \le x \le 1.
\]

\begin{verbatim}
Перш ніж почати розділяти змінні, давайте дамо фізичну інтерпретацію нашої проблеми. Отже, існує стрижень скінченної довжини, кінці якого підтримуються при сталій температурі, рівній нулю (насправді кінці можна підтримувати при значно вищій температурі, значення якої вважається початковою точкою). Додаткові дані про завдання подаються як початкова умова. Наша мета — знайти розподіл температури u\left(x,t\right) в наступний час.
\end{verbatim}

\section{Загальні принципи методу розділення
змінних}\label{ux437ux430ux433ux430ux43bux44cux43dux456-ux43fux440ux438ux43dux446ux438ux43fux438-ux43cux435ux442ux43eux434ux443-ux440ux43eux437ux434ux456ux43bux435ux43dux43dux44f-ux437ux43cux456ux43dux43dux438ux445}

Для найпростішого диференціального рівняння в похідних похідних
розділення змінних є пошуком розв'язків у вигляді

\[
u(x, t) = X(x)T(t),
\]

де \(X(x)\) --- це функція, що залежить лише від змінної x, а \(T(t)\)
--- функція, що залежить лише від t. Це розв'язок у певному сенсі є
найпростішим, оскільки температура \(u(x, t)\), представлена у такій
формі, зберігатиме «форму» профілю у різний час часу t (див. рисунок
5.2).

\pandocbounded{\includegraphics[keepaspectratio]{images/_page_35_Figure_6.jpeg}}

РИСУНОК 5.2 Графік функції X(x)\$\$T(t) у різних точках часу t.

Загальна ідея полягає в тому, щоб знайти нескінченну кількість таких
часткових диференціальних розв'язків (які задовольняють граничні умови

\[
u(0, t) = 0, \qquad u(1, t) = 0.
\]

Для цього підставляємо розв'язки (5.1) у ці граничні умови. У результаті
отримуємо

\[
u(0, t) = B e^{-\lambda^2 \alpha^2 t} = 0 \Rightarrow B = 0,
u(1, t) = A e^{-\lambda^2 \alpha^2 t}\sin \lambda = 0 \Rightarrow \sin \lambda = 0.
\]

Отже,

\[
\lambda_n = n\pi, \qquad n = 1, 2, \ldots
\]

Зверніть увагу, що друга гранична умова могла б виконуватися і при
\(A=0\) , але тоді розв'язок (5.1) був би тривіально нульовим.

\pandocbounded{\includegraphics[keepaspectratio]{images/_page_38_Figure_3.jpeg}}

Фіг. 5.3. Фундаментальне рішення
\(n_n(x, t) = A_n e^{-(\pi n \alpha)^{n/t}} \sin{(\pi n x)}\) .

Тепер, коли ми зробили другий крок, маємо нескінченний набір ознак

(5.2)

\[
u_n(x, t) = A_n e^{-(n\pi\alpha)^2 t} \sin(n\pi x), \quad n = 1, 2 \dots
\]

кожна з яких задовольняє рівняння з частковими похідними та граничними
умовами). Отримані функції --- це будівельні блоки, з яких ми будемо
створювати рішення нашої проблеми. Це рішення буде певною сумою

2) Зверніть увагу, що функції \(u_n\) і \(u_{-n}\) відрізняються лише
знаком.

цих найпростіших функцій. Конкретний тип суми залежить від початкового
стану. Фіг. Рисунок 5.3 показує графіки кількох таких фундаментальних
функцій

КРОК 3. (Пошук розв'язку, який задовольняє рівняння, межу та початкові
умови.)

Останній (і, ймовірно, найцікавіший з математичної точки зору) крок ---
знайти таку суму фундаментальних розв'язків

\[
u(x, t) = \sum_{n=1}^{\infty} A_n e^{-(n\pi\alpha)^n t} \sin(n\pi x),
\]

тобто при виборі таких коефіцієнтів \(A_n\) , що функція задовольнятиме
початкову умову

\[
u\left( x,\ 0\right) =\varphi\left( x\right) .
\]

Підставляючи суму в початковому стані, отримуємо

(5.3)

\[
q(x) = \sum_{n=1}^{\infty} A_n \sin(n\pi x).
\]

Це рівняння підводить нас до цікавого питання: чи можна розкласти
початкову температуру \(\phi(x)\) на ряд елементарних функцій вигляду

\[
A_1 \sin(\pi x) + A_2 \sin(2\pi x) + A_3 \sin(3\pi x) + \dots
\]

?

Позитивну відповідь на це питання дав французький математик Жозеф Фур'є.
Виявилося, що для достатньо хороших функцій така декомпозиція можлива
\(^1\) ). Тоді виникає нове питання: знайти коефіцієнти розкладу \(A_n\)
?

Насправді це дуже просто. Для цього використаємо властивість системи
функцій

\[
\{\sin(n\pi x)\}_{n=1}^{\infty}.
\]

відомий як ортогональність.

У нашому випадку (див. завдання 2) ці функції задовольняють умови

\[
\int_{0}^{1} \sin(m\pi x) \sin(n\pi x) dx = \begin{cases} 0, & m \neq n, \\ 1/2, & m = n. \end{cases}
\]

Умови розкладу функцій у рядах Фур'є наведені в лекції 11. Для глибшого
ознайомлення з цим питанням можна звернутися до численних освітніх
літератур (див., наприклад: Г. П. Толстов, Серія Фур'є. --- Москва:
Наука, 1980). --- Ед. Ед.

Ця властивість чітко показана на рис. 5.4, яка показує графіки деяких
функцій із цієї системи.

Отже, ми хочемо знайти коефіцієнти у розкладі

\[
\varphi(x) = A_1 \sin(\pi x) + A_2 \sin(2\pi x) + A_3 \sin(3\pi x) + \dots
\]

\pandocbounded{\includegraphics[keepaspectratio]{images/_page_40_Figure_5.jpeg}}

РИСУНОК 5.4. Послідовність ортогональних функцій

Помножте обидві сторони цього співвідношення на \(\sin{(m\pi x)}\) (m
--- довільне ціле число) і інтегруйте від нуля до одиниці. У результаті
ми отримуємо

\[
\int_{0}^{1} \varphi(x) \sin(m\pi x) \, dx = A_{m} \int_{0}^{1} \sin^{2}(m\pi x) \, dx = \frac{A_{m}}{2}
\]

(усі інші члени зрівнялися з нулем через ортогональність). Розв'язуючи
рівняння відносно \(A_m\) , отримуємо

\[
A_m = 2 \int_0^1 \varphi(x) \sin(m\pi x) dx.
\]

Отже, отримано, що розв'язок записується у вигляді

(5.4)

\[
u(x, t) = \sum_{n=1}^{\infty} A_n e^{-(n\pi\alpha)^2 t} \sin(\pi n x),
\]

де коефіцієнти \(A_n\) визначаються формулами

\[
A_n = 2\int_0^1 \varphi(x) \sin(\pi nx) dx
\]

Ви можете переконатися, що отримане рішення задовольняє всі умови
початкової задачі. На цьому завершується третій крок.

Багато студентів розчаровані, коли бачать, наскільки складним виглядає
рішення, і майже не намагаються дивитися на нього знову (виглядає дуже
погано). Однак це рішення не здаватиметься таким складним, якщо ви
витратите час на його аналіз. Насправді, складніша форма розв'язку
робить його більш інформативним. Нижче наведено кілька нотаток, які
допоможуть вам інтерпретувати це рішення.

\section{ЗАУВАЖЕННЯ}\label{ux437ux430ux443ux432ux430ux436ux435ux43dux43dux44f-4}

Зверніть увагу, що єдина різниця між синусоїдальним розкладом функції
φ(x) і розв'язком (5.4) --- це наявність часового фактора e-(nπα)31 у
кожному члені ряду. Отже, якщо початкова умова має просту декомпозицію,
наприклад, вигляду

\[
\varphi(x) = \sin(\pi x) + \frac{1}{2}\sin(3\pi x),
\]

тоді рішення можна негайно записати

\[
u(x, t) = e^{-(\pi\alpha)^2 t} \sin(\pi x) + \frac{1}{2} e^{-(3\pi\alpha)^2 t} \sin(3\pi x).
\]

У цьому випадку очевидно, що якщо ми розкладемо \(\phi(x)\) на ряд Фур'є
вздовж синусів, то отримуємо

\[
A_1 = 1,
\]

\(A_2 = 0,\) \(A_3 = 1/2,\) \(A_4 = A_5 = \dots = 0.\) - 2. Розв'язок
(5.4) можна інтерпретувати так: ми представляємо початкову температуру
\(\varphi(x)\) як суму найпростіших функцій виду \(A_n \sin(\pi nx)\) ;
кожна така функція породжує відповідь
\(A_n e^{-(\pi n\alpha)^2 t} \sin(\pi nx)\) . Сумуючи всі такі
відповіді, отримуємо розв'язок, який відповідає початковій умові
\(u(x, 0) = \varphi(x)\) . - 3. Кожен член у декомпозиції

\[
u(x, t) = A_1 e^{-(\pi\alpha)^2 t} \sin(\pi x) + A_2 e^{-(2\pi\alpha)^2 t} \sin(2\pi x) + \cdots
\]

є функцією від x і t. Зверніть увагу, що внесок членів з великими
числами в t\textgreater0 дуже малий через множник
\(e^{-(n\pi\alpha)^2 t}\) . Відповідно, після закінчення досить великого
Р и с

\pandocbounded{\includegraphics[keepaspectratio]{images/_page_42_Figure_2.jpeg}}

. 5.5. У задачах дифузії вищі порядки терміни швидко затухають;
\(\phi(x)\) --- початкова температура. Слід зазначити, що по-перше, у
розчині зникають високі частоти --- по-перше, зникли найбільші коливання
міді.

Повний розв'язок приблизно збігається з першим членом, який є півхвилі
синусоїди, що з часом згасає (рис. 1). 5.5.)

\section{ЗАВДАННЯ}\label{ux437ux430ux432ux434ux430ux43dux43dux44f-4}

\begin{enumerate}
\def\labelenumi{\arabic{enumi}.}
\tightlist
\item
  Показати, що функції
\end{enumerate}

\[
u(x, t) = e^{-\lambda^2 \alpha' t} \left[ A \sin(\lambda x) + B \cos(\lambda x) \right]
\]

задовольняють рівняння \(u_l = \alpha^2 u_{xx}\) для довільних значень
A, B і \(\lambda\) .

\begin{enumerate}
\def\labelenumi{\arabic{enumi}.}
\setcounter{enumi}{1}
\tightlist
\item
  Покажи це
\end{enumerate}

\[
\int_0^1 \sin(m\pi x) \sin(n\pi x) dx = \begin{cases} 0, & m \neq n, \\ 1/2, & m = n. \end{cases}
\]

ПРИМІТКА. Використання ідентичності

\[
\sin(mx)\sin(nx) = \frac{1}{2}[\cos(m-n)x - \cos(m+n)x].
\]

\begin{enumerate}
\def\labelenumi{\arabic{enumi}.}
\setcounter{enumi}{2}
\tightlist
\item
  Знайти розклад у рядах Фур'є синусами функції \(\varphi(x) = 1\) на
  сегменті {[}0, 1{]}. Побудуйте перші три терміни розкладу.
\end{enumerate}

Використовуючи результати завдання 3, знайдіть розв'язок наступної
змішаної задачі:

(PPP)

\[
u_t = u_{xx}
\]

\(0 < x < 1\) , (GI)

\[
\begin{cases} u(0, t) = 0, \\ u(1, t) = 0, \end{cases}
\]

\(0 < t < \infty\) , (НУ) \(u(x, 0) = 1, \quad 0 \le x \le 1\) .

(Зверніть увагу, що ця проблема фізично не має значення, оскільки
припускається, що температура на кінцях стрижня миттєво знижується з
одного до нуля \(\phi(x)\) .

\begin{enumerate}
\def\labelenumi{\arabic{enumi}.}
\setcounter{enumi}{4}
\tightlist
\item
  Знайти розв'язок проблеми 4, якщо початкова умова встановлена у
  вигляді
\end{enumerate}

\[
u(x, 0) = \sin(2\pi x) + \frac{1}{3}\sin(4\pi x) + \frac{1}{5}\sin(6\pi x).
\]

\begin{enumerate}
\def\labelenumi{\arabic{enumi}.}
\setcounter{enumi}{5}
\tightlist
\item
  Розв'яжіть задачу 4, якщо початкова умова
  \(u\left( x,\ 0\right) =x-x^{2},\quad 0< x<1.\) Лекція 6
\end{enumerate}

\section{ПЕРЕТВОРЕННЯ НЕОДНОРІДНИХ ГРАНИЧНИХ УМОВ У ОДНОРІДНІ ГРАНИЧНІ
УМОВИ}\label{ux43fux435ux440ux435ux442ux432ux43eux440ux435ux43dux43dux44f-ux43dux435ux43eux434ux43dux43eux440ux456ux434ux43dux438ux445-ux433ux440ux430ux43dux438ux447ux43dux438ux445-ux443ux43cux43eux432-ux443-ux43eux434ux43dux43eux440ux456ux434ux43dux456-ux433ux440ux430ux43dux438ux447ux43dux456-ux443ux43cux43eux432ux438}

МЕТА ЛЕКЦІЇ: Показати, як виникає змішана задача

\begin{verbatim}
\begin{aligned}
(\text{УЧП})\quad & u_t - a^2 u_{xx} = f(x,t), \\
(\text{ГУ})\quad & \alpha_1 u_x(0,t) + \beta_1 u(0,t) = g_1(t), \\
& \alpha_2 u_x(L,t) + \beta_2 u(L,t) = g_2(t), \\
(\text{НУ})\quad & u(x,0) = \varphi(x).
\end{aligned}
\end{verbatim}

можна звести до задачі вигляду

\[
U_{t} - \alpha^{2}U_{xx} = F(x, t),
\alpha_{1}U_{x}(0, t) + \beta_{1}U(0, t) = 0,
\alpha_{2}U_{x}(L, t) + \beta_{2}U(L, t) = 0,
U(x, 0) = \varphi(x)
\]

з нульовими граничними умовами. Цей новий виклик можна подолати

\begin{itemize}
\tightlist
\item
  \begin{enumerate}
  \def\labelenumi{\arabic{enumi}.}
  \tightlist
  \item
    Шляхом ділення змінних (якщо нова задача виявиться однорідною, тобто
    \(F(x, t)\) = 0).
  \end{enumerate}
\end{itemize}

\begin{verbatim}
- 2. Методом інтегральних перетворень і розкладу власними функціями, якщо F(x, t) \neq 0 .
\end{verbatim}

Метод розділення змінних, про який йшлося в попередній лекції, є досить
потужним, а розв'язки, отримані за його допомогою, подаються у зручній
формі. Однак читач повинен розуміти, що цей метод не застосовний до всіх
задач. Щоб метод розділення змінних був застосовним, граничні умови
мають бути лінійними, однорідними, тобто представленими як

(6.1)

\[
\begin{aligned} \alpha_1 u_x(0, t) + \beta_1 u(0, t) &= 0, \\ \alpha_2 u_x(L, t) + \beta_2 u(L, t) &= 0. \end{aligned}
\]

Мета цієї лекції --- показати, як виникає задача з неоднорідними
граничними умовами форми

(6.2)

\[
\begin{aligned} & (\mathsf{Y}\mathsf{Ch}\Pi) & u_t = \alpha^2 u_{xx}, \\ & (\mathsf{\Gamma}\mathsf{Y}) і \begin{case} \alpha_1 u_x \, (0, \, t) + \beta_1 u \, (0, \, t) = g_1 \, (t) \\ \alpha_2 u_x \, (L, \, t) + \beta_2 u \, (L, \, t) = g_2 \, (t) \end{case} \text{ (гетерогенний } \mathsf{\Gamma}\mathsf{Y}), \\ & (\mathsf{H}\mathsf{Y}) і u \,  (x, \, 0) = \varphi \, (x) \end{aligned}
\]

можна розв'язати, зведивши її до задачі з однорідними граничними
умовами.

Нове завдання можна розв'язати іншими методами

(наприклад, шляхом розкладу за власними функціями).

Почнемо з надзвичайно простого прикладу перетворення гетерогенних
граничних умов на однорідні.

\section{Перетворення неоднорідних граничних умов у однорідні граничні
умови}\label{ux43fux435ux440ux435ux442ux432ux43eux440ux435ux43dux43dux44f-ux43dux435ux43eux434ux43dux43eux440ux456ux434ux43dux438ux445-ux433ux440ux430ux43dux438ux447ux43dux438ux445-ux443ux43cux43eux432-ux443-ux43eux434ux43dux43eux440ux456ux434ux43dux456-ux433ux440ux430ux43dux438ux447ux43dux456-ux443ux43cux43eux432ux438-1}

Розглянемо проблему поширення тепла в термоізольованому стрижні, кінці
якого підтримуються при сталих температурах \(k_1\) і \(k_2\), тобто

\[
\begin{aligned}
(\text{УЧП})\quad & u_t = \alpha^2 u_{xx}, && 0 < x < L,\ 0 < t < \infty, \\
(6.3)\ (\text{ГУ})\quad & u(0,t) = k_1,\quad u(L,t) = k_2, && 0 < t < \infty, \\
(\text{НУ})\quad & u(x,0) = \varphi(x), && 0 \leqslant x \leqslant L.
\end{aligned}
\]

Складність цієї задачі полягає в тому, що, оскільки граничні умови в ній
є гетерогенними, ми не можемо розв'язати її шляхом розділення змінних.
Однак очевидно, що в t \to \infty розв'язок нашої задачі схиляється до
стаціонарного розв'язку, який лінійно (вздовж x) змінюється від
температури k\_1 до температури k\_2 (рис. 6.1). Іншими словами, розумно
припустити, що температуру в нашій задачі можна представити як суму
\textsubscript{}\textasciitilde{}

\section{з двох
термів}\label{ux437-ux434ux432ux43eux445-ux442ux435ux440ux43cux456ux432}

\[
u\left(x,\,t\right)=
\]

\begin{verbatim}
стаціонарна + тимчасова = \uparrow \uparrow Частина розв'язку, що залежить від HY і має тенденцію до нуля з помноженням на =\left[k_1+\frac{x}{L}\left(k_2-k_1\right)\right]+U\left(x,\,t\right). У цьому випадку наше завдання — знайти температуру переходу U\left( x,\,t\right) . Заміна
\end{verbatim}

\[
u(x, t) = \left[k_1 + \frac{x}{L}(k_2 - k_1)\right] + U(x, t)
\]

\begin{verbatim}
У початковій задачі (6.3) ми отримуємо нову задачу відносно невідомої функції U\left(x,\ t\right) . Ви можете вирішити проблему Р и с
\end{verbatim}

\pandocbounded{\includegraphics[keepaspectratio]{images/_page_45_Figure_7.jpeg}}

\begin{verbatim}
. 6.1. Розв'язання задачі (6.3) у різний час часу a- початкова температура u\left(x,\ 0\right)=\phi\left(x\right); b - стаціонарна температура u\left(x,\ \infty\right)=k_{1}+\frac{x}{L}\left(k_{2}-k_{1}\right); b - перехідна температура.
\end{verbatim}

\begin{verbatim}
Відносно нова невідома функція U\left(x,\,t\right) і додає її до стаціонарного розв'язку, що дає бажану функцію u\left(x,\,t\right) . Виконавши ці прості перетворення з (6.3),
\end{verbatim}

Ми отримаємо

(6.4)

\begin{verbatim}
\begin{aligned}
U_t &= \alpha^2 U_{xx}, && 0 < x < L, \\
U(0,t) &= 0,\quad U(L,t) = 0, && 0 < t < \infty, \\
U(x,0) &= \varphi(x) - \left[k_1 + \frac{x}{L}(k_2 - k_1)\right] = \overline{\varphi}(x).
\end{aligned}
\end{verbatim}

\begin{verbatim}
де \overline{\phi}(x) — нова, але відома початкова умова.
\end{verbatim}

Ця задача стосується не лише однорідного рівняння, а, на щастя,
однорідних граничних умов, що дозволяє розв'язувати її методом
розділення змінних. Насправді, читач, ймовірно, пам'ятає рішення:

(6.5)

\[
U(x, t) = \sum_{n=1}^{\infty} a_n e^{-(n\pi\alpha)^4 t} \sin(n\pi x/L),
\]

рей

\[
a_n = \frac{2}{L} \int_{-\infty}^{L} \overline{\phi}(\xi) \sin(\pi n\xi/L) d\xi.
\]

Це все для сердечників із фіксованою температурою на кінцях. А як щодо
більш реалістичних GU з правою стороною, залежною від часу? Основні ідеї
тут ті ж, що й у попередній задачі, але вони набагато складніші.

\section{Перетворення граничних умов, залежних від часу, у
нульові}\label{ux43fux435ux440ux435ux442ux432ux43eux440ux435ux43dux43dux44f-ux433ux440ux430ux43dux438ux447ux43dux438ux445-ux443ux43cux43eux432-ux437ux430ux43bux435ux436ux43dux438ux445-ux432ux456ux434-ux447ux430ux441ux443-ux443-ux43dux443ux43bux44cux43eux432ux456}

Розглянемо типову задачу

(PPP)

\[
u_t = \alpha^2 u_{xx}
\]

\begin{verbatim}
, 0 < x < L , 0 < t < \infty , 0 < t < \infty , 0 < t < \infty , 0 < t < \infty , 0 < t < \infty , 0 < t < \infty , 0 < t < \infty , 0 < t < \infty , 0 < t < \infty , 0 < t < \infty , 0 < t < \infty , 0 < t < \infty , 0 < t < \infty , 0 < t < \infty , 0 < t < \infty , 0 < t < \infty , 0 < t < \infty , 0 < t < \infty , 0 < t < \infty , 0 < t < \infty , 0 < t < \infty , 0 < t < \infty , 0 < t < \infty , 0 < t < \infty , 0 < t < \infty , 0 < t < \infty , , 0 < t < \infty , 0 < t < \infty , 0 < t < \infty , 0 < t < \infty , 0 < t < \infty , 0 < t < \infty , 0 < t < \infty , 0 < t < \infty , 0 < t < \infty , 0 < t < \infty , 0 < t < \infty , 0 < t < \infty , 0 < t < \infty , 0 < t < \infty , 0 < t < \infty , 0 < t < \infty , 0 < t < \infty , 0 < t < \infty , 0 < t < \infty , 0 < t < \infty , 0 < t < \infty , 0 < t < \infty , 0 < t < \infty , 0 < t < \infty , 0 < t < \infty , 0 < t < \infty , 0 < t < \infty , 0 < t < \infty , 0 < t < \infty , 0 < t < \infty , 0 < t < \infty , 0 < t < \infty , 0 < t < \infty , 0 < t < \infty , 0 < t < \infty , 0 < t < \infty , 0 < t < \infty , 0 < t < \infty , 0 < t < \infty , 0 < t < \infty , 0 < t < \infty , 0 < t < \infty , 0 < t < \infty , 0 < t < \infty , 0 < t < \infty , 0 < t < \infty , 0 < t < \infty , 0 < t < \infty , 0 < t < \infty , 0 < t < \infty , 0 < t < \infty , 0 < t < \infty , 0 < t < \infty , 0 < t < \infty , 0 < t < \infty , 0 < t < \infty , 0 < t < \infty , 0 < t < \infty , 0 < t < \infty , 0 < t < \infty , 0 < t < \infty , 0 < t < \infty , 0 < t < \infty , 0 < t < \infty , 0 < t < \infty , 0 < t < \infty , 0 < t < \infty , 0 < t < \infty , 0 < t < \infty , 0 < t < \infty , 0 < t < \infty , 0 < t < \infty , 0 < t < \infty , 0 < t < \infty , 0 < t < \infty , 0 < t < \infty , 0 < t < \infty , 0 < t < \infty , 0 < t < \infty , 0 < t < \infty , 0 < t < \infty , 0 < t < \infty , 0 < t < \infty , 0 < t < \infty , 0 < t < \infty , 0 < t < \infty , 0 < t < \infty , 0 < t < \infty , 0 < t < \infty , 0 < t < \infty , 0 < t < \infty , 0 < t < \infty , 0 < t < \infty , 0 < t < \infty , 0 < t < \infty , 0 < t < \infty , 0 < t < \infty , 0 < t < \infty , 0 < t < \infty , 0 < t < \infty , 0 < t < \infty , 0 < t < \infty , 0 < t < \infty , 0 < t < \infty , 0 < t < \infty , 0 < t < \infty , 0 < t < \infty , 0 < t < \infty , 0 < t < \infty , 0 < t < \infty , , 0 < t < \infty , 0 < t < \infty , 0 < t < \infty , 0 < t < \infty , 0 < t < \infty , 0 < t < \infty , 0 < t < \infty , 0 < t < \infty , 0 < t < \infty , 0 < t < \infty , 0 < t < \infty , 0 < t < \infty , 0 < t < \infty , 0 < t < \infty , 0 < t < \infty , 0 < t < \infty , 0 < t < \infty , 0 < t < \infty , 0 < Щоб перетворити ці граничні умови на нульові умови, ми (після деяких спроб і помилок) зупинилися на такій формі розв'язку:
\end{verbatim}

(6.7)

\[
u(x, t) = A(t)[1-x/L] + B(t)[x/L] + U(x, t),
\]

\begin{verbatim}
де функції A\left(t\right) і B\left(t\right) обираються так, що «квазі-стаціонарна» частина розв'язку (6.7)
\end{verbatim}

(6.8)

\[
S(x, t) = A(t)[1 - x/L] + B(t)[x/L]
\]

\begin{verbatim}
задовольняло граничні умови початкової задачі. У цьому випадку функція U\left(x,\ t\right) задовольнятиме однорідні граничні умови.
\end{verbatim}

Підставляючи функцію \(S(x, t)\) у граничні умови

\[
S(0, t) = g_1(t),
\]

\(S_x\)(L, t) + hS(L, t) = \(g_2\)(t) призводить до двох рівнянь, з яких
можна визначити \(A(t)\) і \(B(t)\). У результаті ми отримуємо

(6.9)

\[
A(t) = g_1(t), B(t) = \frac{g_1(t) + Lg_2(t)}{1 + Lh}.
\]

Отже.

\[
u(x, t) = g_1(t) [1 - x/L] + \frac{g_1(t) + Lg_2(t)}{1 + Lh} [x/L] + U(x, t).
\]

Якщо ми підставимо ці вирази для you(x, t) у початкову задачу (це
просять читача), отримаємо нову задачу для невідомої функції
\(U(x, t)\):

\begin{verbatim}
\begin{array}{llllllllllllllllllllllllllllllllllll
\end{verbatim}

Тепер ми стикаємося з новою проблемою однорідних граничних умов (на
жаль, рівняння стало гетерогенним). Ми не можемо розв'язати цю задачу
методом розділення змінних, але якщо читач погоджується чекати, то за
кілька лекцій він ознайомиться з розв'язком цієї задачі методом цілісних
перетворень і розкладання за власними функціями.

\section{ЗАУВАЖЕННЯ
.}\label{ux437ux430ux443ux432ux430ux436ux435ux43dux43dux44f-.}

\begin{enumerate}
\def\labelenumi{\arabic{enumi}.}
\item
  Як зазначалося раніше, мета цієї лекції --- навчити читача, як
  перетворювати задачі з неунірольними граничними умовами на задачі з
  нульовими граничними умовами. Після цього може статися, що нове
  диференціальне рівняння в похідних похідних стане однорідним. Якщо це
  станеться, ми можемо вважати себе щасливцями, оскільки в цьому випадку
  проблему можна розв'язати методом розділення змінних (див. перший
  приклад з цієї лекції). Якщо перетворене рівняння виявиться
  гетерогенним, тоді нову задачу потрібно розв'язати іншим методом.
\item
  Лінійні неоднорідні граничні умови найзагальної форми
\end{enumerate}

\[
\alpha_1 u_x(0, t) + \beta_1 u(0, t) = g_1(t),
\]

\begin{verbatim}
\alpha_2 u_x(L, t) + \beta_2 u(L, t) = g_2(t) Ви також можете перетворити умови на нульові граничні умови, використовуючи метод, описаний у другому прикладі. Звісно, нове рівняння, ймовірно, буде гетерогенним.
\end{verbatim}

\begin{itemize}
\tightlist
\item
  \begin{enumerate}
  \def\labelenumi{\arabic{enumi}.}
  \setcounter{enumi}{2}
  \tightlist
  \item
    Деякі методи розв'язку не містять вимог щодо унікальності граничних
    умов і, отже, не потребують попередніх перетворень. Пізніше, коли ми
    ознайомимося з перетворенням Лапласа, побачимо, що нульові граничні
    умови є необов'язковими (у випадку нульових граничних умов розв'язок
    дещо простіший, ніж в інших випадках).
  \end{enumerate}
\item
  \begin{enumerate}
  \def\labelenumi{\arabic{enumi}.}
  \setcounter{enumi}{3}
  \tightlist
  \item
    Для граничних умов типу
  \end{enumerate}
\end{itemize}

\[
u(0, t) = g_1(t),
\]

\(u(L, t)\) = \(g_2\)(t) Метод, описаний у другому прикладі, призводить
до наступної форми розв'язання задач:

\[
u(x, t) = \left\{ g_1(t) + \frac{x}{L} \left[ g_2(t) - g_1(t) \right] \right\} + U(x, t).
\]

\section{ЗАВДАННЯ}\label{ux437ux430ux432ux434ux430ux43dux43dux44f-5}

\begin{enumerate}
\def\labelenumi{\arabic{enumi}.}
\tightlist
\item
  Знайти розв'язок наступної змішаної задачі:
\end{enumerate}

\[
\begin{array}{lll} (\text{УЧП}) & u_t = \alpha^2 u_{xx} + x \cos t, & 0 < x < 1, \\ (\Gamma \text{У}) & \begin{cases} u \ (0, \ t) = 1, \\ u_x \ (1, \ t) + hu \ (1, \ t) = 1, \end{cases} & 0 < t < \infty, \\ (\text{HY}) & u \ (x, \ 0) = \sin (\pi x), & 0 \leqslant x \leqslant 1, \end{array}
\]

попередньо перетворюючи його на задачу з нульовими граничними умовами.
Чи відповідає рішення, яке ви отримуєте, вашим задумам щодо цієї
проблеми?

Завдання

(PPP)

\[
u_t = u_{xx}, \quad 0 < x < 1,
\]

(GI)

\[
\begin{cases} u(0, t) = 0, & 0 < t < \infty, \\ u(1, t) = 1, & 0 < t < \infty, \end{cases}
\]

\begin{verbatim}
(НУ) u(x, 0) = x^2, \quad 0 \le x \le 1, Перетворіть у задачу з нульовими граничними умовами і розв'яжіть її. Як виглядатиме рішення у різних
\end{verbatim}

Миті часу? Чи відповідає це вашим інтуїтивним ідеям? Що таке стаціонарне
рішення? Як виглядає перехідне рішення?

\begin{enumerate}
\def\labelenumi{\arabic{enumi}.}
\setcounter{enumi}{2}
\tightlist
\item
  Завдання
\end{enumerate}

\[
\begin{aligned}
(\text{УЧП})\quad & u_t = u_{xx}, && 0 < x < 1,\ 0 < t < \infty, \\
(\text{ГУ})\quad & u_x(0,t) = 0,\quad u_x(1,t) + h u(1,t) = 1, && 0 < t < \infty, \\
(\text{НУ})\quad & u(x,0) = \sin(\pi x), && 0 \leqslant x \leqslant 1.
\end{aligned}
\]

Перетворимо на задачу з нульовими граничними умовами. Чи буде нове
рівняння однорідним?

Лекція 7

\section{РОЗВ'ЯЗАННЯ СКЛАДНІШИХ ЗАДАЧ ШЛЯХОМ РОЗДІЛЕННЯ
ЗМІННИХ}\label{ux440ux43eux437ux432ux44fux437ux430ux43dux43dux44f-ux441ux43aux43bux430ux434ux43dux456ux448ux438ux445-ux437ux430ux434ux430ux447-ux448ux43bux44fux445ux43eux43c-ux440ux43eux437ux434ux456ux43bux435ux43dux43dux44f-ux437ux43cux456ux43dux43dux438ux445}

МЕТА ЛЕКЦІЇ: Показати, як більш складні проблеми теплопровідності можна
розв'язати методом розділення змінних. У цій лекції ми головним чином
розглядаємо такі задачі, розв'язання яких дозволить читачеві краще
ознайомитися з методом ділення змінних. Сподіваємося, що читач зможе
екстраполювати викладені тут ідеї для власних завдань.

У тій же лекції читач ознайомиться з задачею за її власними функціями,
або, як її ще називають, проблемою Штурма-Ліувіля. Обговорюються деякі
властивості проблеми Штурма-Ліувіля.

Як згадувалося вище, мета цієї лекції --- розв'язати змішану задачу
шляхом розділення змінних. Можливо, саме таким методом читач вирішить
свою проблему. Водночас ми сподіваємося, що він зможе перенести метод,
розглянутий у цій лекції, для власних завдань.

Почнемо з одномірної задачі теплопровідності, в якій одна з граничних
умов містить похідну.

\section{Задача теплопровідності з похідною в граничній
умові}\label{ux437ux430ux434ux430ux447ux430-ux442ux435ux43fux43bux43eux43fux440ux43eux432ux456ux434ux43dux43eux441ux442ux456-ux437-ux43fux43eux445ux456ux434ux43dux43eux44e-ux432-ux433ux440ux430ux43dux438ux447ux43dux456ux439-ux443ux43cux43eux432ux456}

Розглянемо пристрій, показаний на рис. 7.1. Припустимо, що верхній
кінець стрижня підтримується при нульовій температурі, тобто \(u(0, t)\)
= 0, а нижній кінець занурений у воду, температура якої також стала і
дорівнює нулю (у цьому випадку --- нуль Р и с

\pandocbounded{\includegraphics[keepaspectratio]{images/_page_50_Picture_4.jpeg}}

. 7.1. Схематичне представлення зсувної задачі.

\pandocbounded{\includegraphics[keepaspectratio]{images/_page_50_Figure_6.jpeg}}

Рис. 7.2. Природа кривих, які задовольняють

\[
(\Gamma Y) = \begin{cases} u(0, t) = 0, \\ u_x(1, t) = -hu(1, t). \end{cases}
\]

відповідає певній еталонній температурі). Теплообмін між стрижнем і
водою відбувається згідно із законом Ньютона, який записується у x=1 як
наступна гранична умова:

\[
u_x(1, t) = -hu(1, t).
\]

Тепер припустимо, що початкова температура стрижня дорівнює \(u(x, 0)\)
= x, але одразу при t \textgreater{} 0 включаємо граничні умови. Щоб
визначити температуру в стрижні, необхідно розв'язати проблему зміщення

\begin{verbatim}
\begin{aligned}
(\text{УЧП})\quad & u_t = \alpha^2 u_{xx}, && 0 < x < 1,\ 0 < t < \infty, \\
(7.1)\ (\text{ГУ})\quad & u(0,t) = 0,\quad u_x(1,t) + h u(1,t) = 0, && 0 < t < \infty, \\
(\text{НУ})\quad & u(x,0) = x, && 0 \leq x \leq 1.
\end{aligned}
\end{verbatim}

Щоб використати метод розділення змінних, розбиємо процес розв'язання на
кілька етапів.

\textbf{КРОК} 1. (Перетворення диференціального рівняння в похідних у
два звичайних диференціальних рівняння)

\begin{verbatim}
Підставляючи u(x, t) = X(x)T(t) у рівняння (7.1), отримуємо XT' = \alpha^2 X''T Поділивши інші частини цього рівняння на \alpha^2 XT отримуємо
\end{verbatim}

\[
\frac{T'}{\alpha^2 T} = \frac{X''}{X} .
\]

\begin{verbatim}
Оскільки ліва частина останнього рівняння залежить лише від часу, а права — лише від x (і оскільки x і t незалежні), обидві сторони рівняння мають бути константами. Прирівнюючи обидві частини до однієї сталої \mu , отримуємо
\end{verbatim}

(7.2)

\[
T' - \mu \alpha^{2} T = 0, X'' - \mu X = 0.
\]

Процес розщеплення змінних завершено.

КРОК 2. (Визначення константи відокремлення)

\begin{verbatim}
По-перше, константа і не повинна бути додатною, інакше ^{1} ) функція T\left( t\right) буде зростати експоненційно, що призведе до нескінченного підвищення температури u=XT. Очевидно, що таке рішення не має фізичного значення.
\end{verbatim}

\begin{verbatim}
По-друге, припустимо, що \mu = 0 . Тоді
\end{verbatim}

\[
X'' = 0
\]

і, отже.

\[
X(x) = A + \hat{B}x.
\]

Однак, оскільки граничні умови вимагають, що

\[
u(0, t) = X(0) T(t) = 0,
\]

\(u(1, t)\) = (hX(1) + X'(1)) \(T(t)\) = 0, Ми одразу дійшли висновку,
що

\[
X(0) = 0
\Rightarrow A = 0 X'(1) + hX(1) = 0 \Rightarrow B = 0 .
\]

\begin{verbatim}
Це означає, що u(x, t) \equiv 0 . Іншими словами, лише нульове розв'язання відповідає нульовій константі розділення. Ми відкинемо його, оскільки нас цікавлять ненульові розв'язки.
\end{verbatim}

\begin{verbatim}
Якщо, нарешті, \mu < 0 , можна поставити \mu = -\lambda^2 і переписати рівняння (7.2) у вигляді
\end{verbatim}

\[
T' + \lambda^2 \alpha^2 T = 0,
X' + \lambda^2 X = 0.
\]

1) Розв'язок \(X(x)\) відповідної задачі з граничним значенням буде лише
тривіальним, \(X(x)\) = 0. --- Ед. Ед.

Загальні розв'язки цих рівнянь легко знайти:

\[
T(t) = Ae^{-(\lambda \alpha)^2 t},
X(x) = B\sin(\lambda x) + C\cos(\lambda x).
\]

Отже, будь-яка функція провідника

(7.3)

\[
u(x, t) = e^{-(\lambda \alpha)^2 t} [A \sin(\lambda x) + B \cos(\lambda x)]
\]

\begin{verbatim}
при будь-яких значеннях \lambda A і B задовольняє рівняння (читач може перевірити це самостійно простою заміною).
\end{verbatim}

Тепер нам потрібно з'ясувати, які з цих функцій задовольняють граничним
умовам

(7.4)

\[
u(0, t) = 0, u_x(1, t) + hu(1, t) = 0.
\]

\begin{verbatim}
Підстановка розв'язків виду (7.3) у граничні умови (7.4) дозволяє знайти відношення, яким мають задовольняти значення \lambda , \Lambda та B:
\end{verbatim}

\[
Be^{-(\lambda\alpha)^2 t} = 0 \Rightarrow B = 0,
\]

\begin{verbatim}
A\lambda e^{-(\lambda\alpha)^2 t} \cos \lambda + hAe^{-(\lambda\alpha)^2 t} \sin \lambda = 0. Після простих алгебранхічних перетворень останнє рівняння зводиться до вигляду
\end{verbatim}

\[
tg \lambda = -\lambda/\hbar
\]

.

\begin{verbatim}
Іншими словами, щоб визначити \lambda , потрібно знайти точки перетину графів функцій tg \lambda і — \lambda/h (рисунок 7.3).
\end{verbatim}

\begin{verbatim}
Числові значення значень \lambda_1 , \lambda_2 , ... можна знайти за допомогою комп'ютера, якщо задано значення h Р и с
\end{verbatim}

\pandocbounded{\includegraphics[keepaspectratio]{images/_page_52_Figure_15.jpeg}}

\begin{verbatim}
. 7.3. Перетин графів функцій tg(\lambda) і -\lambda/h .
\end{verbatim}

викликати власні значення задачі граничного значення

(7.5)

\[
X'' + \lambda^{2}X = 0, \begin{cases} X(0) = 0 \\ X'(1) + hX(1) = 0. \end{cases}
\]

\begin{verbatim}
Іншими словами, з цими значеннями параметра \lambda існують нетривіальні розв'язки задачі (7.5). Власні значення задачі (7.5) визначаються як корені рівняння \log \lambda = -\lambda/h . Їхні числові значення знаходилися на комп'ютері для випадку h = 1. Таблиця 7.1 показує результати розрахунків перших кількох власних значень.
\end{verbatim}

\begin{verbatim}
| n | \lambda_B | rt . | \lambda_n |
\end{verbatim}

\textbar-------------\textbar----------------------\textbar--------\textbar----------------\textbar{}

123 \textbar{} 2,024,917,98 \textbar{} 45 \textbar{} 11,0814,20
\textbar{}

\begin{verbatim}
Розв'язки задачі (7.5), що відповідають власним значенням \lambda_n , називаються власними функціями X_n(x) . Для розглянутого завдання власні функції виглядають так:
\end{verbatim}

\[
X_n(x) = \sin(\lambda_n x)
\]

і зображені на рис. 7.4.

\pandocbounded{\includegraphics[keepaspectratio]{images/_page_53_Figure_10.jpeg}}

Рис. 7.4. Власні функції \(X_n\)(x) завдання (7.5) при h=1; Усі власні
функції задовольняють граничні умови.

КРОК 3, (Пошук фундаментальних рішень)

Тепер у нас є нескінченний набір функцій (фундаментальних розв'язків) у
нашому розпорядженні

\[
u_n(x, t) = X_n(x) T_n(t) = e^{-(\lambda_n x)^2 t} \sin(\lambda_n x),
\]

кожна з яких задовольняє рівняння та граничні умови. Залишився ще один
крок --- загалом

\[
u(x, t) = \sum_{n=1}^{\infty} a_n X_n(x) T_n(t) =
= \sum_{n=1}^{\infty} a_n e^{-(\lambda_n c t^2 t)} \sin(\lambda_n x)
\]

Вибрати коефіцієнти \(a_n\) так, щоб функція \(u(x, t)\) задовольняла
початкову умову \(u(x, 0)\) = x (у цьому випадку функція \(u(x, t)\)
задовольнятиме рівняння та граничні умови, оскільки, по-перше, кожен
член суми їх задовольняє, а по-друге, і рівняння, і граничні умови є
лінійними та однорідними).

Отже, ймовірність \(a_n\) слід визначати з цієї умови

(7 6)

\[
u(x, 0) = x = \sum_{n=1}^{\infty} a_n \sin(\lambda_n x).
\]

Ми маємо зробити останній крок.

КРОК 4. (Розклад початкової умови на ряд за власними функціями)

Щоб знайти коефіцієнти \(a_n\) у розкладі (7,6) на власні функції,
помножте обидві частини цього співвідношення на \(\sin(\lambda_m x)\) і
інтегруйте від 0 до 1. У результаті ми отримуємо

\[
\int_{0}^{2} \xi \sin(\lambda_{m} \xi) d\xi = \sum_{n=1}^{\infty} a_{n} \int_{0}^{1} \sin(\lambda_{n} \xi) \sin(\lambda_{m} \xi) d\xi =
= a_{m} \int_{0}^{1} \sin^{2}(\lambda_{m} \xi) d\xi =
= a_{m} \left( \frac{\lambda_{m} - \sin \lambda_{m} \cos \lambda_{m}}{2\lambda_{m}} \right).
\]

Розв'язавши останнє відношення відносно \(a_m\) і замінивши m на n,
отримуємо бажаний результат:

(7.7)

\[
a_n = \frac{2\lambda_n}{(\lambda_n - \sin \lambda_n \cos \lambda_n)} \int_0^{\infty} \xi \sin (\lambda_n \xi) d\xi.
\]

Таблиця 7.2.

\section{\texorpdfstring{Шанси \(a_n\) у формулі
(7.8)}{Шанси a\_n у формулі (7.8)}}\label{ux448ux430ux43dux441ux438-a_n-ux443-ux444ux43eux440ux43cux443ux43bux456-7.8}

{\def\LTcaptype{none} % do not increment counter
\begin{longtable}[]{@{}llll@{}}
\toprule\noalign{}
n & an & n & \(a_n\) \\
\midrule\noalign{}
\endhead
\bottomrule\noalign{}
\endlastfoot
123 & 0,240,22-0,03 & 45 & -0.11-0.09  \\
\end{longtable}
}

Отже, розв'язок задачі (7.1) записується у вигляді

(7.8)

\[
u(x, t) = \sum_{n=1}^{\infty} a_n e^{-i\lambda_n \alpha_n^{2} t} \sin(\lambda_n x),
\]

де коефіцієнти \(a_n\) задаються формулами (7.7). Перші п'ять
коефіцієнтів \(a_n\) нашої задачі були знайдені чисельно. Їхні значення
наведені в Таблиці 1. 7.2.

\pandocbounded{\includegraphics[keepaspectratio]{images/_page_55_Figure_8.jpeg}}

Рис. 7.5. Графічне представлення функції (7.8); межа умовно
задовольняється у всі моменти часу.

Відповідно, перші три члени розв'язання змішаної задачі

(7.9)

\[
\begin{aligned} u_t &= u_{xx}, & 0 < x < 1, & 0 < t < \infty, \\ u_t &= 0, & 0 < t < \infty, \\ u_x &= 0, & 0 < t < \infty, \\ u_x &= 0, & 0 < t < \infty, \\ u_x &= 0, & 0 < t < \infty, \\ u_x &= 0, & 0 < t < \infty, \\ u_x &= 0, & 0 < t < \infty, \\ u_x &= 0, & 0 < t < \infty, \\ u_x &= 0, & 0 < t < \infty, \\ u_x &= 0, & 0 < t < \infty, \\ u_x &= 0, & 0 < t < \infty, \\ u_x &= 0, & 0 < t < \infty, \\ u_x &= 0, & 0 < t < \infty, \\ u_x &= 0, & 0 < t < \infty, \\ u_x &= 0, & 0 < t < \infty, \\ u_x &= 0, & 0 < t < \infty, \\ u_x &= 0, & 0 < t < \infty, \\ u_x &= 0, & 0 < t < \infty, \\ u_x &= 0, & 0 < t < \infty, \\ u_x &= 0, & 0 < t < \infty, \\ u_x &= 0, & 0 < t < \infty, \\ u_x &= 0, & 0 < t < \infty, \\ u_x &= 0, & 0 < t < \infty, \\ u_x &= 0, & 0 < t < \infty, \\ u_x &= 0, & 0 < t < \infty, \\ u_x &= 0, & 0 < t < \infty, \\ u_x &= 0, & 0 < t < \infty, \\ u_x &= 0, & 0 < t < \infty, \\ u_x &= 0, & 0 < t < \infty, \\ u_x &= 0, & 0 < t < \infty, \\ u_x &= 0, & 0 < t < \infty, \\ u_x &= 0, & 0 < t < \infty, \\ u_x &= 0, & 0 < t < \infty, \\ u_x &= 0, & 0 < t < \infty, \\ u_x &= 0, & 0 < t < \infty, \\ u_x &= 0, & 0 < t < \infty, \\ u_x &= 0, & 0 < t < \infty, \\ u_x &= 0, & 0 < t < \infty, \\ u_x &= 0, & 0 < t < \infty, \\ u_x &= 0, & 0 < t < \infty, \\ u_x &= 0, & 0 < t < \infty, \\ u_x &= 0, & 0 < t < \infty, \\ u_x &= 0, & 0 < t < \infty, \\ u_x &= 0, & 0 < t < \infty, \\ u_x &= 0, & 0 < t < \infty, \\ u_x &= 0, & 0 < t < \infty, \\ u_x &= 0, & 0 < t < \infty, \\ u_x &= 0, & 0 < t < \infty, \\ u_x &= 0, & 0 < t < \infty, \\ u_x &= 0, & 0 < t < \infty, \\ u_x &= 0, & 0 < t < \infty, \\ u_x &= 0, & 0 < t < \infty, \\ u_x &= 0, & 0 < t < \infty, \\ u_x &= 0, & 0 < t < \infty, \\ u_x &= 0, & 0 < t < \infty, \\ u_x &= 0, & 0 < t < \infty, \\ u_x &= 0, & 0 < t < \infty, \\ u_x &= 0, & 0 < t < \infty, \\ u_x &= 0, & 0 < t < \infty, \\ u_x &= 0, & 0 < t < \infty, \\ u_x &= 0, & 0 < t < \infty, \\ u_x &= 0, & 0 < t < \infty, \\ u_x &= 0, & 0 < t < \infty, \\ u_x &= 0, & 0 < t < \infty, \\ u_x &= 0, & 0 < t < \infty, \\ u_x &= 0, & 0 < t < \infty, \\ u_x &= 0, & 0 < t < \infty, \\ u_x &= 0, & 0 < t < \infty, \\ u_x &= 0, & 0 < t < \infty, \\ u_x &= 0, & 0 < t < \infty, \\ u_x &= 0, & 0 < t < \infty, \\ u_x &= 0, & 0 < t < \infty, \\ u_x &= 0, & 0 < t < \infty, \\ u_x &= 0, & 0 < t < \infty, \\ u_x &= 0, & 0 < t < \infty, \\ u_x &= 0, & 0 < t < \infty, \\ u_x &= 0, & 0 < t < \infty, \\ u_x &= 0, & 0 < t < \infty, \\ u_x &= 0, & 0 < t < \infty, \\ u_x &= 0, & 0 < t < \infty, \\ u_x &= 0, & 0 < t \end{aligned}
\]

мають вигляд

\[
u(x, t) = = 0.24e^{-4t} \sin(2.02x) + 0.22e^{-24t} \sin(4.9x) - 0.03e^{-48.3t} \sin(7.98x) + \dots
\]

Фіг. На рисунку 7-5 показані графіки цього розв'язку для кількох точок
часу. Рекомендуємо читачеві відповісти на такі питання: Чи відповідає
тип рішення, яке ви отримуєте, з вашими інтуїтивними ідеями? Чи
виконуються граничні умови?

\section{ЗАУВАЖЕННЯ}\label{ux437ux430ux443ux432ux430ux436ux435ux43dux43dux44f-5}

Задача власних значень (7.5) є особливим випадком більш загальної
задачі, яку зазвичай називають проблемою Штурма-Ліувіля:

(7.10)

\[
\begin{aligned} & (\text{ОДУ}) \quad [p(x)y']' - q(x)y + \lambda r(x)y = 0, \quad 0 < x < 1, \\ & (\alpha_1 y(0) + \beta_1 y'(0) = 0, \\ & (\alpha_2 y(1) + \beta_1 y'(1) = 0. \end{aligned}
\]

Якщо розв'язати задачу з однорідними граничними умовами методом
розділення змінних, то функції \(X(x)\) завжди знаходяться як розв'язки
деякої задачі Штурма--Ліувіля. Очевидно, що задача власних значень (7.5)
є особливим випадком задачі (7.10).

Перелічимо основні властивості задачі (7.10), які були встановлені
Штурмом і Ліувілем за певних природних обмежень.

функції \(p(x)\), \(q(x)\) та \(r(x)\).

\begin{enumerate}
\def\labelenumi{\arabic{enumi}.}
\tightlist
\item
  Існує нескінченна послідовність власних значень, що задовольняють
  нерівності
\end{enumerate}

\[
\lambda_1 < \lambda_2 < \lambda_3 < \ldots < \lambda_n < \ldots \to \infty
\]

.

\begin{enumerate}
\def\labelenumi{\arabic{enumi}.}
\setcounter{enumi}{1}
\item
  Кожне власне значення \(\lambda_n\) відповідає одній (точній до
  константного множника) ненульовій власній функції \(y_n(x)\) .
\item
  Якщо \(y_n(x)\) і \(y_m(x)\) дві різні власні функції, що відповідають
  різним власним значенням \(\lambda_n\) і \(\lambda_m\) , то вони є
  ортогональними з вагою \(r(x)\) на сегменті {[}0, 1{]}, тобто вони
  задовольняють це відношення
\end{enumerate}

\[
\int_{0}^{1} r(x) y_{n}(x) y_{m}(x) dx = 0.
\]

Більше деталей про проблему Штурма-Ліувіля можна знайти в літературі,
наведеній наприкінці книги.

1) Див., наприклад, Тіхоног А. Н., Самарський А. А. Рівняння
математичної фізики. --- Москва: Наука, 1972. --- ред. Ед.

\section{ЗАВДАННЯ}\label{ux437ux430ux432ux434ux430ux43dux43dux44f-6}

\begin{enumerate}
\def\labelenumi{\arabic{enumi}.}
\tightlist
\item
  Розв'язати наступну задачу для рівняння теплопровідності методом
  розкладу
\end{enumerate}

(PPP)

\[
u_1 = u_{xx}, \quad 0 < x < 1, \quad 0 < t < \infty,
\]

(GI)

\[
\begin{cases} u(0, t) = 0, \\ u_x(1, t) = 0, \\ u(x, 0) = x, \quad 0 \le x \le 1. \end{cases}
\]

Чи відповідає рішення, яке ви знаходите, з інтонативними
представленнями? Що таке стаціонарне рішення?

\begin{enumerate}
\def\labelenumi{\arabic{enumi}.}
\setcounter{enumi}{1}
\tightlist
\item
  Знайдіть власні цінності та функції проблеми Ассолта-Ліувіля
\end{enumerate}

(ODY)

\[
X' + \lambda X = 0, \quad 0 < x < 1,
\]

(FY)

\[
\begin{cases} X(0) = 0, \\ X'(1) = 0. \end{cases}
\]

Задайте специфічні вирази для функцій \(p(x)\), \(q(x)\) та \(r(x)\) із
загальної задачі Штурма--Ліувіля для цього випадку.

\begin{enumerate}
\def\labelenumi{\arabic{enumi}.}
\setcounter{enumi}{2}
\tightlist
\item
  Розв'язати задачу теплопровідності для планки з термоізольованими
  межами
\end{enumerate}

(PPP)

\[
u_t = u_{xx}, \quad 0 < x < 1, \quad 0 < t < \infty,
\]

(GI)

\[
\begin{cases} u_x (0, t) = 0, \\ u_x (1, t) = 0, \end{cases} \quad 0 < t < \infty,
\]

(ГУ) \(u(x, 0) = x, \quad 0 \le x \le 1.\) Чи відповідає рішення, яке ви
отримуєте, з вашою інтерпретацією цієї проблеми? Що таке стаціонарне
рішення? Яке його фізичне значення?

\begin{enumerate}
\def\labelenumi{\arabic{enumi}.}
\setcounter{enumi}{3}
\tightlist
\item
  Знайдіть власні цінності та функції завдань
\end{enumerate}

(ОДА)

\[
X'' + \lambda X = 0, \quad 0 < x < 1,
\]

(GI)

\[
\begin{cases} X'(0) = 0, \\ X'(1) = 0. \end{cases}
\]

\#\# ПЕРЕТВОРЕННЯ СКЛАДНИХ РІВНЯНЬ У ПРОСТУ ФОРМУ

МЕТА ЛЕКЦІЇ Показати, як PPP для функції \(u(x, t)\) можна перетворити
на простіший PPP для нової функції \(w(x, t)\). Вибір правильної
трансформації зазвичай базується на ненавмисному намірі. У нашій лекції

\[
u_t = \alpha^2 u_{xx} - \beta u,
u_t = \alpha^2 u_{xx} - v u_x, \quad (v = \text{const})
\]

переходимо до найпростішого рівняння теплопровідності

\[
w_1 = \alpha^2 w_{xx}
\]

внаслідок відповідних перетворень

\[
u(x, t) = e^{-\beta t} w(x, t)
u(x, t) = e^{\epsilon t (1 - \delta t - 2) 2\alpha^2} w(x, t).
\]

Розв'язавши просте рівняння теплопровідності за функцією \(\omega(x,t)\)
, ми можемо знайти розв'язки початкових рівнянь, перетворивши функцію
\(\omega(x,t)\) відповідно до цих формул (звісно, ці перетворення слід
застосовувати і до початкових і граничних умов).

Після прочитання двох останніх лекцій читач може створити враження, що
метод розділення змінних застосовується \emph{лише} до рівняння

\[
u_1 = \alpha^2 u_{xx}
\]

Насправді рівняння теплопотоку є найпростішим, але не єдиним
параболічним рівнянням, яке можна розв'язати методом розділення змінних.
Як уже згадувалося, якщо рівняння є лінійним і однорідним, його можна
розв'язати, поділивши змінні. Наприклад, двовимірна задача
теплопровідності в колі описується рівнянням

\[
u_t = \alpha^2 \left[ u_{rr} + \frac{1}{r} u_r + \frac{1}{r^2} u_{\theta\theta} \right],
\]

І хоча це рівняння з \emph{змінними коефіцієнтами}, його можна звести до
трьох звичайних диференціальних рівнянь.

У цій лекції ми покажемо читачеві, що іноді, замість того, щоб
безпосередньо розв'язувати диференціальне рівняння в похідних, краще
перетворити його на простіше рівняння, яке можна розв'язати методом
розділення змінних або іншим методом. Наступний приклад ілюструє цей
підхід.

\section{Перетворення проблеми з перенесенням бічної поверхні у проблему
з термічно ізольованою бічною
поверхнею}\label{ux43fux435ux440ux435ux442ux432ux43eux440ux435ux43dux43dux44f-ux43fux440ux43eux431ux43bux435ux43cux438-ux437-ux43fux435ux440ux435ux43dux435ux441ux435ux43dux43dux44fux43c-ux431ux456ux447ux43dux43eux457-ux43fux43eux432ux435ux440ux445ux43dux456-ux443-ux43fux440ux43eux431ux43bux435ux43cux443-ux437-ux442ux435ux440ux43cux456ux447ux43dux43e-ux456ux437ux43eux43bux44cux43eux432ux430ux43dux43eux44e-ux431ux456ux447ux43dux43eux44e-ux43fux43eux432ux435ux440ux445ux43dux435ux44e}

Розглянемо наступне завдання:

(8.1)

\[
u_{t} = \alpha^{2}u_{xx} - \beta u, \quad 0 < x < 1, \quad 0 < t < \infty,
\begin{cases} u(0, t) = 0, \\ u(1, t) = 0, \\ u(x, 0) = \varphi(x), \quad 0 \le x \le 1, \end{cases}
\]

(8.1)

\[
(FY) \quad \begin{cases} u(0, t) = 0, \\ u(x, 0) = \varphi(x), \quad 0 \le x \le 1, \end{cases}
\]

\begin{verbatim}
де термін -\beta u описує потік тепла через бічну поверхню стрижня (рис. 8.1).
\end{verbatim}

\pandocbounded{\includegraphics[keepaspectratio]{images/_page_59_Figure_7.jpeg}}

\begin{verbatim}
Рис. 81. Задача теплопровідності, описана рівнянням u_t = \alpha^t u_{\lambda\lambda} - \beta u : \alpha — це втрата гепле через поверхню Боха; \delta – дифузія вздовж стрижня.
\end{verbatim}

\begin{verbatim}
Наша мета — ввести нову температурну w\left(x,\ t\right) замість u\left(x,\ t\right) , щоб рівняння для w\left(x,\ t\right) стало простішим за початкове рівняння
\end{verbatim}

\[
u_t = \alpha^2 u_{xx} - \beta u
\]

.

Це досить загальний метод розв'язання диференціальних рівнянь у
частинних похідних, але вибір відповідного перетворення базується на
інтуїтивних уявленнях про поведінку розв'язків початкового рівняння.
Наприклад, у задачі (8.1) температура \(u(x, t)\) у кожній точці \(x_0\)
змінюється внаслідок двох факторів:

\begin{verbatim}
- 1) дифузія тепла вздовж стрижня (через компонентний \alpha^2 u_{xx} ),
\end{verbatim}

\begin{itemize}
\tightlist
\item
  \begin{enumerate}
  \def\labelenumi{\arabic{enumi})}
  \setcounter{enumi}{1}
  \tightlist
  \item
    теплопередача через бічну поверхню стрижня (через член RI).
  \end{enumerate}
\end{itemize}

\begin{verbatim}
Найважливішим елементом нашого міркування є наступне: якщо не було дифузії вздовж стрижня ( \alpha = 0 ),
\end{verbatim}

тоді температура в кожній точці стрижня експоненційно знижується до нуля
згідно із законом

\[
u(x_0 t) = u(x_0, 0) e^{-\beta t}
\]

.

Виходячи з цього спостереження, запитаємо, чи можна розв'язок проблеми
(8.1) представити як утворення двох факторів

\[
(8.2) u(x, t) = e^{-\beta t}w(x, t)
\]

або:

\begin{verbatim}
\begin{pmatrix}
\end{verbatim}

Температура елемента з \textbackslash{}

Неізольовані боком

Поверхня

де \(w(x, t)\) описує розподіл температури, що виникає лише через
дифузію. Давайте подивимось, що станеться, якщо підставити цей вираз для
температури у задачу (8.1). Після виконання відповідних розрахунків
(читач може зробити їх сам), ми отримаємо

\begin{verbatim}
Нова задача майже точно така ж, як і оригінальна, але більше не має додатку 6u. Тому замість розв'язання задачі (8.1) ми розв'язуємо перетворену задачу (8.3), а потім множимо отриманий розв'язок w(x, t) на e^{-\beta t} . У нашому випадку розв'язок задачі (8.3) вже було отримано раніше методом розділення змінних, тому ми можемо одразу записати
\end{verbatim}

(8.4)

\[
w(x, t) = \sum_{n=1}^{\infty} a_n e^{-(n\pi\alpha)^{\frac{n}{2}}t} \sin(n\pi x),
a_n = 2 \int_0^{\infty} \phi(\xi) \sin(n\pi \xi) d\xi,
\]

і, отже, розв'язок початкової задачі (8.1) має вигляд

\[
u(x, t) = e^{-\beta t}w(x, t).
\]

Нижче, у розділі «Нотатки», наведено ще один приклад проблеми,
розв'язаної цим методом.

\section{ПРИМІТКИ}\label{ux43fux440ux438ux43cux456ux442ux43aux438}

\begin{enumerate}
\def\labelenumi{\arabic{enumi}.}
\tightlist
\item
  Вирішити проблему
\end{enumerate}

\[
\begin{aligned}
u_t &= u_{xx} - u, && 0 < x < 1,\ 0 < t < \infty, \\
u(0,t) &= 0,\quad u(1,t) = 0, && 0 < t < \infty, \\
u(x,0) &= \sin(\pi x) + 0.5\sin(3\pi x), && 0 \leqslant x \leqslant 1.
\end{aligned}
\]

За методом, описаним вище, ми

\begin{itemize}
\tightlist
\item
  \begin{enumerate}
  \def\labelenumi{\alph{enumi})}
  \tightlist
  \item
    тимчасово відкинути конвективний член u у початковому рівнянні.
  \end{enumerate}
\item
  \begin{enumerate}
  \def\labelenumi{\alph{enumi})}
  \setcounter{enumi}{1}
  \tightlist
  \item
    записати розв'язок змішаної задачі без конвективного члена
  \end{enumerate}
\end{itemize}

\[
u(x, t) = e^{-\pi^2 t} \sin(\pi x) + 0.5e^{-(3\pi)^2 t} \sin(3\pi x),
\]

\begin{enumerate}
\def\labelenumi{\alph{enumi})}
\setcounter{enumi}{2}
\tightlist
\item
  помножимо це розв'язання на конвективний фактор
  \(e^{-\beta t}=e^{-t}\) отримаємо остаточне розв'язання початкової
  задачі
\end{enumerate}

\[
u(x, t) = e^{-t} \left[ e^{-\pi^2 t} \sin(\pi x) + 0.5e^{-(3\pi)^2 t} \sin(3\pi x) \right].
\]

\begin{enumerate}
\def\labelenumi{\arabic{enumi}.}
\setcounter{enumi}{1}
\tightlist
\item
  Рівняння конвективної дифузії
\end{enumerate}

\[
u_t = \alpha^2 u_{xx} - v u_x
\]

(коефіцієнт і є константою) також можна перетворити у вигляд

\[
w_t = \alpha^2 w_{xx}.
\]

Відповідне перетворення має вигляд

\[
u(x, t) = e^{v(x-vt/2)/2\alpha^2}w(x, t).
\]

Це перетворення підкреслює експоненціальний фактор, який повністю
зумовлений рухом середовища. Зверніть увагу, що цей множник описує
хвилю, що поширюється зліва направо зі швидкістю \(\upsilon/2\) . Читач
матиме можливість скористатися цією трансформацією.

\section{ЗАВДАННЯ}\label{ux437ux430ux432ux434ux430ux43dux43dux44f-7}

\begin{enumerate}
\def\labelenumi{\arabic{enumi}.}
\tightlist
\item
  Розв'язати задачу конвективної дифузії
\end{enumerate}

\begin{verbatim}
\begin{aligned}
(\text{УЧП})\quad & u_t = u_{xx} - u_x, && 0 < x < 1,\ 0 < t < \infty, \\
(\text{ГУ})\quad & u(0,t) = 0,\quad u(1,t) = 0, && 0 < t < \infty, \\
(\text{НУ})\quad & u(x,0) = e^{x/2}, && 0 \leqslant x \leqslant 1.
\end{aligned}
\end{verbatim}

методом перетворення до простішого рівняння. На що схожий розв'язок?
Можна дати таку інтерпретацію цієї задачі: \(u(x,t)\) --- концентрація
деякої речовини в рухомому середовищі, яке рухається зліва направо зі
швидкістю \(v = 1\); початковий розподіл концентрації задається функцією
\(e^{x/2}\), а на межах концентрація підтримується нульовою. Чи
відповідає ваш розв'язок такій інтерпретації?

\begin{enumerate}
\def\labelenumi{\arabic{enumi}.}
\setcounter{enumi}{1}
\tightlist
\item
  Розв'язати задачу
\end{enumerate}

\begin{verbatim}
\begin{aligned}
(\text{УЧП})\quad & u_t = u_{xx} - u + x, && 0 < x < 1,\ 0 < t < \infty, \\
(\text{ГУ})\quad & u(0,t) = 0,\quad u(1,t) = 1, && 0 < t < \infty, \\
(\text{НУ})\quad & u(x,0) = 0, && 0 \leqslant x \leqslant 1.
\end{aligned}
\end{verbatim}

\begin{itemize}
\tightlist
\item
  а) перетворивши неоднорідні граничні умови на однорідні;
\item
  б) перетворивши вихідне рівняння на нове рівняння без члена \(-u\);
\item
  в) розв'язавши задачу, що виникає після цих перетворень.
\end{itemize}

\begin{enumerate}
\def\labelenumi{\arabic{enumi}.}
\setcounter{enumi}{2}
\tightlist
\item
  Розв'язати задачу
\end{enumerate}

\begin{verbatim}
\begin{aligned}
(\text{УЧП})\quad & u_t = u_{xx} - u, && 0 < x < 1,\ 0 < t < \infty, \\
(\text{ГУ})\quad & u(0,t) = 0,\quad u(1,t) = 0, && 0 < t < \infty, \\
(\text{НУ})\quad & u(x,0) = \sin(\pi x), && 0 \leqslant x \leqslant 1.
\end{aligned}
\end{verbatim}

застосовуючи метод розділення змінних безпосередньо до вихідного
рівняння, не виконуючи жодних попередніх перетворень. Чи збігається
отриманий розв'язок з розв'язком, знайденим раніше після заміни

\[
u(x,t) = e^{-t}w(x,t)
\]

?

\section{РОЗВ'ЯЗАННЯ НЕОДНОРІДНИХ УЧП МЕТОДОМ РОЗКЛАДУ ЗА ВЛАСНИМИ
ФУНКЦІЯМИ}\label{ux440ux43eux437ux432ux44fux437ux430ux43dux43dux44f-ux43dux435ux43eux434ux43dux43eux440ux456ux434ux43dux438ux445-ux443ux447ux43f-ux43cux435ux442ux43eux434ux43eux43c-ux440ux43eux437ux43aux43bux430ux434ux443-ux437ux430-ux432ux43bux430ux441ux43dux438ux43cux438-ux444ux443ux43dux43aux446ux456ux44fux43cux438}

ЦІЛЬ ЛЕКЦІЇ: Навчитися розв'язувати змішані задачі вигляду

\begin{verbatim}
\begin{aligned}
(\text{УЧП})\quad & u_t = \alpha^2 u_{xx} + f(x,t), && 0 < x < 1,\ 0 < t < \infty, \\
(\text{ГУ})\quad & \alpha_1 u_x(0,t) + \beta_1 u(0,t) = 0, && 0 < t < \infty, \\
& \alpha_2 u_x(1,t) + \beta_2 u(1,t) = 0, && 0 < t < \infty, \\
(\text{НУ})\quad & u(x,0) = \varphi(x), && 0 \leqslant x \leqslant 1.
\end{aligned}
\end{verbatim}

Розв'язок неоднорідного УЧП можна шукати у вигляді ряду

\[
u(x,t) = \sum_{n=1}^{\infty} T_n(t) X_n(x).
\]

де \(X_n(x)\) --- власні функції відповідної однорідної задачі

\begin{verbatim}
\begin{aligned}
(\text{УЧП})\quad & u_t = \alpha^2 u_{xx}, && 0 < x < 1,\ 0 < t < \infty, \\
(\text{ГУ})\quad & \alpha_1 u_x(0,t) + \beta_1 u(0,t) = 0, && 0 < t < \infty, \\
& \alpha_2 u_x(1,t) + \beta_2 u(1,t) = 0, && 0 < t < \infty, \\
(\text{НУ})\quad & u(x,0) = \varphi(x), && 0 \leqslant x \leqslant 1.
\end{aligned}
\end{verbatim}

А функції \(T_n(t)\) визначаються як розв'язки деяких звичайних
диференціальних рівнянь.

На лекції 6 ми ознайомилися з тим, як перетворювати неоднорідні граничні
умови на однорідні. На жаль, рівняння залишається гетерогенним, і ми
стоїмо перед завданням

\[
(9.1) \quad u_{t} = \alpha^{2} u_{xx} + f(x, t), \quad 0 < x < 1, \quad 0 < t < \infty,
(9.1) \quad (\Gamma Y) \quad
\begin{cases} \alpha_{1} u_{x}(0, t) + \beta_{1} u(0, t) = 0, \\ \alpha_{2} u_{x}(1, t) + \beta_{2} u(1, t) = 0, \\ u(x, 0) = \varphi(x), \quad 0 \le x \le 1. \end{cases}
\]

Мета нашої лекції --- опанувати метод розв'язання таких задач. Цей метод
аналогічний методу варіації довільних констант.

теорія звичайних диференціальних рівнянь називається методом розкладу за
власними функціями.

Ідея методу надзвичайно проста. Відомо, що розв'язок задачі (9.1) при
\(f(x, t)\) = 0 (так звана відповідна однорідна задача) записується як
ряд

\[
u\left(x,\ t\right)=\sum_{n=1}^{\infty}a_{n}e^{-\left(\lambda_{n}\alpha\right)^{2}}/X_{n}\left(x\right),
\]

де \(X_n(x)\) --- власні функції задачі Штурма-Ліувіля,

(9.2)

\[
X'' + \lambda^2 X = 0, \\ \alpha_1 X'(0) + \beta_1 X(0) = 0, \\ \alpha_2 X'(1) + \beta_2 X(1) = 0.
\]

Запитаймо себе, чи можна також шукати розв'язок гетерогенної проблеми
(9.1) у вигляді серії власних функцій, але більш загальної форми

\[
u(x, t) = \sum_{n=1}^{\infty} T_n(t) X_n(x).
\]

Існують певні фізичні підстави для такого припущення. Насправді терміни
\(f(x,t)\) описують джерело тепла, розташоване всередині стрижня. Це
означає, що залежність рішення від часу не буде виражена показниками

\[
e^{-(\lambda_n\alpha)^2 t}
\]

,

Це стосувалося однорідного завдання, але з деякими іншими функціями.

Щоб показати цей метод у дії і не ускладнювати подання деталями, ми
застосуємо його до розв'язання однієї дуже простої задачі.

\section{Построзв'язання методом демаркації за власними
функціями,}\label{ux43fux43eux441ux442ux440ux43eux437ux432ux44fux437ux430ux43dux43dux44f-ux43cux435ux442ux43eux434ux43eux43c-ux434ux435ux43cux430ux440ux43aux430ux446ux456ux457-ux437ux430-ux432ux43bux430ux441ux43dux438ux43cux438-ux444ux443ux43dux43aux446ux456ux44fux43cux438}

Розглянемо гетерогенну проблему

(9.3) (PY)

\[
\begin{aligned} u_t &= \alpha^2 u_{xx} + f(x, t), & 0 < x < 1, & 0 < t < \infty, \\ u(0, t) &= 0, & 0 < t < \infty, \\ u(1, t) &= 0, & 0 < t < \infty, \\ u(x, 0) &= \varphi(x), & 0 \leqslant x \leqslant 1. \end{aligned}
\]

Процес розв'язання задачі подається у вигляді послідовності наступних
кроків.

КРОК 1. Основна ідея методу полягає в розкладанні густини джерела
\(f(x, t)\) на ряд власними функціями

\[
f(x, t) = f_1(t) X_1(x) + f_2(t) X_2(x) + \dots + f_n(t) X_n(x) + \dots
\]

а також визначити реакції системи \(u_n\left(x,\,t\right)\) на вплив
кожного компонента \(f_n\left(t\right)X_n\left(x\right)\) . Якщо знайти
всі відповіді, рішення проблеми виглядатиме так:

\[
u(x, t) = \sum_{n=1}^{\infty} u_n(x, t).
\]

Головна складність цього методу --- розклад густини джерела на
компоненти \(f_n(t)\,X_n(x)\) . Виявляється, що \(X_n(x)\) факторами в
цій задачі є власними векторами системи Штурма--Ліувіля, які виникають
при розв'язанні відповідної однорідної задачі методом розділення
змінних, а саме

\[
\begin{array}{c} u_t = \alpha^2 u_{xx} \quad \text{(отметим, что } f(x, \ t) = 0), \\ u(0, \ t) = 0, \\ u(1, \ t) = 0, \\ u(x, \ 0) = \varphi(x). \end{array}
\]

Завдання Штурма-Ліувіля в цьому випадку виглядає так:

\[
X'' + \lambda^2 X = 0,
\]

\(X(0) = 0,\) \(X(1) = 0,\) а його розв'язки є функціями

\[
X_n(x) = \sin(n\pi x), \quad n = 1, 2, 3 \dots
\]

Отже, розклад на густину джерела представлений як ряд

\[
(9.5) \ f(x,t) = f_1(t) \sin(\pi x) + f_2(t) \sin(2\pi x) + \dots + f_n(t) \sin(n\pi x) + \dots
\]

\(\mathcal{U}\) нарешті, щоб знайти функції \(f_n(t)\) , помножте обидві
сторони останнього відношення на \(\sin(m\pi x)\) і інтегруйте від 0 до
1 за координатою x.

У результаті ми отримуємо

\[
\int_{0}^{1} f(x, t) \sin(m\pi x) dx = \sum_{n=1}^{\infty} f_{n}(t) \int_{0}^{1} \sin(m\pi x) \sin(n\pi x) dx =
= \frac{1}{2} f_{m}(t)
\]

або (після заміни t на p)

(9.6)

\[
f_{n}(t) = 2 \int_{0}^{1} f(x, t) \sin(n\pi x) dx.
\]

Останнє співвідношення встановлює співвідношення між коефіцієнтами
\(f_n\)(t) і часткою джерела тепла \(f(x, t)\).

КРОК 2. (Пошук \(u_n\)(x, t) = \(X_n\)(x) \(T_n\)(t) реакції на вхідний
вплив \(f_n\)(t) \(X_n\)(x) .)

Замінимо густину джерела на його розклад у рядок

\[
f(x, t) = \sum_{n=1}^{\infty} f_n(t) \sin(n\pi x)
\]

\begin{verbatim}
і намагатися знайти індивідуальні відповіді, тобто функції, T_{\mathbf{n}}(t) у
\end{verbatim}

\[
u(x, t) = \sum_{n=1}^{\infty} T_n(t) \sin(n\pi x).
\]

Оскільки розв'язок задачі має вигляд

\[
u(x, t) = \sum_{n=1}^{\infty} T_n(t) \sin(n\pi x),
\]

потім підставимо цей вираз у початкову задачу

\[
u_{t} = \alpha^{2} u_{xx} + \sum_{n=1}^{\infty} f_{n}(t) \sin(n\pi x),
u(0, t) = 0,
u(1, t) = 0,
u(x, 0) = \varphi(x),
\]

У результаті ми отримаємо

\[
\sum_{n=1}^{\infty} T'_n(t) \sin(n\pi x) = -\alpha^2 \sum_{n=1}^{\infty} (n\pi)^2 T_n(t) \sin(n\pi x) + + \sum_{n=1}^{\infty} f_n(t) \sin(\pi n\lambda),
\]

(9.7)

\[
\sum_{n=1}^{\infty} T_n(t) \sin 0 = 0
\]

(задовольняється ідентично),
\(\sum_{n=1}^{\infty} T_n(t) \sin (n\pi) = 0\) (задовольняється
ідентично), \(\sum_{n=1}^{\infty} T_n(0) \sin (n\pi x) = \phi(x)\) .

Переписування рівняння та початкової умови

(PPP)

\[
\sum_{n=1}^{\infty} [T'_n + (n\pi\alpha)^2 T_n - f_n(t)] \sin(n\pi x) = 0,
\]

(НУ)

\[
\sum_{n=1}^{\infty} T_n(0) \sin(n\pi x) = \varphi(x),
\]

можна побачити, що функції \(T_n(\ell)\) є розв'язками задачі Коші

(9.8)

\[
T_{n} + (n\pi\alpha)^{2} T_{n} = f_{n}(t),
T_{n}(0) = 2 \int_{0}^{1} \varphi(\xi) \sin(n\pi\xi) d\xi = a_{n}.
\]

Усі ці TAC легко розв'язати, а їхні рішення записуються у вигляді

\begin{enumerate}
\def\labelenumi{(\arabic{enumi})}
\setcounter{enumi}{98}
\tightlist
\item
\end{enumerate}

\[
T_n(t) = a_n e^{-(n\pi\alpha)^2 t} + \int_0^t e^{-(n\pi\alpha)^2 (t-\tau)} f_n(\tau) d\tau,
\]

З того місця, де легко отримати повне розв'язання проблеми (9.3):

(9.10)

\[
u(x, t) = \sum_{n=1}^{\infty} T_n(t) \sin(n\pi x) =
= \sum_{n=1}^{\infty} \left[ a_n e^{-(n\pi\alpha)^n t} \sin(n\pi x) \right] + \sum_{n=1}^{\infty} \left[ \sin(n\pi x) \int_0^t e^{-(n\pi\alpha)^n (t-\tau)} f_n(\tau) d\tau \right]
\]

Перехідна частина (обумовлена стаціонарною частиною (умовлена початковою
умовою) правою частиною \(f(x, t)\)).

Розв'язок (9.10) показує, що температурна реакція складається з двох
частин: перша частина зумовлена початковими умовами, а друга ---
густиною джерела \(f(x,t)\). Зверніть увагу, що термін «стаціонарна
частина» не є найкращим описом другої частини температурної реакції,
оскільки він не означає, що ця частина не залежить від часу
(стаціонарний стан може бути періодичним, якщо \(f(x, t)\) є періодичною
функцією часу t).

Проблема повністю вирішена. Перш ніж перейти до наступної теми, давайте
покажемо, як використовувати описаний вище метод для розв'язання

Конкретне завдання.

70

\section{Розв'язання конкретної задачі методом декомпозиції за власними
функціями}\label{ux440ux43eux437ux432ux44fux437ux430ux43dux43dux44f-ux43aux43eux43dux43aux440ux435ux442ux43dux43eux457-ux437ux430ux434ux430ux447ux456-ux43cux435ux442ux43eux434ux43eux43c-ux434ux435ux43aux43eux43cux43fux43eux437ux438ux446ux456ux457-ux437ux430-ux432ux43bux430ux441ux43dux438ux43cux438-ux444ux443ux43dux43aux446ux456ux44fux43cux438}

Розглянемо просту задачу

(9.11)

\[
\begin{aligned} u_t &= \alpha^2 u_{xx} + \sin{(3\pi x)}, & 0 < x < 1, \\ (\mathbf{FY}) & \begin{cases} u(0, t) &= 0, \\ u(1, t) &= 0, \end{cases} & 0 < t < \infty, \\ (\mathbf{HY}) & u(x, 0) &= \sin{(\pi x)}, & 0 \leq x \leq 1. \end{aligned}
\]

Наша мета --- обчислити коефіцієнти \(T_n(t)\) при розкладі

\[
u\left(x,\,t\right)=\sum_{n=1}^{\infty}\,T_{n}\left(t\right)\sin\left(n\pi x\right)
\]

(Власні функції цієї задачі \(X_n(x)\) збігатися з власними функціями
попереднього). Якщо підставити цю декомпозицію на початкову задачу,
отримаємо наступну задачу Коші для визначення функцій \(T_n(t)\) :

\[
T_n + (n\pi\alpha)^2 T_n = \int_n (t) = 2 \int_0^1 \sin(3\pi x) \sin(n\pi x) dx =
\begin{cases} 1 &n = 3, \\ 0 & n \neq 3, \end{cases}
T_n(0) = 2 \int_0^1 \sin(n\xi) \sin(n\pi\xi) d\xi =
\begin{cases} 1 &n = 1, \\ 0 & n \neq 1. \end{cases}
\]

Перепишем эти уравнения последовательно для n=1, 2\ldots{}

\begin{verbatim}
\begin{aligned}
(n=1)\quad & T_1' + (\pi\alpha)^2 T_1 = 0, && T_1(0) = 1, && T_1(t) = e^{-(\pi\alpha)^2 t}, \\
(n=2)\quad & T_2' + (2\pi\alpha)^2 T_2 = 0, && T_2(0) = 0, && T_2(t) = 0, \\
(n=3)\quad & T_3' + (3\pi\alpha)^2 T_3 = 4, && T_3(0) = 0, && T_3(t) = \frac{1}{(3\pi\alpha)^2}\left[1 - e^{-(3\pi\alpha)^2 t}\right], \\
(n\geqslant 4)\quad & T_n' + (n\pi\alpha)^2 T_n = 0, && T_n(0) = 0, && T_n(t) = 0.
\end{aligned}
\end{verbatim}

Отже, розв'язок задачі можна записати у вигляді

\[
u(x, t) = e^{-(\pi\alpha)^2 t} \sin(\pi x) + \frac{1}{(3\pi\alpha)^2} [1 - e^{-(3\pi\alpha)^2 t}] \sin(3\pi x).
\]

(9.12)

\section{ЗАУВАЖЕННЯ}\label{ux437ux430ux443ux432ux430ux436ux435ux43dux43dux44f-6}

\begin{enumerate}
\def\labelenumi{\arabic{enumi}.}
\item
  Метод розкладу власних функцій є одним із найпотужніших методів
  розв'язання неоднорідних диференціальних рівнянь у похідних похідних.
  Пізніше, коли ми ознайомимося з інтегральними перетвореннями,
  побачимо, що існують й інші методи розв'язання задач такого типу.
\item
  Власні функції \(X_n\)(x) , за допомогою якої виконується розклад,
  залежать від задачі, яку потрібно розв'язати, точніше від рівнянь і
  граничних умов. Читача просять розв'язати задачу 4 наприкінці лекції,
  щоб перевірити, чи дійсно він навчився знаходити власні функції
  \(X_n\)(x) .
\item
  Якщо читач пам'ятає основні факти з теорії звичайних диференціальних
  рівнянь, то він знає, що розв'язки, що відповідають неоднорідним
  членам вигляду
\end{enumerate}

\[
P_n(x) e^{\alpha x}
\begin{cases} \sin{(\beta x)}, \\ \cos{(\beta x)}, \end{cases}
\]

можно найти методом неопределенных коэффициентов. То же самое и здесь.
Задача (9.11) может быть решена таким же способом. \# ЗАДАЧІ \# 1.
Решение задачи

\[
\begin{array}{lll} (\text{УЧП}) & u_t = u_{xx} + \sin{(3\pi x)}, & 0 < x < 1, \\ (\Gamma \text{У}) & \begin{cases} u \ (0, \ t) = 0, \\ u \ (1, \ t) = 0, \end{cases} & 0 < t < \infty \\ (\text{HV}) & u \ (x, \ 0) = \sin{(\pi x)}, \ 0 \leqslant x \leqslant 1 \end{array}
\]

дается формулой (9.12). Соответствует ли это решение вашим интуитивным
представлениям? Что представляет собой полученное решение?

\begin{enumerate}
\def\labelenumi{\arabic{enumi}.}
\setcounter{enumi}{1}
\tightlist
\item
  Решите задачу
\end{enumerate}

\[
\begin{array}{llll} (\text{YUII}) & u_t = u_{xx} + \sin{(\pi x)} + \sin{(2\pi x)}, & 0 < x < 1, & 0 <t < \infty, \\ (\text{\GammaY}) & \begin{cases} u\left(0,\ t\right) = 0, \\ u\left(1,\ t\right) = 0, \end{cases} & 0 < t < \infty, \\ (\text{HY}) & u\left(x,\ 0\right) = 0, & 0 \leqslant x \leqslant 1. \end{array}
\]

\# 3. Розв'яжіть задачу

(PPP)

\[
u_t = u_{xx} + \sin(\pi x), \quad 0 < x < 1, \quad 0 < t < \infty,
\]

(GI)

\[
\begin{cases} u(0, t) = 0, \\ u(1, t) = 0, \\ u(x, 0) = 1, \quad 0 < x < 1, \end{cases}
\]

методом разложения по собственным функциям.

\begin{enumerate}
\def\labelenumi{\arabic{enumi}.}
\setcounter{enumi}{3}
\tightlist
\item
  Найдите решение задачи
\end{enumerate}

\[
\begin{aligned}
(\text{УЧП})\quad & u_t = u_{xx} + \sin(\lambda_1 x), && 0 < x < 1,\ 0 < t < \infty, \\
(\text{ГУ})\quad & u(0,t) = 0,\quad u_x(1,t) + u(1,t) = 0, && 0 < t < \infty, \\
(\text{НУ})\quad & u(x,0) = 0, && 0 \leqslant x \leqslant 1.
\end{aligned}
\]

методом разложения по собственным функциям, если \(\lambda_1\) ---
первый корень уравнения \(\operatorname{tg}\lambda = -\lambda\). Каковы
собственные функции этой задачи?

\section{5. Решите
задачу}\label{ux440ux435ux448ux438ux442ux435-ux437ux430ux434ux430ux447ux443}

\[
\begin{aligned}
(\text{УЧП})\quad & u_t = u_{xx}, && 0 < x < 1,\ 0 < t < \infty, \\
(\text{ГУ})\quad & u(0,t) = 0,\quad u(1,t) = \cos t, && 0 < t < \infty, \\
(\text{НУ})\quad & u(x,0) = \varphi(x), && 0 \leqslant x \leqslant 1.
\end{aligned}
\]

Спочатку перетворивши граничні умови в нуль. Отримана задача
розв'язується методом розкладу за власними функціями.

\section{ІНТЕГРАЛЬНІ ПЕРЕТВОРЕННЯ (СИНУСОЇДАЛЬНІ ТА КОСИНУСОВІ
ПЕРЕТВОРЕННЯ)}\label{ux456ux43dux442ux435ux433ux440ux430ux43bux44cux43dux456-ux43fux435ux440ux435ux442ux432ux43eux440ux435ux43dux43dux44f-ux441ux438ux43dux443ux441ux43eux457ux434ux430ux43bux44cux43dux456-ux442ux430-ux43aux43eux441ux438ux43dux443ux441ux43eux432ux456-ux43fux435ux440ux435ux442ux432ux43eux440ux435ux43dux43dux44f}

МЕТА ЛЕКЦІЇ: Ввести концепцію інтегрального перетворення та показати, як
використання інтегральних перетворень дозволяє перейти від CPP з n
змінними до нового рівняння з n---1 змінними.

Лекція показує, що інтегральні перетворення можна інтерпретувати як
розклад вхідної дії на елементарні (частотний розклад), визначення
реакції системи на елементарні впливи та подальше підсумовування всіх
отриманих результатів.

У інтегральних перетвореннях операція диференціювання переходить у
операцію множення, і, відповідно, деякі часткові похідні переходять у
алгебраїчні вирази.

Для розв'язання задач дифузії на необмежених інтервалах вводять
синусоїдальні та косинусні перетворення. Розв'язок, отриманий у цьому
випадку, цікавий тим, що містить функцію, яка є додатковою до
ймовірнісного інтеграла.

Інтегральне перетворення зазвичай називається перетворенням, яке кожній
функції \(f(t)\) призначає нову функцію \(F(s)\) відповідно до формули

\[
F(s) = \int_{A}^{B} K(s, t) f(t) dt.
\]

Зверніть увагу, що початкова функція залежить від t, а отримана функція
--- від s. Функція \(K(s,t)\) називається ядром перетворення. Це
головний елемент, який відрізняє одну трансформацію від іншої. Зазвичай
ядро обирається так, щоб перетворення мало певні задані властивості.
Межі інтегрування також залежать від типу трансформації.

Голосною, що лежить в основі інтегральних перетворень, є здатність
позбутися часткових похідних однієї з незалежних змінних. Це означає, що
перетворене рівняння матиме на одну змінну менше, ніж початкове.
Наприклад, якщо застосувати інтегральне перетворення до рівняння

Щоб позбутися похідної часу, отримуємо звичайне диференціальне рівняння
з незалежною змінною X.

\[
u_{xx} + u_{cy} + u_{zz} = 0
\]

застосувати перетворення Фур'є до змінної x, тоді похідна \(u_{xx}\)
буде усунена, і в новому рівнянні залишаться лише дві незалежні змінні y
та z. Тоді отримуємо звичайне диференціальне рівняння з однією
незалежною змінною. Іншими словами, інтегральні перетворення спрощують
початкове завдання. Потім розв'язується перетворення задачі, і,
застосовуючи обернене перетворення до її розв'язку, отримується
розв'язок початкового перетворення. Схема описаного алгоритму показана
на рис. 10.1

\pandocbounded{\includegraphics[keepaspectratio]{images/_page_72_Figure_5.jpeg}}

Рис. 10.1, Загальна схема застосування інтегрованих перетворень.

Фіг. На рисунку 10.1 видно, що кожне перетворення прямого інтеграла
пов'язане з оберненим перетворенням, яке відновлює початкову функцію з
трансформованої.

Прямі та обернені перетворення зазвичай називають парою перетворень.
Таблиця 1. Рисунок 10.1 показує кілька пар перетворень, які ми будемо
використовувати при розв'язуванні рівнянь із факторами.

Слід зазначити, що в цих перетвореннях використовуються два типи
символів. Наприклад, у випадку перетворення Лапласа позначення
\(\mathcal{L}[f]\) вказує, що функція f перетворена, тоді як позначення
\(F(s)\) означає, що перетворення породжує функцію, яка залежить від s.

У цій лекції ми не прагнемо вивчати всі пари перетворень. Ми
зосередимося лише на перетвореннях Фур'є синуса та косинуса (пари 1 і
2). Трохи пізніше ми вивчаємо іншу з цих пар. Ознайомляючись із
трансформаціями, ми дізнаємося про їхні переваги та недоліки,
дізнаємося, як

Таблиця 10.1. Деякі пари інтегральних перетворень

\begin{enumerate}
\def\labelenumi{\arabic{enumi}.}
\tightlist
\item
  Синус-перетворення Фур'є:
\end{enumerate}

\[
\mathcal{F}_s[f](\omega) = \frac{2}{\pi}\int_0^\infty f(t)\sin(\omega t)\,dt,
\mathcal{F}_s^{-1}[F](t) = \int_0^\infty F(\omega)\sin(\omega t)\,d\omega.
\]

\begin{enumerate}
\def\labelenumi{\arabic{enumi}.}
\setcounter{enumi}{1}
\tightlist
\item
  Косинус-перетворення Фур'є:
\end{enumerate}

\[
\mathcal{F}_c[f](\omega) = \frac{2}{\pi}\int_0^\infty f(t)\cos(\omega t)\,dt,
\mathcal{F}_c^{-1}[F](t) = \int_0^\infty F(\omega)\cos(\omega t)\,d\omega.
\]

\begin{enumerate}
\def\labelenumi{\arabic{enumi}.}
\setcounter{enumi}{2}
\tightlist
\item
  Перетворення Фур'є:
\end{enumerate}

\[
\mathcal{F}[f](\omega) = \frac{1}{\sqrt{2\pi}}\int_{-\infty}^{\infty} f(x)e^{-i\omega x}\,dx,
\mathcal{F}^{-1}[F](x) = \frac{1}{\sqrt{2\pi}}\int_{-\infty}^{\infty} F(\omega)e^{i\omega x}\,d\omega.
\]

\begin{enumerate}
\def\labelenumi{\arabic{enumi}.}
\setcounter{enumi}{3}
\tightlist
\item
  Скінченне синус-перетворення:
\end{enumerate}

\[
S_n = \frac{2}{L}\int_0^L f(x)\sin\!\left(\frac{n\pi x}{L}\right) dx,
f(x) = \sum_{n=1}^{\infty} S_n \sin\!\left(\frac{n\pi x}{L}\right).
\]

\begin{enumerate}
\def\labelenumi{\arabic{enumi}.}
\setcounter{enumi}{4}
\tightlist
\item
  Скінченне косинус-перетворення:
\end{enumerate}

\[
C_n = \frac{2}{L}\int_0^L f(x)\cos\!\left(\frac{n\pi x}{L}\right) dx,
f(x) = \frac{C_0}{2} + \sum_{n=1}^{\infty} C_n \cos\!\left(\frac{n\pi x}{L}\right).
\]

\begin{enumerate}
\def\labelenumi{\arabic{enumi}.}
\setcounter{enumi}{5}
\tightlist
\item
  Перетворення Лапласа:
\end{enumerate}

\[
\mathcal{L}[f](s) = \int_0^\infty e^{-st}f(t)\,dt,
\mathcal{L}^{-1}[F](t) = \frac{1}{2\pi i}\int_{\gamma-i\infty}^{\gamma+i\infty} e^{st}F(s)\,ds.
\]

\begin{enumerate}
\def\labelenumi{\arabic{enumi}.}
\setcounter{enumi}{6}
\tightlist
\item
  Перетворення Ханкеля:
\end{enumerate}

\[
H\{f\}(\xi) = F_n(\xi) = \int_0^\infty r J_n(\xi r) f(r)\,dr,
H^{-1}[F_n](r) = \int_0^\infty \xi J_n(\xi r) F_n(\xi)\,d\xi.
\]

Однак перед тим, як перейти до вивчення інтегральних перетворень, буде
корисно зрозуміти, що зазвичай називають спектром функції (або
спектральним розкладом функції).

\section{Спектральна
функція}\label{ux441ux43fux435ux43aux442ux440ux430ux43bux44cux43dux430-ux444ux443ux43dux43aux446ux456ux44f}

Інтегральні перетворення та спектр функції тісно пов'язані. Насправді
інтегральне перетворення можна інтерпретувати як розклад функції на
певний спектр компонент. Цей спектр залежить від обраного перетворення,
але кожна функція має певний спектр відносно цього перетворення Р и с

\pandocbounded{\includegraphics[keepaspectratio]{images/_page_74_Picture_7.jpeg}}

. 10 2 Розклад періодичної функції на ряд синусами та косинусами;
Послідовність квадратних імпульсів апроксимується синусами та косинусами

Наприклад, розглянемо розклад періодичної функції на ряд синусами та
косинусами (ряд Фур'є) 1:

\[
f(x) = \sum_{n=0}^{\infty} \left[ A_n \cos(nx) + B_n \sin(nx) \right]
\]

(див. рис. 10.2).

Тут коефіцієнти \(A_n\) і \(B_n\) позначають внесок \(\cos(nx)\) і
\(\sin(nx)\) у функцію \(f(x)\), тоді як величина

\[
V\overline{A_n^2+B_n^3}
\]

1) Фур'є-ряди детальніше розглядаються у лекції 11,

(називається спектром функції) показує частку амплітуди функції
\(f(x)\), що припадає на частоту n.

Наприклад, якщо функція \(f(x)\) --- це просто сума

синуси та косинуси

\[
f(x) = 1 + \sin x + \frac{1}{5}\sin(3x) + \cos x + \frac{1}{2}\cos(2x) + \frac{1}{4}\cos(4x),
\]

тоді його спектр (дискретний) виглядатиме так, як показано на рис. 10.3.

Зчитуючи значення значення \(\sqrt{A_n^2 + B_n^2}\) з графіка, можна
визначити амплітудне значення компоненти функції \(f(x)\), яка має
частоту n.

\pandocbounded{\includegraphics[keepaspectratio]{images/_page_75_Figure_8.jpeg}}

Рис. 10 3. Дискретний спектр функції \(f(r)\).

Періодичні функції розкладаються на нескінченні ряди (вони мають
дискретний спектр), тоді як неперіодичні функції повинні мати
неперервний спектр. Якщо функція визначена лише на скінченному
інтервалі, її можна періодично розширювати на всю вісь і розкладати на
ряд Фур'є, який представлятиме нашу функцію в межах визначення області.

І хоча неперіодичну функцію не можна розкласти на нескінченний ряд
синусами та косинусами, можна спробувати записати неперервний аналог
ряду Фур'є:

\[
f(x) = \int_{-\infty}^{+\infty} [C(\omega)\cos(\omega x) + S(\omega)\sin(\omega x)] d\omega,
\]

де функції \(S(\omega)\) і \(C(\omega)\) слугують мірою спектральної
густини синуса і косинуса у функції \(f(x)\), а також значення

\[
V \overline{S^2(\omega) + C^2(\omega)}
\]

Він слугує мірою загального внеску частотного \(\omega\) у функцію
\(f(x)\) і називається спектром (неперервним) функції \(f(x)\). Виходячи
з наведеного вище інтуїтивного визначення спектра функції, давайте
проаналізуємо інтегральні перетворення «гвинт за гвинтом». Спершу
запишемо кілька властивостей цих перетворень, які ми будемо широко
використовувати.

\section{Синусоїдальні та косинусні перетворення
похідних}\label{ux441ux438ux43dux443ux441ux43eux457ux434ux430ux43bux44cux43dux456-ux442ux430-ux43aux43eux441ux438ux43dux443ux441ux43dux456-ux43fux435ux440ux435ux442ux432ux43eux440ux435ux43dux43dux44f-ux43fux43eux445ux456ux434ux43dux438ux445}

\begin{enumerate}
\def\labelenumi{\arabic{enumi}.}
\tightlist
\item
\end{enumerate}

\[
\mathcal{F}_s[f'] = -\omega \mathcal{F}_c[f]
\]

(доведено інтегруванням по частинах).

\begin{enumerate}
\def\labelenumi{\arabic{enumi}.}
\tightlist
\item
\end{enumerate}

\[
\mathcal{F}_s[f'] = -\omega \mathcal{F}_c[f]
\]

\begin{enumerate}
\def\labelenumi{\arabic{enumi}.}
\setcounter{enumi}{1}
\tightlist
\item
  \(\mathcal{F}_s[f''] = \frac{2}{\pi} \omega f(0) - \omega^2 \mathcal{F}_s[f]\)
  (10.1) 3.
\end{enumerate}

\[
\mathcal{F}_{\rho}[f'] = \frac{2}{\pi} f(0) + \omega \mathcal{F}_{s}[f]
\]

.

\begin{enumerate}
\def\labelenumi{\arabic{enumi}.}
\setcounter{enumi}{3}
\tightlist
\item
\end{enumerate}

\[
\mathcal{F}_{c}[f'') = -\frac{2}{\pi} f_{c}(0) - \omega^{2} \mathcal{F}_{c}[f].
\]

Деякі інші властивості прямих і обернених синусоїдальних і спеційних
перетворень наведені в таблиці наприкінці книги. Тепер покажемо, як
синусоїдальне перетворення можна використати для розв'язання важливої
змішаної задачі.

\section{Розв'язок задачі дифузії на напівпрямому синусоїдальному
метрі}\label{ux440ux43eux437ux432ux44fux437ux43eux43a-ux437ux430ux434ux430ux447ux456-ux434ux438ux444ux443ux437ux456ux457-ux43dux430-ux43dux430ux43fux456ux432ux43fux440ux44fux43cux43eux43cux443-ux441ux438ux43dux443ux441ux43eux457ux434ux430ux43bux44cux43dux43eux43cux443-ux43cux435ux442ux440ux456}

Задача, яку ми хочемо розв'язати, формулюється так:

(PPP)

\[
u_t = \alpha^2 u_{xx}
\]

, \(0 < x < \infty\) , \(0 < t < \infty\) , (GU) \(u(0, t) = A\) ,
\(0 < t < \infty\) , (NU) \(u(x, 0) = 0\) , \(0 \le x < \infty\) .

\pandocbounded{\includegraphics[keepaspectratio]{images/_page_76_Figure_14.jpeg}}

Рис. 10.4. Задача дифузії в напівнескінченній області;
\(u\left(0,\,t\right)=\) = A, \(u(x, 0)\) = 0.

КРОК 1. Розберемо розв'язання проблеми на три досить прості кроки.
Діятимемо відповідно до наступного плану. На першому кроці виконуємо
перетворення синуса Фур'є на змінній x.

тоді отримуємо звичайне диференціальне рівняння для змінної t. Для
реалізації цього плану застосовуємо синусоїдальне перетворення до обох
частин рівняння з частковими похідними. У результаті ми отримуємо

\[
\mathcal{F}_s[u_t] = \alpha^2 \mathcal{F}_s[u_{xx}].
\]

Розглянемо кожну частину цього співвідношення. \(\mathcal{F}_s[u_t]\) :
часткову похідну \(u_t\) в цій задачі можна назвати зовнішньою похідною,
оскільки ми використовуємо перетворення на x. У цьому випадку можна
записати

\[
\mathcal{F}_{s}\left[u_{t}\right] = \frac{2}{\pi} \int_{0}^{\infty} u_{t}\left(x, t\right) \sin\left(\omega x\right) dx =
= \frac{\partial}{\partial t} \left[\frac{2}{\pi} \int_{0}^{\infty} u\left(x, t\right) \sin\left(\omega x\right) dx\right] =
= \frac{d}{dt} \mathcal{F}_{s}\left[u\right] =
= \frac{d}{dt} U\left(t\right).
\]

Тут ми взяли операцію диференціації поза знаком інтеграла, застосовуючи
відповідну теорему інтегрального числення. Зверніть увагу, що початкова
функція u залежить від змінних x і t, а також від її образу

\[
\mathcal{F}_{s}[u] = U(\omega, t)
\]

\begin{verbatim}
залежить від змінних \omega і t. Нова змінна \omega відіграє роль параметра в новій задачі, і тому можна припустити, що образи функцій залежать лише від однієї змінної t:
\end{verbatim}

\[
\mathcal{F}_s[u] = U(t)
\]

\begin{verbatim}
Для члена \mathcal{F}_s[u_{xx}] маємо
\end{verbatim}

\[
\mathcal{F}_s[u_{xx}] = \frac{2}{\pi}\omega u(0,t) - \omega^2 \mathcal{F}_s[u]
= \frac{2A\omega}{\pi} - \omega^2 U(t).
\]

\begin{verbatim}
Зверніть увагу, що при виведенні відношень (10.1) ми розглядали функцію f як функцію однієї змінної x. Тут випадок дещо інший, оскільки u\left(x,\ t\right) залежить і від x, і від t. Однак ті ж формули можна використовувати, якщо перетворення виконується однією змінною, а решта вважаються константами.
\end{verbatim}

\begin{verbatim}
У нашому випадку перетворення виконується змінною x, тобто змінну t можна вважати сталою. Звернімо також увагу на той факт, що вже на цьому етапі розв'язання задачі використовується гранична умова u\left(0,\ t\right)=A . Підставляючи знайдені вирази у рівняння u_{t}=\alpha^{2}u_{xx} , отримуємо звичайне диференціальне рівняння
\end{verbatim}

\[
\frac{dU}{dt} = \alpha^2 \left[ -\omega^2 U(t) + \frac{2A\omega}{\pi} \right].
\]

Єдине, що існує --- це початкова умова функції \(U(t)\). Отримуємо його,
застосовуючи перетворення до початкової умови \(u(x, 0)\) = 0:

\[
\mathcal{F}_{s}[u(x, 0)] = U(0) = 0.
\]

Це завершує перший крок у побудові розв'язку методом інтегрального
перетворення: замість початкової змішаної задачі ми отримали задачу Коші
для звичайного диференціального рівняння:

(10.2)

\[
\begin{aligned} (O \Pi V) & \frac{dU}{di} + \omega^2 \alpha^2 U = \frac{2A\omega \alpha^2}{\pi}, \\ (H V) & U(0) = 0. \end{aligned}
\]

КРОК 2. Щоб розв'язати отриману задачу Коші, можна використати будь-який
елементарний метод розв'язання звичайних диференціальних рівнянь. У
результаті ми отримуємо

\[
U(t) = \frac{2A}{\pi\omega} (1 - e^{-\omega^2 \alpha^2 t}).
\]

\begin{verbatim}
Отже, перед нами маємо результат застосування синусоїдального перетворення до бажаної функції u\left(x,\ t\right) . Останній крок — застосувати обернене перетворення до функції U\left(t\right) і отримати потрібне рішення u\left(x,\ t\right) :
\end{verbatim}

\[
u(x, t) = \mathcal{F}_{x}^{-1}[U].
\]

КРОК 3. Щоб отримати розв'язок, можна або безпосередньо обчислити
інтеграл, що входить до оберненого перетворення, або використати таблиці
перетворень. Використовуючи таблиці, ми знаходимо

\[
u(x, t) = A \operatorname{erfc}(x/2\alpha \sqrt{t}),
\]

де erfc(x) --- функція, доповнююча до ймовірнісного інтегралу, а
визначається відношенням

\[
\operatorname{erfc}(x) = \frac{2}{\sqrt{2\pi}} \int_{x}^{\infty} e^{-t^{2}} dt.
\]

Графік цієї функції показано на рис. 10.5.

Ймовірнісний інтеграл і додаткова функція до нього не виражаються через
елементарні функції. Вони добре популярні, і відповідні таблиці можна
знайти в довідниках з фізики та хімії.

\pandocbounded{\includegraphics[keepaspectratio]{images/_page_79_Figure_4.jpeg}}

Фіг. 10.5. Графи функцій erf (x) та eric (x).

\section{Інтерпретація
рішення}\label{ux456ux43dux442ux435ux440ux43fux440ux435ux442ux430ux446ux456ux44f-ux440ux456ux448ux435ux43dux43dux44f}

Розв'язок

\[
u(x, t) = A \operatorname{erfc}[x/2\alpha V \overline{t}]
\]

дає багато їжі для роздумів. На рисунку 10.6 показані графіки
ймовірнісної комплементарної функції на різних значеннях часу t.

\pandocbounded{\includegraphics[keepaspectratio]{images/_page_79_Figure_10.jpeg}}

Рис. 10.6. Розподіл температури в напівнескінченному стрижні, на кінці
якого підтримується постійна температура.

Вони відрізняються один від одного лише характеристичною шкалою вздовж
осі x.

\section{ЗАВДАННЯ}\label{ux437ux430ux432ux434ux430ux43dux43dux44f-8}

\begin{itemize}
\tightlist
\item
  \begin{enumerate}
  \def\labelenumi{\arabic{enumi}.}
  \tightlist
  \item
    Довести ідентичності (10.1).
  \end{enumerate}
\item
  \begin{enumerate}
  \def\labelenumi{\arabic{enumi}.}
  \setcounter{enumi}{1}
  \tightlist
  \item
    Розв'яжіть задачу Коші (10.2).
  \end{enumerate}
\item
  \begin{enumerate}
  \def\labelenumi{\arabic{enumi}.}
  \setcounter{enumi}{2}
  \tightlist
  \item
    Використайте синусоїдальне або косинусне перетворення для
    розв'язання задачі
  \end{enumerate}
\end{itemize}

\begin{verbatim}
\begin{aligned}
(\text{УЧП})\quad & u_t = \alpha^2 u_{xx}, && 0 < x < \infty,\ 0 < t < \infty, \\
(\text{ГУ})\quad & u_x(0,t) = 0, && 0 < t < \infty, \\
(\text{НУ})\quad & u(x,0) = H(1-x), && 0 \leqslant x < \infty.
\end{aligned}
\end{verbatim}

де \(H(x)\) --- функція Хевісайду

\[
H(x) =
\begin{cases} 0, & x < 0, \\ 1, & 0 \le x. \end{cases}
\]

Входящая в начальное условие функция изображена на рисунке.

\pandocbounded{\includegraphics[keepaspectratio]{images/_page_80_Picture_7.jpeg}}

Как будет выглядеть график решения в различные моменты времени?

\begin{enumerate}
\def\labelenumi{\arabic{enumi}.}
\setcounter{enumi}{4}
\tightlist
\item
  Найдите прообраз константы при синус-преобразовании, т. е. найдите
  \(f(x)\), если \(F(\omega) = c = \mathrm{const.}\)
\end{enumerate}

\section{\texorpdfstring{РЯДЫ И ПРЕОБРАЗОВАНИЯ ФУРЬЕ МЕТА ЛЕКЦІЇ: Ввести
понятие ряда Фурье и показать, каким образом периодическую функцию
\(f(x)\) можно представить в гиде ряда по синусам и
косинусам:}{РЯДЫ И ПРЕОБРАЗОВАНИЯ ФУРЬЕ МЕТА ЛЕКЦІЇ: Ввести понятие ряда Фурье и показать, каким образом периодическую функцию f(x) можно представить в гиде ряда по синусам и косинусам:}}\label{ux440ux44fux434ux44b-ux438-ux43fux440ux435ux43eux431ux440ux430ux437ux43eux432ux430ux43dux438ux44f-ux444ux443ux440ux44cux435-ux43cux435ux442ux430-ux43bux435ux43aux446ux456ux457-ux432ux432ux435ux441ux442ux438-ux43fux43eux43dux44fux442ux438ux435-ux440ux44fux434ux430-ux444ux443ux440ux44cux435-ux438-ux43fux43eux43aux430ux437ux430ux442ux44c-ux43aux430ux43aux438ux43c-ux43eux431ux440ux430ux437ux43eux43c-ux43fux435ux440ux438ux43eux434ux438ux447ux435ux441ux43aux443ux44e-ux444ux443ux43dux43aux446ux438ux44e-fx-ux43cux43eux436ux43dux43e-ux43fux440ux435ux434ux441ux442ux430ux432ux438ux442ux44c-ux432-ux433ux438ux434ux435-ux440ux44fux434ux430-ux43fux43e-ux441ux438ux43dux443ux441ux430ux43c-ux438-ux43aux43eux441ux438ux43dux443ux441ux430ux43c}

\[
f(x) = \frac{a_0}{2} + \sum_{n=1}^{\infty} [a_n \cos(n\pi x) + b_n \sin(n\pi x)].
\]

Если же функция \(f(x)\) непериодическая и определена на
\((-\infty, +\infty)\) , то показать, как ряд Фурье переходит в
\emph{преобразование Фурье} и как функцию \(f(x)\) можно разложить в
непрерывную совокупность простейших функций. Эго разложение (интеграл
Фурье) можно записать в комплексной форме:

\[
f(x) = \frac{1}{\sqrt{2\pi}} \int_{-\infty}^{+\infty} \left[ \frac{1}{\sqrt{2\pi}} \int_{-\infty}^{+\infty} f(x) e^{-i\xi x} dx \right] e^{i\xi x} d\xi,
\]

откуда получаются прямое и обратное преобразования Фурье

\[
\mathcal{F}[f] = F(\xi) = \frac{1}{\sqrt{2\pi}} \int_{-\infty}^{+\infty} f(x) e^{-i\xi x} dx \qquad \text{(пряме перетворення Фур'є),}
\mathcal{F}^{-1}[F] = f(x) = \frac{1}{\sqrt{2\pi}} \int_{-x}^{+x} F(x) e^{i\xi x} d\xi \qquad \text{(обернене перетворення Фур'є)}.
\]

\begin{verbatim}
Важлива роль рядів Фур'є в теорії диференціальних рівнянь у похідних пов'язана з тим, що періодичні функції, визначені в (-\infty, +\infty) або функції, визначені на скінченному інтерваті, можна розкласти на ряд за допомогою сіпусів і косинусів. Це дозволяє спростити початкове завдання. Наприклад, функція пилкоподібної форми
\end{verbatim}

\[
f(x) = x, -L < x < L,
\]

\(f(x + 2L)\) = \(f(x)\) (условие периодичности),

зображено на рис. 11.1 можна представити як ряд Фур'є

(11.1)

\[
f(x) = \frac{a_0}{2} + \sum_{n=1}^{\infty} \left[ a_n \cos(n\pi x/L) + b_n \sin(n\pi x/L) \right],
\]

у яких коефіцієнти Фур'є \(a_n\) і \(b_n\) знаходяться відповідно до
формул Ейлера

\[
a_n = \frac{1}{L} \int_{-L}^{L} f(x) \cos(n\pi x/L) dx = 0, \quad n = 0, 1, 2, \dots,
b_n = \frac{1}{L} \int_{-L}^{L} f(x) \sin(n\pi x/L) dx = -(-1)^n (2L/n\pi), \quad n = 1, 2, \dots.
\]

Эти ингегралы легко вычисляются. Чтобы вывести формулы Эйлера для
коэффициентов \(a_n\) и \(b_n\) , умножим обе части разложения (11.1) на
\(\sin{(n\pi x/L)}\) или \(\cos{(n\pi x/L)}\) и проинтегрируем получив-
Р И С

\pandocbounded{\includegraphics[keepaspectratio]{images/_page_82_Figure_7.jpeg}}

. 11.1. Разложение пилообразной функции в ряд Фурье.

шееся соотношение от -L до L. Ортогональность функций
\(\{\sin(n\pi x/L)\}\) и \(\{\cos(n\pi x/L)\}\) позволяет срязу получить
формулы для коэффициентов \(a_n\) и \(b_n\) (см. задачу 6 в конце
лекции). Для пилообразной функции разложение в ряд Фурье принимает вид

(11.3)

\[
f(x) = \frac{2L}{\pi} \left[ \sin(\pi x/L) - \frac{1}{2} \sin(2\pi x/L) + \frac{1}{3} \sin(3\pi x/L) - \dots \right].
\]

Каждое слагаемое ряда принято называть гармоникой, частота каждой
следующей гармоники выше предыдущей. Частоты всех гармоник кратны
основной частоте, период которой совпадает

с периодом функции \(f(x)\). Один из недостатков рядов Фурье состоит в
том, что в ряд Фурье можно разложить только периодическую функцию.
Конечно, если функция определена на конечном интервале (например,
\(f(x)\) = x при \(0 \le x \le 1\) ), ее можно периодически продолжить
на всю числовую ось и затем использовать разложение (11.1).

Тот факт, что сумма ряда Фурье определена и периодична вне интервала
{[}0, 1{]}, ничему не противоречит, поскольку мы интересуемся поведением
функции только внутри интервала. На самом деле для функции, определенной
внутри интервала, можно указать много различных разложений в ряд Фурье,
поскольку функцию можно продолжить за пределы интервала различными
способами. Сразу же отметим, что полученные ряды будут сходиться к
исходной функции с различными скоростями.

Читатель не должен думать, что всякую периодическую функцию можно
разложить в ряд Фурье. Пока мы знаем только, что если функцию \(f(x)\)
можно разложить в ряд Фурье (11.1), то коэффициенты \(a_n\) и \(b_n\)
находятся по формулам Эйлера (11.2). Более того, даже если функция
разложима в ряд Фурье, мы не знаем, можно ли ряд почленно
дифференцировать. И в самом деле, легко видеть, что производную функции
\(f(x)\) = x (пилообразной функции) нельзя найти почленным
дифференцированием ряда Фурье (11.3).

В литературе, перечисленной в конце книги, можно найти точные условия
разложимости функции \(f(x)\) в ряд Фурье, а также условия почленной
дифференцируемости получающегося ряда. Нам будет достаточно ознакомиться
с известной теоремой Дирихле.

\emph{Теорема Дирихле} (достаточные условия разложимости функции в ряд
Фурье):

Если \(f(x)\)---ограниченная периодическая функция, имеющая на каждом
периоде конечное число максимумов, минимумов и точек разрыва, то ряд
Фурье функции \(f(x)\) сходится к \(f(x)\) в каждой точке непрерывности
функции и к среднему значению пределов функции слева и справа в каждой
точке разрыва функции \(f(x)\).

Например, на рис. 11.1 ряд Фурье сходится к функции \(f(x)\) во всех
точках, за исключением точк \(x=\pm L, \pm 3L, \ldots\) (точки разрыва).
В каждой точке разрыва ряд сходится к нулю (среднее арифметическое
значение функции в точках \(\pm L\) и \(\pm L\) ).

Теперь мы педголовлены для знакомства с преобразованием Фурье. Однако
целесообразно будет познакомиться сначала с понятием частотного спектра
периодической функции.

\section{Дискретный частотный спектр периодической
функции.}\label{ux434ux438ux441ux43aux440ux435ux442ux43dux44bux439-ux447ux430ux441ux442ux43eux442ux43dux44bux439-ux441ux43fux435ux43aux442ux440-ux43fux435ux440ux438ux43eux434ux438ux447ux435ux441ux43aux43eux439-ux444ux443ux43dux43aux446ux438ux438.}

Если \(f(x)\) --- периодическая функция, то разложение в ряд Фурье можно
интерпретировать как представление функции \(f(x)\) последовательностью
чисел \(\{c_n\}\) вида

\[
c_n = \sqrt{a_n^2 + b_n^2}, \quad n = 0, 1, 2, \dots,
\]

где коэффициенты \(c_n\) служат мерой вклада соответствующих частотных
компонент в функцию \(f(x)\). Например, пилообразная функция \(f(x)\)
разлагается в ряд Фурье

\[
f(x) = \frac{2L}{\pi} \left[ \sin(\pi x/L) - \frac{1}{2} \sin(2\pi x/L) + \frac{1}{3} \sin(3\pi x/L) - \dots \right]
\]

и, следовательно, частотный спектр \(\{c_n\}\) записывается в виде
\(c_n=2L/n\pi\) для \(n=1,\ 2,\ \dots\) и \(c_0=|a_0|=0\) (см. рис.
11.2).

\pandocbounded{\includegraphics[keepaspectratio]{images/_page_84_Figure_8.jpeg}}

PNC. 11 2. Дискретный частотный спектр пилообразной функции a---
основная частота (период 2L); b---n-ая гармоника (период 2L/n).

Последовательность \(\{c_n\}\) напоминает спектральное разложение белого
света по цветам в спектроскове.

Введем теперь понятие преобразования Фурье.

\section{Преобразование
Фурье}\label{ux43fux440ux435ux43eux431ux440ux430ux437ux43eux432ux430ux43dux438ux435-ux444ux443ux440ux44cux435}

Как уже отмечалось рансе, непериодические функции, заданные на
\((-\infty, +\infty)\) , не разлагаются в ряды Фурье Однако для
некоторых таких функций можно построить аналог ряда Фурье. Не вдаваясь в
детали доказательства, можно утверждать, что для непериодических
функций, определенных на \((-\infty, \infty)\) , ряд Фурье

\[
f(x_{l} = \frac{a_{0}}{2} + \sum_{n=1}^{\infty} \left[ a_{n} \cos(n\pi x/L) + b_{n} \sin(n\pi x/L) \right]
\]

переходит в интеграл Фирье (разложение по непрерывному спектру частот)

(11.4)

\[
f(x) = \int_{0}^{\infty} a(\xi) \cos(\xi x) d\xi + \int_{0}^{\infty} b(\xi) \sin(\xi x) d\xi,
\]

гле

(11.5)

\[
a(\xi) = \frac{1}{\pi} \int_{0}^{\infty} f(x) \cos(\xi x) dx
\]

, \(b(\xi) = \frac{1}{\pi} \int_{0}^{\infty} f(x) \sin(\xi x) dx\) .

Из (11.4) видно, что интеграл Фурье осуществляет разложение функции
\(f(x)\) в спектр по всем частотам \(0 < \xi < \infty\) (а не Р И С

\pandocbounded{\includegraphics[keepaspectratio]{images/_page_85_Figure_7.jpeg}}

11.3. Частотные спектры некоторых функций.

только по частотам, кратным основной частоте, как это было у
периодических функций). По аналогии с рядами Фурье определим частотный
спектр по формуле

\[
C(\xi) = \sqrt{a^2(\xi) + b^2(\xi)}.
\]

Частотний спектр є мірою внеску частоти \(\xi\) у функ-

иню \(f(x)\).

На рис. 11.3 изображены некоторые функции \(f(x)\) и их спектры. Из рис.
11.3 видно, что функции, графики которых имеют угловые точки, обладают
широким частотным спектром. Объясняется это тем, что для передачи излома
графкка требуются высокочастотные спектральные компоненты. С другой
стороны, если \(f(x) = \sin \xi_0 x\) ---простая периодическая функция,
то ее частотный спектр обращается в нуль всюду, за исключением точки
\(\xi = \xi_0\) .

Теперь мы готовы к тому, чтобы познакомиться с экспоненциальной формой
преобразования Фурье (формулы 11.5 известны как синус- и
косинус-преобразования Фурье). Воспользуемся фор-

мулой Эйлера

\[
e^{i\theta} = \cos \theta + i \sin \theta
\]

и носле несложных преобразований приведем (11.4) к виду

(11.6)

\[
f(x) = \frac{1}{\sqrt{2\pi}} \int_{-\infty}^{+\infty} \left[ \frac{1}{\sqrt{2\pi}} \int_{-\infty}^{+\infty} f(x) e^{-i\xi x} dx \right] e^{i\xi x} d\xi,
\]

которое принято называть интегральным представлением Фурье. Отсюда
получаем пару интегральных преобразований

(11.7)

\[
\mathcal{F}[f] = F(\xi) = \frac{1}{\sqrt{2\pi}} \int_{-\infty}^{+\infty} f(x) e^{-i\xi x} dx
\]

(прямое преобразование Фурье),

\[
\mathcal{F}^{-1}[F] = f(x) = \frac{1}{\sqrt{2\pi}} \int_{-\infty}^{+\infty} F(\xi) e^{i\xi x} d\xi
\]

(обратное преобразование Фурье).

В следующей лекции мы остановимся на изучении свойств этих
преобразований. \# ЗАУВАЖЕННЯ 1. Фурье-образ \(F(\xi)\) функции \(f(x)\)
может оказаться комплекснозначной функцией. Например, для функции

\[
f(x) =
\begin{cases} 0, & x \le 0, \\ e^{-x}, & x > 0, \end{cases}
\]

Фур'є-образ має вигляд

\[
F(\xi) = \frac{1}{\sqrt{2\pi}} \frac{1 - i\xi}{1 + \xi^2}.
\]

\begin{verbatim}
2. Абсолютне значення зображення Фур'є |F(\xi)| називається частотним спектром функції f(x). Наприклад, частотний спектр функції f(x) з попереднього зауваження визначається формулою
\end{verbatim}

\textbar F(\xi)\textbar{} = \sqrt{\frac{1}{2$\pi$(1+\xi^2)}}

(зчитувач повинен мати змогу знайти модуль комплексного числа).

\begin{verbatim}
3. Не всі функції піддаються перетворенню Фур'є (інтеграл (11.7) може не існувати). Наприклад, функції f(x) = c, \sin x , e^x , не мають x^2 образу Фур'є. Лише ті функції, які мають тенденцію швидко нульовою на |x| \to \infty 1). У прикладних задачах ми застосовуємо перетворення Фур'є до температурних, хвильових та інших фізичних полів, які мають тенденцію до нуля на |x| \to \infty .
\end{verbatim}

\section{ЗАВДАННЯ}\label{ux437ux430ux432ux434ux430ux43dux43dux44f-9}

\begin{enumerate}
\def\labelenumi{\arabic{enumi}.}
\tightlist
\item
  Розташуйте функцію у рядку Фур'є над синусами
\end{enumerate}

\[
f(x) =
\begin{cases} -1, & -1 < x < 0, \\ 1, & 0 \le x < 1, \end{cases}
f(x+2) = f(x)
\]

(умова періодичності).

Побудуйте суму перших 2, 3 і 4 членів, щоб слідувати збіжності ряду до
\(f(x)\). Побудуйте графіки частотного спектра функції \(f(x)\).

\begin{itemize}
\tightlist
\item
  \begin{enumerate}
  \def\labelenumi{\arabic{enumi}.}
  \setcounter{enumi}{1}
  \tightlist
  \item
    Покажіть, що якщо розрізнити ряд Фур'є (11.3) функції пилки за
    членом, отримаємо послідовність, яка не збігається до похідної
    пилоподібної функції.
  \end{enumerate}
\item
  \begin{enumerate}
  \def\labelenumi{\arabic{enumi}.}
  \setcounter{enumi}{2}
  \tightlist
  \item
    Побудуйте частотний спектр наступних періодичних функцій:
  \end{enumerate}
\end{itemize}

\begin{verbatim}
- a) f(x) = \sin x ,
 - 6) f(x) = \sin x + \cos 2x, c) f(x) = \sin x + \cos x + 0.5 \sin 3x .
\end{verbatim}

\begin{verbatim}
4. Знайти F(\xi) зображення Фур'є та частотний спектр C(\xi) = |F(\xi)| функцію
\end{verbatim}

\[
f(x) =
\begin{cases} 1, & -1 < x < 1, \\ 0 & \text{в остальных точках.} \end{cases}
\]

Тут ми говоримо лише про перетворення Власіка-Фур'є. Ед.

\begin{enumerate}
\def\labelenumi{\arabic{enumi}.}
\setcounter{enumi}{4}
\tightlist
\item
  Показати, що абсолютне значення функції \(F(\xi) = \frac{1}{1+i\xi}\)
  дорівнює
\end{enumerate}

\[
\left|F(\xi)\right| = \sqrt{\frac{1}{1+\xi^2}}
\]

ПРИМІТКА. Спочатку помножте чисельник і знаменник на \(1-i\xi\) , а
потім позбавтеся комплексної одиниці i у знаменнику. \(^1\) )

\begin{enumerate}
\def\labelenumi{\arabic{enumi}.}
\setcounter{enumi}{5}
\tightlist
\item
  Перевірити ортогональність функцій \(\{\sin(n\pi x/L)\}\) та
  \(\{\cos(n\pi x/L)\}\) на сегменті {[}-L, L{]}:
\end{enumerate}

\[
\int_{-L}^{L} \sin(m\pi x/L) \sin(n\pi x/L) dx =
\begin{cases} 0, & m \neq n, \\ L, &m = n, \end{cases}
\int_{-L}^{L} \cos(m\pi x/L) \cos(n\pi x/L) dx =
\begin{cases} 0, & m \neq n, \\ L, &m = n, \end{cases}
\int_{-L}^{L} \sin(m\pi x/L) \cos(n\pi x/L) dx = 0 \text{ npn Bcex } m, n = 1, 2, \dots
\]

\section{Лекція 12}\label{ux43bux435ux43aux446ux456ux44f-12}

\section{ПЕРЕТВОРЕННЯ ФУР'Є ТА ЙОГО ЗАСТОСУВАННЯ ДО РОЗВ'ЯЗАННЯ
ДИФЕРЕНЦІАЛЬНИХ РІВНЯНЬ У ПОХІДНИХ
ПОХІДНИХ}\label{ux43fux435ux440ux435ux442ux432ux43eux440ux435ux43dux43dux44f-ux444ux443ux440ux454-ux442ux430-ux439ux43eux433ux43e-ux437ux430ux441ux442ux43eux441ux443ux432ux430ux43dux43dux44f-ux434ux43e-ux440ux43eux437ux432ux44fux437ux430ux43dux43dux44f-ux434ux438ux444ux435ux440ux435ux43dux446ux456ux430ux43bux44cux43dux438ux445-ux440ux456ux432ux43dux44fux43dux44c-ux443-ux43fux43eux445ux456ux434ux43dux438ux445-ux43fux43eux445ux456ux434ux43dux438ux445}

МЕТА ЛЕКЦІЇ: Проілюструвати важливі властивості перетворення Фур'є та
показати, як ці властивості можна використовувати для розв'язання PPP.
Зокрема, щоб показати, як під впливом перетворення Фур'є операція
диференціації замінюється множенням, а диференціальні рівняння
перетворюються на алгебранхічні рівняння. У цій лекції вводиться поняття
згортки.

2) Легше
\(|F(\xi)| = \frac{|1|}{|1+i\xi|} = \frac{1}{\sqrt{1+\xi^2}} - \Pi_{PVM}, pcd.\)
Якщо функція \(f(x)\) визначена в точці \(-\infty < x < \infty\) , то її
чотиризображення визначається формулою

(12.1)

\[
\mathcal{F}[f] = F(\xi) = \frac{1}{\sqrt{2\pi}} \int_{-\infty}^{+\infty} f(x) e^{-i\xi x} dx.
\]

Отже, якщо у нас є функція \(f(x)\), визначена на дійсній осі, ми
підставляємо її у формулу (12.1) і отримуємо нову функцію, \(F(\xi)\)
визначену для \(-\infty < \xi < \infty\) . Водночас не слід забувати про
конвенції існування раннього перетворення Фур'є (у цьому випадку
класичного).

Таблиця 1. 12.1 показує деякі функції та їхні фур'є-образи. Більш
розгорнута таблиця функцій та їхніх фур'є-образів наведена в додатку. З
наведеної вище таблиці видно, що образ Фур'є функції може виявитися
комплекснозначною функцією аргументу \(\xi\) . У першому прикладі цієї
таблиці функція \(F(\xi)\) містить як множник уявну одиницю i, і таким
чином є комплекснозначною функцією екзистенційного аргументу ( \(\xi\)
змінюється від \(-\infty\) до \(+\infty\) ). Іншими словами, аргумент
\(\xi\) є реальним, а значення функції є складними.

Перетворення Фур'є (як і більшість інших інтегральних перетворень)
виявилося дуже корисним, оскільки воно замінює операцію диференціації на
операцію множення і таким чином перетворює диференціальні рівняння на
алгебранхічні рівняння. Існує багато інших властивостей, які роблять
перетворення Фур'є зручним інструментом для розв'язання задач. Тут ми
перелічимо лише найважливіші з них.

Таблиця 12.1. Перетворення Фур'є деяких функцій

{\def\LTcaptype{none} % do not increment counter
\begin{longtable}[]{@{}
  >{\raggedright\arraybackslash}p{(\linewidth - 2\tabcolsep) * \real{0.4740}}
  >{\raggedright\arraybackslash}p{(\linewidth - 2\tabcolsep) * \real{0.5260}}@{}}
\toprule\noalign{}
\begin{minipage}[b]{\linewidth}\raggedright
Функція / (a)
\end{minipage} & \begin{minipage}[b]{\linewidth}\raggedright
\(\Gamma\) Перетворення Фур'є \(F\) ( \(\xi\) )
\end{minipage} \\
\midrule\noalign{}
\endhead
\bottomrule\noalign{}
\endlastfoot
1. & \\
\end{longtable}
}

\[
f(x) = \begin{cases} e^{-x}, & x \ge 0, \\ -e^{-x}, & x < 0. \end{cases}
\]

\textbar{} \(F(\xi) = -i \sqrt{\frac{2}{\pi}} \frac{\xi}{1 + \xi^2}\)
(Xomnekches funktsia). \textbar{} \textbar{} 2.

\[
f(x) = \begin{cases} 1, & -1 < x < 1, \\ 0, & x \le -1, x \ge 1 \end{cases}
\]

\textbar{} \(F(\xi) = \sqrt{\frac{2}{\pi}}\,\frac{\sin \xi}{\xi}\)
\textbar{} \textbar{} 3. \(f(x) = e^{-x^2}\) \textbar{}
\(F(\xi) = \frac{1}{\sqrt{2}} e^{-\xi^2/4}\) \textbar{}

\section{Властивості Властивості перетворення Фур'є Властивість 1 (Пара
перетворень
Фур'є)}\label{ux432ux43bux430ux441ux442ux438ux432ux43eux441ux442ux456-ux432ux43bux430ux441ux442ux438ux432ux43eux441ux442ux456-ux43fux435ux440ux435ux442ux432ux43eux440ux435ux43dux43dux44f-ux444ux443ux440ux454-ux432ux43bux430ux441ux442ux438ux432ux456ux441ux442ux44c-1-ux43fux430ux440ux430-ux43fux435ux440ux435ux442ux432ux43eux440ux435ux43dux44c-ux444ux443ux440ux454}

Після застосування перетворення Фур'є до функції \(f(x)\), визначеної
для \(-\infty < x < \infty\) , отримуємо нову функцію \(F(\xi)\) за
формулою

\[
\mathscr{F}[f] = F(\xi) = \frac{1}{\sqrt{2\pi}} \int_{-\infty}^{+\infty} f(x) e^{-i\xi x} dx.
\]

Якщо застосувати до функції, \(F(\xi)\) визначеної для
\(-\infty < \xi < \infty\) , обернене перетворення Фур'є

\[
\mathscr{F}^{-1}[F] = f(x) = \frac{1}{\sqrt{2\pi}} \int_{-\infty}^{+\infty} F(\xi) e^{i\xi x} d\xi,
\]

\textbf{Тоді} отримуємо початкову функцію \(f(x)\).

\pandocbounded{\includegraphics[keepaspectratio]{images/_page_90_Figure_7.jpeg}}

Рис. 12.1. Графи функції та її фур'є-образу.

Наприклад, (див. рис. 12.1)

\[
e^{-|x|} \xrightarrow{\mathcal{F}} \sqrt{\frac{2}{\pi}} \xrightarrow{1 + \xi^2} \xrightarrow{\mathcal{F}^{-1}} e^{-|x|}.
\]

\section{Властивість 2 (Перетворення
лінійності)}\label{ux432ux43bux430ux441ux442ux438ux432ux456ux441ux442ux44c-2-ux43fux435ux440ux435ux442ux432ux43eux440ux435ux43dux43dux44f-ux43bux456ux43dux456ux439ux43dux43eux441ux442ux456}

Перетворення Фур'є --- це лінійне перетворення, тобто

\[
\mathscr{F}[af+bg]=a\mathscr{F}[f]+b\mathscr{F}[g].
\]

Це легко побачити. Читач може витратити кілька хвилин, щоб ознайомитися
з цією нерухомістю, яка використовується «на кожному кроці». Наприклад,
перетворення Фур'є виразу

\[
e^{-x} + 3e^{-x^2}
\]

має вигляд

\[
\mathscr{F}[e^{-c}] + 3\mathscr{F}[e^{-c^*}].
\]

\section{Властивість 3 (Часткова конвертація
деривативів)}\label{ux432ux43bux430ux441ux442ux438ux432ux456ux441ux442ux44c-3-ux447ux430ux441ux442ux43aux43eux432ux430-ux43aux43eux43dux432ux435ux440ux442ux430ux446ux456ux44f-ux434ux435ux440ux438ux432ux430ux442ux438ux432ux456ux432}

Якщо нас цікавить, як перетворення Фур'є діє на похідні, то перш за все
слід звернути увагу на змінні, за якими диференційована функція.
Наприклад, якщо перетворення Фур'є виконується на змінній
\(\boldsymbol{x}\) (тобто змінна \(\boldsymbol{x}\) інтегрована у
перетворення), і якщо часткова похідна функції
\(\boldsymbol{u}(\boldsymbol{x},t)\) на тій самій змінній
\(\boldsymbol{x}\) перетворена, то наступні формули перетворення
похідних є дійсними

\[
\begin{split} \mathcal{F}\left\{u_{x}\right\} &= \frac{1}{\sqrt{2\pi}} \int_{-\infty}^{+\infty} u_{x}\left(x,\,t\right) e^{-i\xi x} dx = i\xi \mathcal{F}\left[u\right], \\ \mathcal{F}\left\{u_{xx}\right\} &= \frac{1}{\sqrt{2\pi}} \int_{-\infty}^{+\infty} u_{xx}\left(x,\,t\right) \, \delta^{-i\xi x} dx = -\xi^{2} \mathcal{F}\left[u\right]. \end{split}
\]

С другой стороны, если подвергнуть преобразованию частную производную
\(u_t(x, t)\) (а переменная интегрирования в преобразовании, как и
прежде, будет x), получаем следующие формулы:

\[
\begin{aligned} \mathcal{F}\left[u_{t}\right] &= \frac{1}{\sqrt{2\pi}} \int_{-\infty}^{+\infty} u_{t}\left(x, \ t\right) e^{-i\frac{\pi}{2}x} dx = \frac{\partial}{\partial t} \mathcal{F}\left[u\right], \\ \mathcal{F}\left[u_{tt}\right] &= \frac{1}{\sqrt{2\pi}} \int_{-\infty}^{+\infty} u_{tt}\left(x, \ t\right) e^{-i\frac{\pi}{2}x} dx = \frac{\partial^{2}}{\partial t^{2}} \mathcal{F}\left[u\right]. \end{aligned}
\]

\# Властивість 4 (Властивість згортки)

Кожне інтегральне перетворення має так звану властивість згортки. Факт у
тому, що образ Фур'є добутку двох функцій загалом не відрізняється від
добутку образів Фур'є цих функцій, тобто

\[
\mathcal{F}[f(x)g(x)] \neq \mathcal{F}[f]\mathcal{F}[g].
\]

Однак у теорії перетворень існує операція, яка називається згорткою
функцій f і g (позначена як f*g), яка певною мірою виконує роль добутку,
а саме відношення

(12.2)

\[
\mathcal{F}[/*g] = \mathcal{F}[/]\mathcal{F}[g].
\]

Що це за загадкова операція під назвою згортка? Вона визначається
формулою

(12.3)

\[
(f*g)(x) = \frac{1}{\sqrt{2\pi}} \int_{-\infty}^{+\infty} f(x-\xi) g(\xi) d\xi.
\]

Неважко показати, що співвідношення (12,2) дійсно виконується. З
визначення видно, що згортка двох функцій \(f(x)\) і \(g(x)\) є новою
функцією (f*g)(x) аргументу x.

\section{Приклад обчислення згортки двох
функцій}\label{ux43fux440ux438ux43aux43bux430ux434-ux43eux431ux447ux438ux441ux43bux435ux43dux43dux44f-ux437ux433ux43eux440ux442ux43aux438-ux434ux432ux43eux445-ux444ux443ux43dux43aux446ux456ux439}

Нехай задано дві функції

\[
f(x) = x,
g(x) = e^{-x^2}.
\]

Тоді їхня згортка дорівнює

\[
(f \circ g)(x) = \frac{1}{\sqrt{2\pi}} \int_{-\infty}^{\infty} (x - \xi) e^{-\xi^2} d\xi = x/\sqrt{2}.
\]

Тут ми використали відому формулу

\[
\int_{-\infty}^{+\infty} e^{-\xi^2} d\xi = \sqrt{\pi}.
\]

Важливість концепції згортки (12.3) полягає в тому, що дуже часто,
розв'язуючи диференціальні рівняння, ми стикаємося з необхідністю
застосовувати обернене перетворення Фур'є до виразу, який є добутком
двох образів Фур'є, тобто потрібно знайти

(12.4)

\[
\mathcal{F}^{-1}\left\{\mathcal{F}\left[f\right]\mathcal{F}\left[g\right]\right\}.
\]

Застосовуючи обернене перетворення до обох частин (12.2), отримуємо

\[
f * g = \mathcal{J}^{-1} \{ \mathcal{F} [f] \mathcal{F} [g] \}.
\]

Отже, щоб визначити (12.4), потрібно знайти прототипи кожного з
факторів, тобто функцій f і g, а потім обчислити їхню згортку. Тепер ми
готові розв'язати одну важливу проблему: теорію диференціальних рівнянь
у похідних похідних.

\section{Розв'язання проблеми
Коші}\label{ux440ux43eux437ux432ux44fux437ux430ux43dux43dux44f-ux43fux440ux43eux431ux43bux435ux43cux438-ux43aux43eux448ux456}

\begin{verbatim}
Розглянемо проблему поширення тепла в *нескінченному* стрижні, якщо задано початкову температуру u(x, 0) = \varphi(x) .
\end{verbatim}

Цю задачу зазвичай називають задачою Коші для рівняння теплопровідності
або задачою початкових умов

\[
\begin{array}{lll} (\text{УЧП}) & u_t = \alpha^2 u_{xx}, & -\infty < x < \infty, & 0 < t < \infty, \\ (\text{12.6}) & (\text{HV}) & u(x,0) = q(x), & -\infty < x < \infty. \end{array}
\]

Розв'язання задачі можна поділити на три основні етапи.

КРОК 1 (Конверсія завдань).

\begin{verbatim}
Оскільки просторова змінна x змінюється між -\infty і +\infty , ми піддамо рівняння та початкову умову (12,6) перетворенню Фур'є на змінній x^{(1)} (x — інтегрована змінна у перетворенні). Після цього ми отримуємо
\end{verbatim}

\[
\mathcal{F}[u_t] = \alpha^2 \mathcal{F}[u_{xx}],
\mathcal{F}[u(x, 0)] = \mathcal{F}[\varphi(x)].
\]

Воспользуемся свойствами преобразования Фурье. Тогда

\[
\frac{\frac{dU(t)}{dt} = -\alpha^2 \xi^2 U(t), }{U(0) = \Phi(\xi)} = (\text{here } \Phi - \phi) \text{rje-image funcin } \phi),
\]

где \(U(t) = \mathcal{F}[u(x, t)]\) . Читате ть должен обратить внимание
на то, что функция \(U(t)\) на салот доле зависит не только от t, но и
от \(\xi\) , но, поскольку в дифференциальном уравнении (12.7) величина
\(\xi\) играет роль константы, мы не будем записывать ее в число
аргументов, т. е. будем считать, что U = \(U(t)\).

ШАГ 2 (Решение преобразованной задачи).

Поскольку новая переменная играет роль гонстанты в задаче (12.7), ее
решение легко находится:

(12.8)

\[
U(t) = \Phi(\xi) e^{-x^2 e^{tt}}.
\]

ШАГ 3 (Нахождение обратного преобразования). Искомое решение \(u(x, t)\)
находится по формуле

\[
u(x, t) = \mathcal{F}^{-1} [U(\xi, t)] =
= \mathcal{F}^{-1} [\Phi(\xi) e^{-\alpha \xi^2 t}].
\]

Припускається, що існує перетворення Фур'є цих функцій. Примітка. Ед.

Вот здесь нас выручит теорема о свертке (12.5). Используя эту теорему,
получаем

(12.9)

\[
u(x, t) = \mathcal{F}^{-1} \left[ \Phi(\xi) e^{-\alpha r \xi^{2} t} \right] =
= \mathcal{F}^{-1} \left[ \Phi(\xi) \right] * \mathcal{F}^{-1} \left[ e^{-\alpha^{2} \xi^{2} t} \right] =
= \Phi(x) * \left[ \frac{1}{2\alpha \sqrt{\pi t}} e^{-x^{2}/4\alpha^{2} t} \right] \text{ (знайдено в таблиці)} =
= \frac{1}{2\alpha \sqrt{\pi t}} \int_{-\infty}^{+\infty} \Phi(\xi) e^{-(x-\xi)^{2}/4\alpha^{2} t} d\xi.
\]

Дело сделано! Формула (12.9) дает решение нашей задачи. Прежде чем
поставить точку, давайте проанализируем полученный результат. Отметим,
что почынтегральная функция состои с

из двух сомножителей:

\begin{enumerate}
\def\labelenumi{\arabic{enumi}.}
\item
  Початкової температури \(\phi(\xi)\),
\item
  Функції
\end{enumerate}

\[
G(x, t) = \frac{1}{2\alpha \sqrt{\pi t}} e^{-(x-\xi)^2/4\alpha^2 t},
\]

яка називається функцією Гріна або функцією джерела.

Можна показати, що функція джерела \(G(x, t)\) описує відгук системи на
початковий одиничний температурний імпульс у точці \(x=\xi\).

\pandocbounded{\includegraphics[keepaspectratio]{images/_page_94_Figure_9.jpeg}}

Рис. 12.2. Відгук \(G(x,t)\) на температурний імпульс у точці \(x=\xi\).
Графік \(G(x,t)\) подібний до кривої нормального розподілу; час відіграє
роль середньоквадратичного відхилення: при малих \(t\) функція вузька, а
зі зростанням \(t\) помітно розширюється.

Іншими словами, функція \(G(x, t)\) описує розподіл температури в
стрижні в момент часу \(t\), якщо на точку \(x=\xi\) подіяв одиничний
тепловий імпульс (див. рис. 12.2).

Тепер формулі (12.9) можна дати таку інтерпретацію: початкову
температуру \(u(x,0)=\phi(x)\) можна розглядати як неперервну сукупність
точкових імпульсів величини \(\phi(\xi)\) у точці \(x=\xi\). Кожен
точковий імпульс дає розподіл температури \(\phi(\xi)G(x,t)\).
Підсумковий розподіл знаходиться сумуванням (інтегруванням) температур
від точкових джерел за формулою (12.9). Пізніше ми побачимо, що це один
із проявів загального принципу суперпозиції.

\section{ЗАМЕЧАНИЕ}\label{ux437ux430ux43cux435ux447ux430ux43dux438ux435}

Главный недостаток преобразования Фурье в том, что его можно применять
не ко всем функциям. Например, функции вида

\[
f(x) = \text{const},
f(x) = e^x,
f(x) = \sin x
\]

не имеют фурье-образа, так как интеграт

\[
\mathcal{F}[f] = \frac{1}{\sqrt{2\pi}} \int_{-\pi}^{\pi} f(x) e^{-i\xi x} dx
\]

расходится. Фурье-образ имеют только те функции, которые достаточно
быстро сгремятся к нулю при \(|x| \to \infty\) . Кроме того, нельзя
применять преобразование Фурье по времени для решения задач с начальными
условиями, так как \(0 < t < \infty\) (если не рассматривать продолжения
функций на полуось в соответствующем классе функций). \# ЗАДАЧІ 1.
Найдите фурье-образ функции

\[
f(x) =
\begin{cases} 0, & x \leq 0, \\ e^{-x}, & 0 < x, \end{cases}
\]

используя таблицы, приведенные в приложении.

\begin{itemize}
\tightlist
\item
  \begin{enumerate}
  \def\labelenumi{\arabic{enumi}.}
  \setcounter{enumi}{1}
  \tightlist
  \item
    На 10м переконайтеся, що прямі та обернені перетворення Фур'є є
    лінійними перетвореннями.
  \end{enumerate}
\item
  \begin{enumerate}
  \def\labelenumi{\arabic{enumi}.}
  \setcounter{enumi}{2}
  \tightlist
  \item
    Використовуючи перетворення Фур'є, розв'яжете задачу Коші
  \end{enumerate}
\end{itemize}

\[
\begin{array}{lll} (\text{УЧП}) & u_t = \alpha^2 u_{xx}, & -\infty < x < \infty, \\ (\text{HV}) & u(x, 0) = e^{-x^2}, & -\infty < x < \infty. \end{array}
\]

\begin{enumerate}
\def\labelenumi{\arabic{enumi}.}
\setcounter{enumi}{3}
\tightlist
\item
  Проверьте справедливость соотношений
\end{enumerate}

\[
\begin{aligned} \mathcal{F}\left[u_{x}\right] &= i\xi \mathcal{F}\left[u\right], \\ \mathcal{F}\left[u_{xx}\right] &= -\xi^{2} \mathcal{F}\left[u\right]. \end{aligned}
\]

УКАЗАНИЕ, Воспользуйтесь интегрированием по частям.

\begin{enumerate}
\def\labelenumi{\arabic{enumi}.}
\setcounter{enumi}{4}
\tightlist
\item
  Перевірити, чи можна записати згортку двох функцій f і g у двох
  еквівалентних формах
\end{enumerate}

\[
(f * g)(x) = \frac{1}{\sqrt{2\pi}} \int_{-\infty}^{+\infty} f(\xi) g(x - \xi) d\xi
\]

или

\[
(f * g)(x) = \frac{1}{\sqrt{2\pi}} \int_{-\infty}^{+\infty} f(x-\xi) g(\xi) d\xi.
\]

Лекция 13

\section{МЕТАЛЕКЦІЇ З ТРАНСФОРМАЦІЇ
ЛАПЛАСА}\label{ux43cux435ux442ux430ux43bux435ux43aux446ux456ux457-ux437-ux442ux440ux430ux43dux441ux444ux43eux440ux43cux430ux446ux456ux457-ux43bux430ux43fux43bux430ux441ux430}

\[
\mathscr{L}[f] = F(s) = \int\limits_0^\infty f(t)\,e^{-st}dt
\]

(преобразование Лапласа),

ознайомтеся з формулою зворотного перетворення Лапласа

\[
\mathscr{L}^{-1}[F] = f(t) = \frac{1}{2\pi t} \int_{c-t\infty}^{c+t\infty} F(s) e^{st} ds
\]

(обратное преобразование Лапласа)

и проиллюстрировать их свойства.

\begin{verbatim}
Перетворення Лапласа має важливу перевагу над перетворенням Фур'є, оскільки субінтегративний вираз містить швидко затухаючий коефіцієнт e^{-st} , і тому перетворення Лапласа застосовне до ширшого класу функцій порівняно з перетворенням Фур'є. Крім того, перетворення Лапласа працює з функціями, визначеними в інтервалі [0, \infty) , що робить його особливо зручним для змінної t.
\end{verbatim}

Ознайомившись із базовими властивостями перетворення Лапласа ми
розв'яжемо кілька важливих задач для диференціальних рівнянь у похідних
113 усіх інтегральних перетворень, розглянутих у цій книзі ---
перетворення Лапласа

(13 1)

\[
\mathscr{Z}[f] = \int_{0}^{\infty} f(t) e^{-st} dt,
\]

\begin{verbatim}
По-перше, існує один із найскладніших методів перекладу найважливіших диференціальних рівнів на алгебранхічних рівнях. Крім того, відтворення Лапласа дозволяє перетворити диференціальний рівень у попередній диференціальний рівень, і можна змінити його настільки, x, y, z, t \ldots що можливо змінити з 0 на \infty (найпоширеніше відтворення Лапласа — змінювати з 0 на ).
\end{verbatim}

Якщо перетворення Лапласа зупинити до найвищої диференціальної рівності,
то у випадку падіння часткових диференціальних рівнів перетворення
Лапласа можливо перевести його на новий диференціальний рівень з меншою
кількістю змінних або на значущий диференціальний рівень. Перше, що
потрібно зробити --- це розвинути вже виконані завдання, дії влади
роботи Лапласа

\section{Властивості перетворення
Лапласа}\label{ux432ux43bux430ux441ux442ux438ux432ux43eux441ux442ux456-ux43fux435ux440ux435ux442ux432ux43eux440ux435ux43dux43dux44f-ux43bux430ux43fux43bux430ux441ux430}

Свойство 1 (Прямое и обратное преобразования)

Перетворення Лапласа та формула його інверсії мають вигляд

\[
\mathscr{L}[f] = F(s) = \int_{0}^{\infty} f(t) \, e^{-st} \, dt \quad \text{(Перетворення Лапласа),}
\]

(13.2)

\[
\mathscr{L}^{-1}[F] = f(t) = \frac{1}{2\pi i} \int_{c-i\infty}^{c+i\infty} F(s) \, e^{st} \, ds \quad \text{(формула інверсії).}
\]

\begin{verbatim}
Перетворення Лапласа має важливу перевагу над Фур'є-перетворенням. Ця перевага зумовлена наявністю затухаючого e^{-st} множника в субінтегративних членах, що дозволяє застосовувати перетворення Лапласа до ширшого класу функцій ( e^{t\xi t} множник у перетворенні Фур'є не згас, оскільки модуль цього множника дорівнює одиниці). Більш точне уявлення про цей клас функцій дає наступна теорема:
\end{verbatim}

\section{Достатні умови для існування перетворення
Лапласа}\label{ux434ux43eux441ux442ux430ux442ux43dux456-ux443ux43cux43eux432ux438-ux434ux43bux44f-ux456ux441ux43dux443ux432ux430ux43dux43dux44f-ux43fux435ux440ux435ux442ux432ux43eux440ux435ux43dux43dux44f-ux43bux430ux43fux43bux430ux441ux430}

Ести

\begin{verbatim}
1) для будь-якого додатного числа A функція f(t) є кусочно-неперервною в інтервалі 0 \le t \le A ;
\end{verbatim}

\begin{verbatim}
2) існують такі константи M, a і T, які |f(t)| \leq Me^{at} для всіх t > T, тоді перетворення Лапласа
\end{verbatim}

\[
\mathscr{L}[f(\ell)] = F(s) = \int_{0}^{\infty} f(\ell) e^{-st} dt
\]

существует при всех s \textgreater{} a.

Приведем примеры некоторых функций и их преобразований Лапласа

\begin{verbatim}
У визначенні перетворення Лапласа мається на увазі, що s є дійсною змінною, що належить інтервалу (0, \infty) . Однак можливо (і часто бажано) розширити це визначення на комплексні значення s
\end{verbatim}

\[
\mathscr{L}^{-1}[F] = \frac{1}{2\pi i} \int_{\epsilon - i\infty}^{\epsilon + i\infty} F(s) e^{st} ds.
\]

Використання цієї формули передбачає знайомство з технікою інтегрування
вздовж контурів у комплексній площині та теорією дедукцій. Ми не будемо
відволікатися на вивчення цих питань тут, і коли йдеться про
перетворення Лапласа, ми використаємо таблицю, наведену в додатку.

\section{Властивість 2 (Перетворення часткових
похідних)}\label{ux432ux43bux430ux441ux442ux438ux432ux456ux441ux442ux44c-2-ux43fux435ux440ux435ux442ux432ux43eux440ux435ux43dux43dux44f-ux447ux430ux441ux442ux43aux43eux432ux438ux445-ux43fux43eux445ux456ux434ux43dux438ux445}

\begin{verbatim}
Нехай u(x, t) — функція змінних, і ми хочемо знайти результати застосування перетворення Лапласа до різних часткових похідних u_t , u_{tt} , u_x , u_{\lambda x} . . . . Оскільки перетворення Лапласа виконується змінною t (t — змінна інтегрування), правила перетворень мають вигляд
\end{verbatim}

\[
\mathcal{L}\left[u_{t}\right] = \int_{0}^{\infty} u_{t}(x, t) e^{-st} dt = sU(x, s) - u(x, 0),
\mathcal{L}\left[u_{tt}\right] = \int_{0}^{\infty} u_{tt}(x, t) e^{-st} dt = s^{2}U(x, s) - su(x, 0) - u_{t}(x, 0),
\mathcal{L}\left[u_{x}\right] = \int_{0}^{\infty} u_{x}(x, t) e^{-st} dt = \frac{\partial U}{\partial x}(x, s),
\mathcal{L}\left[u_{xx}\right] = \int_{0}^{\infty} u_{xx}(x, s) e^{-st} dt = \frac{\partial^{2}U}{\partial x^{2}}(x, s),
\]

\begin{enumerate}
\def\labelenumi{\arabic{enumi}.}
\tightlist
\item
\end{enumerate}

\[
f(t) = 1
\]

, \(0 < t < \infty\) \((M = 1, a = 0)\)

\pandocbounded{\includegraphics[keepaspectratio]{images/_page_99_Figure_3.jpeg}}

\[
F(s) = \int_{0}^{\infty} e^{-st} dt = \frac{1}{s}
\]

\pandocbounded{\includegraphics[keepaspectratio]{images/_page_99_Figure_5.jpeg}}

2

\[
f(t) = e^{2t}
\]

, \(0 < t < \infty\)\\
(M = 1, \(n = 2\) )

\pandocbounded{\includegraphics[keepaspectratio]{images/_page_99_Figure_7.jpeg}}

\[
F(s) = \frac{1}{s-2}
\]

, \(s > 2\)

a

\pandocbounded{\includegraphics[keepaspectratio]{images/_page_99_Figure_9.jpeg}}

\begin{enumerate}
\def\labelenumi{\arabic{enumi}.}
\setcounter{enumi}{2}
\tightlist
\item
\end{enumerate}

\[
f(t) = \sin(\omega t)
\]

\((M=1, \alpha=0)\)

\pandocbounded{\includegraphics[keepaspectratio]{images/_page_99_Figure_11.jpeg}}

\pandocbounded{\includegraphics[keepaspectratio]{images/_page_99_Figure_12.jpeg}}

\pandocbounded{\includegraphics[keepaspectratio]{images/_page_99_Figure_13.jpeg}}

\begin{enumerate}
\def\labelenumi{\arabic{enumi}.}
\setcounter{enumi}{3}
\tightlist
\item
  \(f(t) = e^{t^2}\) (не имеет преобразования Лапласа) Рис. 13.1.
  Графики некоторых функций и их изображений по Лапласу.
\end{enumerate}

где \(U\left(x,\,s\right)=\mathcal{L}\left[u\left(x,\,t\right)\right]\)
. Правила преобразования производных \(u_{x}\) и \(u_{xy}\) вытекают из
теоремы о дифференцировании интеграла по параметру

\[
\frac{d}{dx}\int_{a}^{b}f\left(x,\ y\right)dy=\int_{a}^{b}\frac{df}{dx}\left(x,\ y\right)dy,
\]

а формулы для преобразования производных \(u_t\) и \(u_{tt}\) можно
получить интегрированием по частям.

\section{Свойство 3 (Теорема о
свертке)}\label{ux441ux432ux43eux439ux441ux442ux432ux43e-3-ux442ux435ux43eux440ux435ux43cux430-ux43e-ux441ux432ux435ux440ux442ux43aux435}

Свертка играет здесь ту же роль, что и в теории преобразования Фурье, но
вводится несколько иначе.

\section{Определение конечной
свертки}\label{ux43eux43fux440ux435ux434ux435ux43bux435ux43dux438ux435-ux43aux43eux43dux435ux447ux43dux43eux439-ux441ux432ux435ux440ux442ux43aux438}

Конечная свертка двух функций f и g определяется формулами

\[
(f * g) (t) =
\begin{cases} \int_0^t f(\tau) g(t-\tau) d\tau, \\ \int_0^t f(t-\tau) g(\tau) d\tau \end{cases}
\]

\begin{verbatim}
(Можна довести, що ці інтеграли рівні). Тут інтегрування здійснюється вздовж скінченного інтервалу від 0 до t, а не вздовж нескінченної ( \infty , +\infty ), як це було раніше з визначенням нескінченної згортки.
\end{verbatim}

Ось приклад скінченної згортки двох функцій. Нехай

\[
f(t) = t, \quad g(t) = t;
\]

тоді

\[
(f * g)(t) = \int_{0}^{t} \tau(t - \tau) d\tau = t^{3}/6.
\]

Як і у випадку нескінченної згортки, її найважливіша властивість
визначається формулою

(13.4)

\[
\mathscr{L}[f * g] = \mathscr{L}[f] \mathscr{L}[g]
\]

або еквівалентна формула

(13.5)

\[
\mathscr{L}^{-1}\left\{ \left[\mathscr{L}\left[f\right]\mathscr{L}\left[g\right]\right\} = f*g. \right.
\]

\begin{verbatim}
Ці формули дозволяють змінити перетворення Лапласа у випадках, коли образ можна представити як добуток \mathscr{L}[f] і \mathscr{L}[g] таких множників, для яких
\end{verbatim}

Функції F і G легко знайти. Після пошуку функцій f і g можна знайти їхню
згортку. Наприклад,

\[
\mathcal{Z}^{-1}\left[\frac{1}{s} \cdot \frac{1}{s^2+1}\right] = \int_0^t \sin \tau \, d\tau = 1 - \cos t,
F(s) = \frac{1}{s} \xrightarrow{\mathcal{Z}^{-1}} f(t) = 1 \quad G(s) \xrightarrow{\mathcal{Z}^{-1}} g(t) = \sin t.
\]

Тепер ми готові до важливого змішаного виклику.

\section{Теплопровідність у напівнескінченному
середовищі}\label{ux442ux435ux43fux43bux43eux43fux440ux43eux432ux456ux434ux43dux456ux441ux442ux44c-ux443-ux43dux430ux43fux456ux432ux43dux435ux441ux43aux456ux43dux447ux435ux43dux43dux43eux43cux443-ux441ux435ux440ux435ux434ux43eux432ux438ux449ux456}

Розглянемо глибокий резервуар з рідиною, і нехай бічна поверхня
резервуара буде теплоізольованою. Нехай \(u_0\) буде початковою
температурою рідини, а температура повітря над рідиною ---

\pandocbounded{\includegraphics[keepaspectratio]{images/_page_101_Picture_7.jpeg}}

Рис. 13.2. Схематичне представлення проблеми теплопровідності: a ---
тепло тече всередину, якщо ви(0, t) \textless{} 0, і назовні, якщо ви(0,
t) \textgreater{} 0; b --- резервуар настільки глибокий, що гранична
умова внизу не впливає на розв'язок при значеннях x.

кістка дорівнює нулю (орієнтир температур) (див. рис. 13.2). Наше
завдання --- знайти температурне поле в рідині на різних глибинах і в
різний час. Іншими словами, нам потрібно розв'язати проблему

(13.6)

\begin{verbatim}
\begin{aligned}
(\text{УЧП})\quad & u_t = u_{xx}, && 0 < x < \infty,\; 0 < t < \infty, \\
(\Gamma \text{У})\quad & u_x(0,t) - u(0,t) = 0, && 0 < t < \infty, \\
(\text{НУ})\quad & u(x,0) = u_0, && 0 < x < \infty.
\end{aligned}
\end{verbatim}

Щоб розв'язати цю задачу, ми застосуємо перетворення Лапласа до змінної
t. Зверніть увагу, що можна також застосувати перетворення Лапласа до
змінної x, оскільки воно змінюється від 0 до \(\infty\) . Після
застосування перетворення

Ми отримуємо звичайне диференціальне рівняння за змінною X.

(13.7)

\[
(ODY) \quad sU(x) - u_0 = \frac{d^2U}{dx^2}, \quad 0 < x < \infty,
(\Gamma Y) \quad \frac{dU}{dx}(0) = U(0)
\]

\begin{verbatim}
(Ми перетворили лише рівняння та граничну умову, а не початкову умову.) Маємо звичайне диференціальне рівняння другого порядку, але з однією граничною умовою при x=0. Насправді існує ще одна гранична умова. Він визначається фізичними міркуваннями і фіксується як U(x) \to 0 , коли x \to +\infty . Зверніть увагу, що для спрощення позначення ми опустили параметр s із аргументів функції U(x) всюди, оскільки рівняння (13.7) є диференціальним рівнянням лише відносно x.
\end{verbatim}

Щоб розв'язати задачу (13.7), запишіть загальне рішення (загальний
розв'язок однорідного + часткового розв'язку гетерогенного) TAC:

\[
U(x) = c_1 e^{V_{5x}} + c_2 e^{-V_{5x}} + u_0/s.
\]

Підставляючи цей вираз у граничну умову (13.7), можна визначити
константи \(c_1\) та \(c_2\) (одразу зрозуміло, що \(c_3 = 0\),

\pandocbounded{\includegraphics[keepaspectratio]{images/_page_102_Figure_7.jpeg}}

Рис. 13.3. Граф ймовірнісного інтеграла (eri) та його комплементарної
функції (eric).

\begin{verbatim}
Інакше температура буде зростати нескінченно зі збільшенням координати X). Визначаючи сталу c_2 з граничної умови при x=0, отримуємо фінальний вираз для U\left(x\right) (13.8)
\end{verbatim}

\[
U(x) = u_0 \left\{ \frac{e^{-V \widehat{s} x}}{s(V \widehat{s} + 1)} \right\} + \frac{1}{s}.
\]

\begin{verbatim}
Залишається зробити останній крок. Щоб визначити температурне поле, потрібно обчислити u\left(x,\,t\right)
\end{verbatim}

\[
u(x, t) = \mathcal{L}^{-1}[U(x, s)]
\]

(ми повернулися до написання всіх аргументів функції \(U(x, s)\)). Щоб
повернути перетворення Лапласа, ми використаємо таблиці, наведені в
додатку. У результаті ми отримуємо

(13.9)

\[
u(x, t) = u_0 - u_0 \left[ \text{eric} \left( \frac{x}{2} \sqrt{t} \right) + \text{eric} \left( \sqrt{t} + \frac{x}{2} \sqrt{t} \right) e^{x+t} \right],
\]

де

\[
\operatorname{eric}(x) = \frac{2}{\sqrt{\pi}} \int_{x}^{\infty} e^{-\xi^{2}} d\xi
\]

є функцією, доповнюючою до ймовірнісного інтеграла. Графіки цієї функції
показані на рис. 13.3. Аналіз цього рішення P H C

\pandocbounded{\includegraphics[keepaspectratio]{images/_page_103_Figure_7.jpeg}}

13.4. Температура всередині напівнескінченного середовища в різний час.

Потребує дуже малих зусиль. Фіг. (13.4) представлено результати
чисельного аналізу, отриманого на комп'ютері, оснащеному плоттером.

\section{ПРИМІТКИ}\label{ux43fux440ux438ux43cux456ux442ux43aux438-1}

\begin{itemize}
\tightlist
\item
  \begin{enumerate}
  \def\labelenumi{\arabic{enumi}.}
  \tightlist
  \item
    Перетворення Лапласа також можна звести до неоднорідного
    диференціального рівня з частковими множниками (для виконання методу
    диференціації різних рівнів можна звести до одного рівня), але можна
    звести його до рівня різних множників (метод диференціації можна
    звести до рівня коефіцієнтів заміщення). Таблиця 1. 13.2 У сфері
    захисту шкіри існує два методи.
  \end{enumerate}
\item
  \begin{enumerate}
  \def\labelenumi{\arabic{enumi}.}
  \setcounter{enumi}{1}
  \tightlist
  \item
    Для виконання цих і межових завдань також проводиться трансформація
    Ганкеля і Мелліна, а також трансформація Лапласа в одному аспекті.
    Відтворення Лапласа підстановки оперативного множення формули
  \end{enumerate}
\end{itemize}

\[
\mathscr{L}[y'] = s\mathscr{L}[y] - y(0),
\]

Таблиця 13.2 Порівняння методів розділення перетворення Лапласа та
змінних

\textbar{} Метод \textbar{} \textbar{}

\textbar-------------------------\textbar---------------------------\textbar----------------------------------\textbar{}

\textbar{} Перетворення Лапласа \textbar{} Розділення \textbar{}

Гетерогенні PPP \textbar{} Так \textbar{} Так \textbar{}

Гетерогенні ГІ \textbar{} Так \textbar{} Так \textbar{}

Змінні ймовірності \textbar{} Так \textbar{} de' \textbar{}

Нелінійні рівняння \textbar{} Так \textbar{} Так \textbar{}

а перетворення Ганкеля і Мелліна замінюють диференціальний оператор на
множник. Наприклад, для перетворення Ганкеля, визначеного формулою

\[
H[y] = \int_{0}^{\infty} r J_{o}(\xi r) y(r) dr,
\]

справедливе співвідношення

\[
H\left[y''\left(r\right)+\frac{1}{r}\ y'\left(r\right)\right]=-\xi^{2}H\left[y\right].
\]

Таким чином можна розв'язати спеціальне диференціальне рівняння з
змінними коефіцієнтами (рівняння Бесселя). 3. Перетворення Лапласа на
змінній t можна інтерпретувати як проекцію площини xt на осі x, що
призводить до У ц ь о м у в

\pandocbounded{\includegraphics[keepaspectratio]{images/_page_104_Figure_9.jpeg}}

ипадку початкові рівняння, граничні та початкові умови перетворюються на
нове диференціальне рівняння та нові граничні умови (див. діаграму
вище).

\section{ЗАВДАННЯ}\label{ux437ux430ux432ux434ux430ux43dux43dux44f-10}

\begin{enumerate}
\def\labelenumi{\arabic{enumi}.}
\tightlist
\item
  Переконайтеся, що наступна формула для перетворення часткової похідної
  \(u_t\) є істинною:
\end{enumerate}

\[
\mathscr{L}\left[u_{t}(x, t)\right] = sU(x, s) - u(x, 0).
\]

\begin{enumerate}
\def\labelenumi{\arabic{enumi}.}
\setcounter{enumi}{1}
\tightlist
\item
  Використати перетворення Лапласа для розв'язання задачі Коші
\end{enumerate}

\begin{verbatim}
\begin{aligned}
(\text{УЧП})\quad & u_t = \alpha^2 u_{xx}, && -\infty < x < \infty,\ 0 < t < \infty, \\
(\text{НУ})\quad & u(x,0) = \sin x, && -\infty < x < \infty.
\end{aligned}
\end{verbatim}

\begin{enumerate}
\def\labelenumi{\arabic{enumi}.}
\setcounter{enumi}{2}
\tightlist
\item
  За допомогою перетворення Лапласа за змінною t розв'язати задачу
\end{enumerate}

\begin{verbatim}
\begin{aligned}
(\text{УЧП})\quad & u_t = u_{xx}, && 0 < x < \infty,\ 0 < t < \infty, \\
(\text{ГУ})\quad & u(0,t) = \sin t, && 0 < t < \infty, \\
(\text{НУ})\quad & u(x,0) = 0, && 0 \leqslant x < \infty.
\end{aligned}
\end{verbatim}

Дайте фізичну інтерпретацію цієї задачі.

\begin{enumerate}
\def\labelenumi{\arabic{enumi}.}
\setcounter{enumi}{3}
\tightlist
\item
  Розв'язати задачу з граничним значенням
\end{enumerate}

(ОДА)

\[
\frac{d^2U}{dx^2} - sU = A, \quad 0 < x < 1,
(\Gamma Y) \quad
\begin{cases} \frac{dU}{dx} (0) = 0 \\ U(1) = 0. \end{cases}
\]

Лекція 14

\section{ПРИНЦИП
ДЮАМЕЛЯ}\label{ux43fux440ux438ux43dux446ux438ux43f-ux434ux44eux430ux43cux435ux43bux44f}

МЕТА ЛЕКЦІЇ: Показати, як переклад Лапласа можна використовувати для
ідентифікації характеристик диференціальних рівнів. Зокрема, за
допомогою алгебранічного перетворення образів Лапласа, розвиток рівня
часткового оберненого можна відкинути принципом погляду Дюамеля.
Інтерпретація цього принципу в теорії найвищих диференціальних рівнів
світу використовується для розробки певних завдань УПП.

технології.

Давайте подивимось, як трансформація Лапласа здійснюється спробою
розвитку часткових диференціальних ліній, що дозволяє нам розкрити
фізичну сутність різних розвитків. За допомогою трансформації Лапласа ми
вже розробили принцип Дюамеля. Перше, що потрібно зробити --- звернутися
до принципу Дюамеля, ми бачимо проблему, яка часто є проблемою

\section{Теплопровідність у стрижні, кінці якого підтримуються при
заданій
температурі}\label{ux442ux435ux43fux43bux43eux43fux440ux43eux432ux456ux434ux43dux456ux441ux442ux44c-ux443-ux441ux442ux440ux438ux436ux43dux456-ux43aux456ux43dux446ux456-ux44fux43aux43eux433ux43e-ux43fux456ux434ux442ux440ux438ux43cux443ux44eux442ux44cux441ux44f-ux43fux440ux438-ux437ux430ux434ux430ux43dux456ux439-ux442ux435ux43cux43fux435ux440ux430ux442ux443ux440ux456}

Дуже часто необхідно знайти температуру всередині області, якщо відомо,
що на межі області вона змінюється з часом відповідно до певного закону
(граничні умови, залежні від часу).

На лівому кінці підтримується умова \(u(0,t)\)=0, а на правому кінці
температура задається функцією \(f(t)\).

Рис. 14.1. Граничні умови, що залежать від часу.

Наприклад, розглянемо термоізольований стержень, правий кінець якого
підтримується у \(f(t)\) (див. рисунок 14.1). Потрібно вирішити цю
проблему

(14.1)

\[
u_{t} = u_{xx}, \quad 0 < x < 1, \quad 0 < t < \infty,
(\Gamma Y) \quad \begin{cases} u(0, t) = 0, \\ u(1, t) = f(t), \end{cases} \quad 0 < t < \infty,
(HY) \quad u(x, 0) = 0, \quad 0 \le x \le 1.
\]

Сподіваємося, що розв'язання цієї задачі (14.1) буде легким, оскільки ми
знаємо розв'язок простішої задачі, де температура в кінці є сталою:

Насправді, якщо обидві задачі ((14.1) і (14.2)) розв'язуються одночасно
за допомогою перетворення Лапласа, то отримаємо визначний результат:
розв'язок задачі (14.1) виражається через розв'язання задачі (14.2).

Отже, розв'язуючи обидві задачі одночасно, ми отримуємо

\section{Легке завдання (14.2) (Постійні
солдати)}\label{ux43bux435ux433ux43aux435-ux437ux430ux432ux434ux430ux43dux43dux44f-14.2-ux43fux43eux441ux442ux456ux439ux43dux456-ux441ux43eux43bux434ux430ux442ux438}

Застосування перетворення -Лапласа до задачі (14.2)

\[
\frac{d^2W}{dx^2} - sW(x) = 0,
W(0) = 0,
W(1) = 1/s
\]

.

Пемеріль ODU

\[
W(x, s) = \frac{1}{s} \left[ \frac{\sinh(x \sqrt{s})}{\sinh(\sqrt{s})} \right].
\]

Пошук оберненого перетворення

\[
w(x, t) =
= x + \frac{2}{\pi} \sum_{n=1}^{\infty} \frac{(-1)^n}{n} e^{-(n\pi)^2 t} \times
\]

\(\times \sin(n\pi x)\) (розв'язання задачі з постійним ГІ).

\section{Важке завдання (14.1) (Солдати, що залежать від
часу)}\label{ux432ux430ux436ux43aux435-ux437ux430ux432ux434ux430ux43dux43dux44f-14.1-ux441ux43eux43bux434ux430ux442ux438-ux449ux43e-ux437ux430ux43bux435ux436ux430ux442ux44c-ux432ux456ux434-ux447ux430ux441ux443}

Застосування перетворення Лапласа до задачі (14.1)

\[
\frac{d^2U}{dx^2} - sU(x) = 0,
U\left( \mathbf{0}\right) =0,
U(1) = F(s)
\]

.

\[
U(x, s) = F(s) \left[ \frac{\sinh(x \sqrt{s})}{\sinh(\sqrt{s})} \right].
\]

Розділіть і помножте на z

\[
U(x, s) = F(s) \left\{ s \left[ \frac{\sinh(x \sqrt{s})}{s \sinh(s)} \right] \right\}.
\]

Використовуючи співвідношення \(\mathscr{L}[w_t] = sW - w(x, 0),\)
отримуємо

\[
U(x, s) = F(s) \mathcal{L}[w_t].
\]

Відповідно,

\[
u(x, t) = \mathcal{L}^{-1} \{ F(s) \mathcal{L}[w_t] \} =
= \mathcal{L}^{-1} [ F(s) ] * \mathcal{L}^{-1}[w_t] =
= f(t) * w_t(t) =
= \int_0^t w_t(x, t-\tau) f(\tau) d\tau =
\]

(після інтегрування частинами)

\[
= \int_{0}^{t} w(x, t-\tau) f'(\tau) d\tau + f(0) w(x, t).
\]

(Розв'язок задачі з часово-залежними ГІ виражається через розв'язок
задачі з постійними ГІ

розв'язано у задачах з константою

Отже, ми висловили розв'язок задач із граничними умовами, що змінюються
в часі, через розв'язання задачі з постійними граничними умовами.
Відповідні формули такі:

(14.3)

плутанина у наступних формулах. Записати диференціальне рівняння в
похідних у нових координатах (~, т), мы воспользуемся следующими
формулами перехода;

Наприклад, якщо dy/dx = B/2A = 3, то y - 3x = c, а отже
\(\xi(x, y) = y - 3x\) . Ця функція задовольняє умову
\([\xi_x/\xi_y] = -3\) і, відповідно, повертає коефіцієнт
\(\overline{A}\) до нуля.

Отже, половина роботи виконана: ми знайшли нову координатну
\(\xi = \xi(x, y)\) , яка забезпечила рівність \(\overline{A} = 0\) .
Залишається знайти \(\eta(x, y)\) таким чином, щоб перетворити
коефіцієнт \(\overline{B}\) на нуль.

Час зупинитися тут. Виявляється, що якщо вибрати \(B^2-4AC=0\) і \(\xi\)
так, щоб \(\overline{A}=0\) , то коефіцієнт \(\overline{B}\) автоматично
стане нулем 1). Тепер ми побачимо це.

Коефіцієнт \(\overline{B}\) задається виразом

\[
\overline{B}^2 - 4\overline{A}\overline{C} = (B^2 - 4AC) (\xi_x \eta_y - \xi_y \eta_x)^2 - \Pi \rho u u. \ \rho e \partial.
\overline{B} = 2 \left( V \overline{A} \xi_x + V \overline{C} \xi_y \right) \left( V \overline{A} \eta_x + V \overline{C} \eta_y \right).
\]

Функциональная диаграмма на рис. 15.5 делает эти формулы совсем
очевидными. Эта диаграмма особенно полезна при вычислении частных
производных функций и и и по переменным Р И С

\pandocbounded{\includegraphics[keepaspectratio]{images/_page_115_Picture_14.jpeg}}

15.5 Диаграмма, иллюстрирующая функциональную зависимость переменных

\begin{verbatim}
\xi   и   \tau  , поскольку   \xi   зависит и от x, и от t. Переменная   \tau   зависит только от t. Это все, что касается преобразований. Теперь подставим найденные выражения   u_t  ,   u_x   и   u_{vx}   в уравнение (15 2) и получим
\end{verbatim}

\[
\overline{B} = 2 \sqrt{\overline{A}} \left[ \sqrt{\overline{A}} \, \eta_x + \sqrt{\overline{C}} \, \eta_y \right].
\]

откуда

\[
u_{xy} + 2u_{xy} + u_{yy} = 0.
\]

.

\begin{verbatim}
Следовательно, новая задача с начальным условием в персменных   \xi   и   \tau   имеет вид
\end{verbatim}

(УЧП)

\[
\frac{dy}{dx} = -[\xi_x/\xi_y] = B/2A = 1.
\]

\begin{verbatim}
(НУ)   u(\xi, 0) = 1 - H(\xi), \quad -\infty < \xi < \infty.
\end{verbatim}

\begin{verbatim}
Так как   \xi = x   при t = 0, то новые граничные условия совпадают со старыми. Эта задача была решена в лекции 12 с помощью интегрального преобразования Фурье, и ее решение записывается в виде
\end{verbatim}

\[
y = x + c
\]

\begin{verbatim}
где   \phi(\beta)   — начальное условие. Следовательно, в нашем случае
\end{verbatim}

\[
\eta = y
\]

После подстановки

\begin{verbatim}
\begin{cases}
 \xi = y - x, \\
 \eta = y
 \end{cases}
\end{verbatim}

получаем окончательный результат

(15.3)

\[
\overline{C} = A \eta_x^2 + B \eta_x \eta_y + C \eta_y^2 = 1,
\overline{E} = A \eta_{xx} + B \eta_{xy} + C \eta_{yy} + D \eta_x + E \eta_y = 0,
\begin{cases} \frac{1}{2} \left[ 1 + \operatorname{erf} \left( \frac{-\xi}{2\sqrt{D\tau}} \right) \right], & \xi < 0, \\ \frac{1}{2} \operatorname{erfc} \left( \frac{\xi}{2\sqrt{D\tau}} \right), & 0 \leq \xi. \end{cases}
\overline{G} = G = 0.
\overline{A}u_{\xi\xi} + \overline{B}u_{\xi\eta} + \overline{C}u_{\eta\eta} + \overline{D}u_{\xi} + \overline{E}u_{\eta} + \overline{F}u = \overline{G}
\]

Графики этого решения для различных моментов времени изображены на рис.
15.6. Осталось сделать последний шаг --- записать решение нашей задачи в
координатах x и t. Подставляя

\[
u_{\eta} = f(\xi),
,   \tau = t
\]

в формулы (15.3), получаем

(15.4)

\[
u(x, t) = \begin{cases} \frac{1}{2} \left[ 1 + \operatorname{erf}\left(\frac{Vt - x}{2\sqrt{Dt}}\right) \right], & Vt > x, \\ \frac{1}{2} \operatorname{erfc}\left(\frac{x - Vt}{2\sqrt{Dt}}\right), & Vt \leq x. \end{cases}
\]

Перед нами розв'язок задачі конвективної дифузії (15.2). Профіль
концентрації рухається праворуч зі швидкістю \(V\) і водночас
розмивається дифузією з коефіцієнтом \(D\).

\paragraph{ЗАВДАННЯ}\label{ux437ux430ux432ux434ux430ux43dux43dux44f-11}

\begin{enumerate}
\def\labelenumi{\arabic{enumi}.}
\tightlist
\item
  Розв'яжіть задачу Коші
\end{enumerate}

\[
\begin{array}{ll} u_t = u_{xx} - 2u_x, & -\infty < x < \infty, \ 0 < t < \infty, \\ u(x, 0) = \sin x, & -\infty < x < \infty. \end{array}
\]

\begin{enumerate}
\def\labelenumi{\arabic{enumi}.}
\setcounter{enumi}{1}
\tightlist
\item
  Знайдіть розв'язок задачі Коші для рівняння конвективної дифузії
\end{enumerate}

\[
\begin{array}{ll} u_t = u_{xx} - 2u_x, & -\infty < x < \infty, \ 0 < t < \infty, \\ u(x, 0) = e^x \sin x, & -\infty < x < \infty. \end{array}
\]

\begin{enumerate}
\def\labelenumi{\arabic{enumi}.}
\setcounter{enumi}{2}
\tightlist
\item
  Знайдіть розв'язок задачі переносу, використовуючи перетворення
  координат з лекції 8.
\end{enumerate}

\subsection{Приведення еліптичних рівнянь до канонічної
форми}\label{ux43fux440ux438ux432ux435ux434ux435ux43dux43dux44f-ux435ux43bux456ux43fux442ux438ux447ux43dux438ux445-ux440ux456ux432ux43dux44fux43dux44c-ux434ux43e-ux43aux430ux43dux43eux43dux456ux447ux43dux43eux457-ux444ux43eux440ux43cux438}

Розглянемо загальне рівняння

\[
Au_{xx} + Bu_{xy} + Cu_{yy} + Du_x + Eu_y + Fu = G
\]

але тепер у випадку \(B^2 - 4AC < 0\). Переходом до нових незалежних
змінних ми хочемо звести його до форми

\[
u_{\xi\xi} + u_{\eta\eta} = \varphi(\xi, \eta, u, u_{\xi}, u_{\eta}).
\]

Щоб знайти ці нові змінні, виконуємо ті самі обчислення, що й у двох
попередніх випадках, отримуємо перетворене рівняння

\[
\overline{A}u_{\xi\xi} + \overline{B}u_{\xi\eta} + \overline{C}u_{\eta\eta} + \overline{D}u_{\xi} + \overline{E}u_{\eta} + \overline{F}u = \overline{G},
\]

і вимагаємо, щоб \(\xi\) та \(\eta\) задовольняли умовам
\(\overline{A} = \overline{C}\) і \(\overline{B} = 0\). Безпосередньо
знайти \(\xi\) та \(\eta\) тут складніше, тому шукаємо перетворення як
композицію двох кроків.

\section{Перетворення
1}\label{ux43fux435ux440ux435ux442ux432ux43eux440ux435ux43dux43dux44f-1}

Спочатку введемо комплексні координати \(\xi\) та \(\eta\), щоб надати
рівнянню вигляду

\[
u_{\xi\eta} = \psi(\xi, \eta, u, u_{\xi}, u_{\eta}).
\]

Для цього розв'язуємо характеристичні рівняння

\[
\frac{dy}{dx} = \frac{B - \sqrt{B^2 - 4AC}}{2A},
\]

де \(B^2 - 4AC < 0\), і отримуємо

\[
\xi(x, y) = \text{const}, \qquad \eta(x, y) = \text{const}.
\]

\section{Перетворення
2}\label{ux43fux435ux440ux435ux442ux432ux43eux440ux435ux43dux43dux44f-2}

Далі переходимо від \((\xi,\eta)\) до \((\alpha,\beta)\) за формулами

\[
\alpha = \frac{\xi + \eta}{2}, \qquad \beta = \frac{\xi - \eta}{2i}.
\]

У результаті отримуємо канонічну форму еліптичного рівняння.

\section{Зведення рівняння y\^{}2u\_\{xx\} + x\^{}2u\_\{yy\} = 0 до його
канонічної
форми}\label{ux437ux432ux435ux434ux435ux43dux43dux44f-ux440ux456ux432ux43dux44fux43dux43dux44f-y2u_xx-x2u_yy-0-ux434ux43e-ux439ux43eux433ux43e-ux43aux430ux43dux43eux43dux456ux447ux43dux43eux457-ux444ux43eux440ux43cux438}

Розглянемо рівняння

у яких \(A=y^2\) , B=0, \(C=x^2\) , D=E=F=G=0. Дискримінантний
\(B^2-4AC\) дорівнює \(-4x^2y^2\) , і ми зведемо це рівняння до
канонічної форми в першому квадранті x\textgreater0, y\textgreater0.

КРОК 1 (Перше перетворення)

Запишемо характеристичні рівняння

\[
\frac{dy}{dx} = \frac{B + \sqrt{B^2 - 4AC}}{2A} = \frac{\sqrt{-4x^2y^2}}{2y^2} = i\frac{x}{y}
\]

Замініть змінні \(\xi = x - Vt\) , \(\tau = t\) . У нашому випадку
інтеграл можна обчислити явно

\[
y^2 + ix^2 = \text{const},
\]

Такий інтеграл уже зустрічався у лекції 12. Можливо, читачу буде
зручніше звернутися до таблиць перетворення Фур'є.

\section{ОДНОВИМІРНЕ ХВИЛЬОВЕ РІВНЯННЯ (ГІПЕРБОЛІЧНІ
РІВНЯННЯ)}\label{ux43eux434ux43dux43eux432ux438ux43cux456ux440ux43dux435-ux445ux432ux438ux43bux44cux43eux432ux435-ux440ux456ux432ux43dux44fux43dux43dux44f-ux433ux456ux43fux435ux440ux431ux43eux43bux456ux447ux43dux456-ux440ux456ux432ux43dux44fux43dux43dux44f}

МЕТАЛЕКЦІЇ: Показати, що коливання навантаження описуються рівнянням

Від

і що це рівняння є наслідком законів Ньютона. Також розглядаються інші
типи хвильового рівняння:

\begin{verbatim}
\eta(x, y) = y^2 - ix^2. (Неважливо, яка з двох функцій називається \xi , а яка — \eta ; ви можете поміняти їх місцями, якщо хочете.) Після цього перетворення початкове рівняння набуває вигляду
\end{verbatim}

Ми не звертатимемо уваги на це рівняння (яке є комплексним гіперболічним
рівнянням), а одразу перейдемо до другого перетворення.

КРОК 2. (Друга трансформація)

При здійсненні другої трансформації ми маємо

\[
\begin{split} u_{tt} &= \alpha^2 u_{xx} + F\left(x, \ t\right), \\ u_{tt} &= \alpha^2 u_{xx} - \beta u_t - \lambda u + F\left(x, \ t\right). \end{split}
\]

(реальна частина \(\xi\) і \(\eta\) ),
\(\beta = \frac{\xi - \eta}{2\iota} = x^2\) (уявна частина \(\xi\) і
\(\eta\) ).

Щоб спростити нотацію, найкраще перепозначити \((\alpha,\beta)\)
\((\xi,\eta)\) окремо і записати отримане \(\epsilon\) перетворення як

КРОК 3 (Отримання нового рівняння)

\begin{verbatim}
Нову канонічну форму можна отримати, обчисливши всі коефіцієнти \overline{A}, \overline{B}, \overline{C}, \overline{D}, \overline{E}, \overline{F} і \overline{G} у рівнянні
\end{verbatim}

До цього часу ми розглядали фізичні процеси, які описуються одномірними
параболічними рівняннями (задачі дифузії). Тепер ми почнемо вивчати
інший клас диференціальних рівнянь у похідних --- гіперболічні рівняння.
Почнемо з одновимірного хвильового рівняння, яке описує, зокрема,
поперечні коливання струни.

\section{Вібрації
струн}\label{ux432ux456ux431ux440ux430ux446ux456ux457-ux441ux442ux440ux443ux43d}

Розглянемо невеликі коливання струни з нерухомими кінцями. Припустимо,
що струна щільно натягнута, зроблена з однорідного матеріалу і
коливається в одній площині (рис. 16.1).

\pandocbounded{\includegraphics[keepaspectratio]{images/_page_119_Figure_10.jpeg}}

Фіг. 16.1. Поперечні коливання струни: a --- зміщення струни з положення
рівноваги.

Ігноруємо вплив гравітації на струну. Щоб побудувати математичну модель
коливань струн, розглянемо всі сили, що діють на невелику частину струни
(рис. 16.2). Виявляється, що хвильове рівняння --- це не більше ніж
закон руху Ньютона (зміна імпульсу дорівнює сумі діючих сил),
застосований до струни. Дивлюся на рис. На рисунку 16-1 ми можемо
представити сили, що діють на струну, у напрямках, перпендикулярних осі
x.

\pandocbounded{\includegraphics[keepaspectratio]{images/_page_120_Figure_2.jpeg}}

Фіг. 16.2. Невелика частина \((x, x + \Delta x)\) осцилюючою струною.

\begin{enumerate}
\def\labelenumi{\arabic{enumi}.}
\tightlist
\item
  Сила, що виникає через натяг струни ( \(\alpha^2 u_{xx}\) ). Значення
  поперечної компоненти сили натягу визначається відповідно до формули
\end{enumerate}

Сила =

\begin{verbatim}
за формулами (41.5) з \xi = y^2 та \eta = x^2 . Зробивши це, отримуємо канонічну форму еліптичного рівняння
\end{verbatim}

\(\simeq T [u_x(x + \Delta x, t) - u_x(x, t)].\) 2. Внешняя сила
\(F (x, t)\).

Зовнішня сила \(F(x, t)\), прикладена до струни, може довольно залежати
від x і t. Нижче наведено деякі приклади зовнішніх сил:

а) гравитационная сила \(F(x, t)\) = -mg;

\begin{enumerate}
\def\labelenumi{\alph{enumi})}
\setcounter{enumi}{1}
\item
  імпульси, розподілені вздовж струни та діючи на неї в різний час:
\item
  сили, що виникають внаслідок дії звукової хвилі на шкіру барабана
  (пізніше ми детально вивчаємо коливання мембрани, які описуються
  хвильовим рівнянням).
\end{enumerate}

\begin{enumerate}
\def\labelenumi{\arabic{enumi}.}
\setcounter{enumi}{2}
\item
  Сили тертя, що діють на струну (--- \(\beta u_t\) ). Якщо струна
  вібрує у середовища, виникає сила опору, яка пропорційна швидкості
  \(u_t\) .
\item
  Возвращающая сила (- үи).
\end{enumerate}

Ця сила спрямована протилежно зміщенню струни. Якщо зміщення u додатне,
то сила від'ємна.

Якщо тепер застосувати рівняння руху Ньютона до невеликої ділянки
струни, отримаємо

\section{ПРИМІТКИ}\label{ux43fux440ux438ux43cux456ux442ux43aux438-2}

\begin{enumerate}
\def\labelenumi{\arabic{enumi}.}
\tightlist
\item
  Аналогічно можна класифікувати лінійні рівняння другого порядку з
  \emph{трьох або більше змінних}, але для цього потрібен матричний
  аналіз. Наприклад, рівняння другого порядку з трьома змінними
\end{enumerate}

де \(\rho\) --- лінійна густина струни. Поділимо обидві частини цього
співвідношення на \(\Delta x\) і поставимо \(\Delta x\) --- 0.
Результатом є добре відоме телеграфне рівняння

(16.1)

згідно з цією класифікацією, буде параболічним, а рівняння

У цьому рівнянні слід поділити \(\alpha\) , \(\beta\) , \(\gamma\) і
\(F(x,t)\) на \(\rho\) , але для простоти залишимо старі позначки для
цих величин. Отримано бажане рівняння. Ось інтерпретація найпростішого
хвильового рівняння

\begin{itemize}
\tightlist
\item
  Гіперболічно.
\end{itemize}

\begin{enumerate}
\def\labelenumi{\arabic{enumi}.}
\setcounter{enumi}{1}
\tightlist
\item
  Наш інтерес до класифікації диференціальних рівнянь у похідних
  пов'язаний із тим, що три основні типи рівнянь описують три різні типи
  фізичних явищ, і ми хотіли б отримати математичну класифікацію цих
  трьох класів фізичних явищ.
\end{enumerate}

\section{ЗАВДАННЯ}\label{ux437ux430ux432ux434ux430ux43dux43dux44f-12}

\begin{enumerate}
\def\labelenumi{\arabic{enumi}.}
\tightlist
\item
  Які з цих параболічних і еліптичних рівнянь записані у канонічній
  формі?
\end{enumerate}

\[
a) u_t = u_{xx} + hu,
б) u_{xy} + u_{xx} + 3u = \sin x,
в) u_{xx} + 2u_{yy} = 0,
г) u_{xx} = \sin^2 x?
\]

\begin{enumerate}
\def\labelenumi{\arabic{enumi}.}
\setcounter{enumi}{1}
\item
  Перетворіть параболічне рівняння \(u_{xx}\) + 2u\_\{xy\} + \(u_{yy}\)
  + u = 2 до канонічної форми.
\item
  Перетворіть еліптичне рівняння \(u_{xx}\) + 2u\_\{yy\} + x\^{}2
  \(u_x = e^{-x^2/2}\)до канонічної форми.
\end{enumerate}

\section{МЕТОД МОНТЕ-КАРЛО
(ВСТУП)}\label{ux43cux435ux442ux43eux434-ux43cux43eux43dux442ux435-ux43aux430ux440ux43bux43e-ux432ux441ux442ux443ux43f}

МЕТА ЛЕКЦІЇ: Пояснити основні ідени, що лежать в основі методу
Монте-Карло, і показати, як цей метод можна використовувати для
розв'язання різних задач. Виявляється, що за допомогою азартних ігор
можна змоделювати (загалом, на комп'ютері) наближені розв'язки реальних
задач. Як простий приклад можна вказати обчислення інтеграла

основанную на интуитивных соображениях.

\section{Інтуїтивна інтерпретація хвильового
рівняння}\label{ux456ux43dux442ux443ux457ux442ux438ux432ux43dux430-ux456ux43dux442ux435ux440ux43fux440ux435ux442ux430ux446ux456ux44f-ux445ux432ux438ux43bux44cux43eux432ux43eux433ux43e-ux440ux456ux432ux43dux44fux43dux43dux44f}

Читач може запитати: чому рівняння типу (16.2) має описувати щось
подібне до вібрацій скрипкової струни? Щоб зрозуміти це, необхідно
зрозуміти Ф і г

\pandocbounded{\includegraphics[keepaspectratio]{images/_page_121_Picture_9.jpeg}}

. 16 3. Інтерпретація хвильового рівняння \(u_{tt} = \alpha^2 u_{xx}\)
\(u_{tt}\) Equation (16.2) можна інтерпретувати так: прискорення струни,
викликане натягом, у кожній точці більше, чим більша увігнутість струни
\(u_{xx}\) у заданій точці (коефіцієнт пропорційності сталий
\(\alpha^2 = T/\rho\) ) (див. рис. 133). \# ПОВАГА 1. Хвильове рівняння
\(u_{tt} = \alpha^2 u_{xx}\) також описує поздовжні та крутильні
коливання стрижня. У поздовжніх осцилях зміщення паралельні елементу, і
значення поздовжнього зміщення відносно положення рівноваги позначається
через \(u\left(x,\,t\right)\) . Такі коливання виникають, наприклад,
коли молоток б'ють по кінці стрижня. Торсіональний

ные колебания

кидаючи дротик у одиничний квадрат \{(x, y): 0 \textless{} x \textless{}
1, 0 \textless{} y \textless{} 1\} Після приблизно 100 кидків відносна
частина дротиків, що опускається нижче кривої y = x\^{}2 , вважається
приблизним значенням інтеграла. Узагальнення такої гри (з випадковим
киданням дрогіка) завжди містить процедуру генерації послідовності
випадкових чисел.

У цій лекції ми коротко опишемо метод Монте-Карло, а в наступній лекції
покажемо, як його використовувати для розв'язання диференціальних
рівнянь у похідних

Методи Монте-Карло зазвичай називають групою методів розв'язання
детермінованих задач (тобто задач без випадковості), у яких суттєво
використовуються елементи випадковості. Загальна філософія етнічних
методів показана на рис. 42 1.

Ймовірнісна гра

\begin{verbatim}
Деторміноване завдання \rho — розв'язання задачі (наприклад, обчислення інтеграла, розв'язання UBP тощо).
\end{verbatim}

\begin{verbatim}
Результат = \hat{p} Наближення p Рішення = p Рис. 42.1. Загальна теорія методу Монте-Карло.
\end{verbatim}

\section{Обчислення
інтеграла}\label{ux43eux431ux447ux438ux441ux43bux435ux43dux43dux44f-ux456ux43dux442ux435ux433ux440ux430ux43bux430}

Припустимо, ми хочемо обчислити інтеграл

Тут \emph{k} --- це модуль Юнга, який характеризує еластичність
матеріалу. Матеріали з великим модулем Юнга коливаються на вищих
частотах. Звукові хвилі --- це поздовжні хвилі.

\begin{enumerate}
\def\labelenumi{\arabic{enumi}.}
\setcounter{enumi}{1}
\tightlist
\item
  Якщо лінійна густина струни \(\rho(x)\) залежить від координати, то
  хвильова аугура записується як
\end{enumerate}

(детерміноване завдання). Щоб використати метод Монте-Карло, потрібно
винайти гру на удачу, яка дає приблизне значення цього інтеграла.

\begin{verbatim}
Звісно, можна придумати багато різних ігор. Остаточна версія залежить від точності розрахунків, простоти гри тощо. Очевидний варіант гри для обчислення інтегралу — кинути дротик у прямокутний R = \{(x,y): a \le x \le b, 0 \le y \le \max f(x)\} (див. рис. 42.2)
\end{verbatim}

\pandocbounded{\includegraphics[keepaspectratio]{images/_page_315_Figure_7.jpeg}}

Рис. 42 2. Обчислення інтегралу методом Монте-Карло.

Очевидно, якщо випадковим чином кинути близько 100 дротиків у
прямокутник R, що містить графік функції, можна отримати оцінку значення
інтеграла, помноживши площу прямокутника R на відносну кількість
дротиків під графіком функції. Отже, результат нашої гри

т. е. получается УЧП с переменными коэффициентами.

\begin{enumerate}
\def\labelenumi{\arabic{enumi}.}
\setcounter{enumi}{2}
\tightlist
\item
  Оскільки хвильове рівняння \(u_{tt} = \alpha^2 u_{xx}\) містить
  похідну другого порядку часу, щоб отримати єдине розв'язання в точці
  t\textgreater0, необхідно вказати \(\partial sa\) початкові умови:
\end{enumerate}

--- це інтегральний бал 1.

Щоб виконати ці обчислення на комп'ютері, нам потрібно якимось чином
отримати послідовність випадкових точок (подивимось, як це буде зроблено
пізніше), тобто наказати комп'ютеру «кидати дротики». Блок-схема на рис.
42.3 показує, як комп'ютер вирішить цю проблему.

\pandocbounded{\includegraphics[keepaspectratio]{images/_page_316_Figure_2.jpeg}}

\begin{verbatim}
**РИС.** 42.3. Блок-схема розрахунку інтегрального \int_{a}^{b} f(x) dx методом Монте-Карло (100 кидків); M = \max_{a} f(x) .
\end{verbatim}

\section{Випадкові
числа}\label{ux432ux438ux43fux430ux434ux43aux43eux432ux456-ux447ux438ux441ux43bux430}

Перш ніж перейти до рівняння Монте-Карло з частковими похідними,
зупинимося на важливому питанні випадкових чисел. Щойно, обчислюючи
інтеграл, нам потрібно було побудувати послідовність випадкових точок
\(P_i = (x_i, y_i)\)всередині прямокутника R. Іншими словами,
x-координата точки P має бути випадковим числом з інтервалу {[}a, b{]},
а \(y_i\) від інтервалу {[}0, M{]}. Щоб знайти випадкові числа з цих
інтервалів, звертаємося до послідовності випадкових чисел, рівномірно
розподілених по сегменту {[}0, 1{]}. Тоді числа, рівномірно розподілені
на сегменті {[}a, b{]}, можна обчислити за формулою

(початкове зміщення струни), \(u_t(x, 0) = g(x)\) (початкова швидкість
струни).

Таким чином, хвильове рівняння відрізняється від рівняння
теплопровідності, для якого потрібно встановити лише одну початкову
умову.

\begin{enumerate}
\def\labelenumi{\arabic{enumi}.}
\setcounter{enumi}{3}
\tightlist
\item
  За допомогою хвильового рівняння можна описати розподіл електричного
  струму в дроті. З законів Кірхгофа отримуємо таку систему двох рівнянь
  у похідних похідних першого порядку:
\end{enumerate}

(16.3)

\begin{verbatim}
Тепер, звичайно, виникає питання: як отримати послідовність випадкових чисел, \{r_i: i=1,2,\ldots\} рівномірно розподілених на [0,1]? Найпоширенішим методом сьогодні є метод дедукцій (метод порівнянь). Використовуючи цей метод, отримуємо послідовність випадкових цілих чисел (наприклад, 2120, 1401, 177, 3013,...), потім ліворуч від кожного числа розміщується десяткова крапка, так що всі числа знаходяться між нулем і одним (0,212; 0,1401; 0,0177; 0,3013; \ldots) .
\end{verbatim}

Отже, щоб отримати послідовність випадкових цілих чисел (між 0 і P), ми
використаємо метод віднімання.

\section{Отримання випадкових чисел за допомогою методу
відрахування.}\label{ux43eux442ux440ux438ux43cux430ux43dux43dux44f-ux432ux438ux43fux430ux434ux43aux43eux432ux438ux445-ux447ux438ux441ux435ux43b-ux437ux430-ux434ux43eux43fux43eux43cux43eux433ux43eux44e-ux43cux435ux442ux43eux434ux443-ux432ux456ux434ux440ux430ux445ux443ux432ux430ux43dux43dux44f.}

\begin{enumerate}
\def\labelenumi{\arabic{enumi}.}
\item
  Як перше випадкове ціле число, обирайте будь-яке число, що знаходиться
  між 0 і P (число P обирається наперед).
\item
  Помножте це випадкове число на певне фіксоване число.
\end{enumerate}

Попередньо обраний множник M.

\begin{enumerate}
\def\labelenumi{\arabic{enumi}.}
\setcounter{enumi}{2}
\item
  Додайте до добутку деяке фіксоване ціле число \emph{K}, обране
  заздалегідь.
\item
  Поділити отриману суму на P і залишок
\end{enumerate}

Виберіть нове випадкове число.

\begin{verbatim}
Метод виведення можна записати у формі формули r_{i+1} \equiv (Mr_i + K) \bmod P , що означає, що для отримання наступного r_{i+1} потрібно взяти попередній r_i , помножити його на M, додати результат на K, поділити суму на P і взяти решту поділу. Наприклад, якщо P = 100, M = 37, K = 16, r_0 = 15 , то отримується наступна послідовність випадкових чисел:
\end{verbatim}

где

ж --- координата вдоль провода,

\emph{t} --- время,

\(i(x, t)\) --- распределение тока вдоль провода,
\(v\left(x,\ t\right)\) --- распределение потенциала вдоль провода,

С -- емкость провода на единицу длины,

G --- утечка на единицу длины,

Р --- сопротивление на единицу длины,

-- индуктивность провода на единицу длины.

Рівняння (16.3) зазвичай називають системою телеграфних рівнянь. У такій
формі вони зазвичай не використовуються. Щоб перетворити цю систему,
диференціюйте перше рівняння на x, друге на t, помножте обидві частини
на C і відніміть від першого. У результаті ми отримуємо

\[
Mr_1 + K = 2643
\]

Воспользуемся теперь вторым уравнением из (16.3)

\[
r = 71
\]

окончательно получаем

(16.4)

Y

Це рівняння для струму, відоме як телеграфне рівняння, є гіперболічним
рівнянням другого порядку (якщо, звісно, C і L не дорівнюють нулю,
інакше рівняння стає параболічним).

Напряжение описывается точно таким же уравнением:

(16.5)

\[
I = \int_{0}^{1} \int_{0}^{1} \int_{0}^{1} \int_{0}^{1} e^{-(x^{2}+y^{2}+z^{2}+\omega^{2})} dx dy dz dw,
\]

Если G = R = 0, то уравнения (16.4)---(16.5) упрощаются:

(16.6)

\[
I = \int_{0}^{1} e^{\sin x} dx.
\]

\section{\texorpdfstring{ЗАВДАННЯ - 1. Отримайте рівняння (16.5) для
\emph{and} з системи з двох рівнянь першого порядку
(16.3)}{ЗАВДАННЯ - 1. Отримайте рівняння (16.5) для and з системи з двох рівнянь першого порядку (16.3)}}\label{ux437ux430ux432ux434ux430ux43dux43dux44f---1.-ux43eux442ux440ux438ux43cux430ux439ux442ux435-ux440ux456ux432ux43dux44fux43dux43dux44f-16.5-ux434ux43bux44f-and-ux437-ux441ux438ux441ux442ux435ux43cux438-ux437-ux434ux432ux43eux445-ux440ux456ux432ux43dux44fux43dux44c-ux43fux435ux440ux448ux43eux433ux43e-ux43fux43eux440ux44fux434ux43aux443-16.3}

\begin{itemize}
\tightlist
\item
  \begin{enumerate}
  \def\labelenumi{\arabic{enumi}.}
  \setcounter{enumi}{1}
  \tightlist
  \item
    Знаючи фізичне значення кожного члена хвильового рівняння, що ви
    можете сказати про поведінку в момент розв'язання наступної задачі:
  \end{enumerate}
\end{itemize}

(YYII)

\[
r_{i+1} \equiv (3r_i + 4) \mod 7, \quad r_0 = 0.
\]

, \(0 < x < 1\) , \(0 < t < \infty\) ,\\
(ГУ)

\[
I = \int_{0}^{1} \int_{0}^{1} \int_{0}^{1} e^{-(x^{2}+y^{2}+z^{2})} dx dy dz.
u(x,y) = g(x,y) =
\begin{cases} u(0, t) = 0, & 0 < t < \infty, \\ u(1, t) = 0, & 0 < t < \infty, \end{cases}
\]

Після того, як цю конкретну задачу буде розв'язано методом Монте-Карло,
ми розглянемо, як розв'язувати більш загальні задачі за допомогою цього
методу.

Щоб проілюструвати метод Монте-Карло, розглянемо гру під назвою
«Блукаючий п'яниця». Щоб грати в цю гру, потрібна таблиця з сіткою
ліній, як показано на рис. 43.1.

Перейдем к описанию правил игры.

\section{Как играть в «Блуждающего
пьяницу»?}\label{ux43aux430ux43a-ux438ux433ux440ux430ux442ux44c-ux432-ux431ux43bux443ux436ux434ux430ux44eux449ux435ux433ux43e-ux43fux44cux44fux43dux438ux446ux443}

КРОК 1. Блукання п'яниці починаються з довільної точки сітки (у нашому
випадку --- точки А).

\pandocbounded{\includegraphics[keepaspectratio]{images/_page_321_Figure_4.jpeg}}

\begin{itemize}
\tightlist
\item
  Початкова точка - - внутрішні точки сітки - □ - Кінцева точка p;
\item
  \(g_i\) винагороди в m. \(p_I\) Фіг. 43.1. Ігрове поле для гри
  «Блукаючий п'яниця».
\end{itemize}

КРОК 2. У кожному ході гри п'яний випадково переходить до однієї з
чотирьох суміжних точок сітки (у нашому випадку сусідні точки для A ---
це точки B, C, D і E, і ймовірність досягти кожної з цих точок становить
0,25).

Таблица 43.1

{\def\LTcaptype{none} % do not increment counter
\begin{longtable}[]{@{}
  >{\raggedright\arraybackslash}p{(\linewidth - 8\tabcolsep) * \real{0.2523}}
  >{\raggedright\arraybackslash}p{(\linewidth - 8\tabcolsep) * \real{0.4541}}
  >{\raggedright\arraybackslash}p{(\linewidth - 8\tabcolsep) * \real{0.2752}}
  >{\raggedright\arraybackslash}p{(\linewidth - 8\tabcolsep) * \real{0.0092}}
  >{\raggedright\arraybackslash}p{(\linewidth - 8\tabcolsep) * \real{0.0092}}@{}}
\toprule\noalign{}
\begin{minipage}[b]{\linewidth}\raggedright
Вероятность окончания случайного блуждания в граничной точке \(p_i\) и
величина вознаграждения \(g_i\)
\end{minipage} & \begin{minipage}[b]{\linewidth}\raggedright
Граничная точка р і
\end{minipage} & \begin{minipage}[b]{\linewidth}\raggedright
\(P_{\mathcal{A}}\left(p_{i}\right)\) --- относительная доля блужданий,
закончившихся в точке \(p_{i}\)
\end{minipage} & \begin{minipage}[b]{\linewidth}\raggedright
\(g_i\) --- величина вознаграждения при достижении точки \(p_i\)
\end{minipage} & \begin{minipage}[b]{\linewidth}\raggedright
\end{minipage} \\
\midrule\noalign{}
\endhead
\bottomrule\noalign{}
\endlastfoot
1234567891011 & 0,040,150,030,060,170,050,060,150,030,060,160,04 &
110000000 & & \\
\end{longtable}
}

ti

КРОК 3. Після переходу до наступної точки процес блукання відновлюється.
Таким чином, п'яниця ходить від точки до точки, поки випадково не
опиняється на межі \(p_t\) . Тут це зупиняється, і ми записуємо номер
точки \(p_t\) . Це кінець однієї невимушеної прогулянки.

КРОК 4. Повторимо 1-3 рази достатньо і знайдемо відносну кількість
прибуттів п'яниця в кожній з меж \(p_i\) . Таблиця 1. 43.1 показує
типові результати, отримані внаслідок 100 випадкових прогулянок.

КРОК B: Припустимо, що п'яниця отримує винагороду
\(g_i = g\left(p_i\right)\) (значення граничної умови в точці \(p_i\) ,
якщо він досягає точки \(p_i\) ) і припустимо, що мета гри --- обчислити
середній \(R\left(A\right)\) винагороди за всі випадкові прогулянки,
починаючи з точки A. Цей середній виграш визначається формулою

\begin{verbatim}
Игра завершается, когда величина  R\left(A\right)  найдена. В соответствии с данными табл. 43.1 для  R\left(A\right)  получаем
\end{verbatim}

Теперь объясним, зачем нам понадобилось играть в такую игру.

\section{Для чего играть в «Блуждающего
пьяницу»?}\label{ux434ux43bux44f-ux447ux435ux433ux43e-ux438ux433ux440ux430ux442ux44c-ux432-ux431ux43bux443ux436ux434ux430ux44eux449ux435ux433ux43e-ux43fux44cux44fux43dux438ux446ux443}

Виявляється, що середня винагорода, про яку щойно йдеться, є приблизним
розв'язком проблеми Діріхле в точці A. Це цікаве спостереження базується
на двох фактах.

\begin{enumerate}
\def\labelenumi{\arabic{enumi}.}
\item
  Припустимо, що п'яниця почав свій шлях з точки, що лежить на Єрениці.
  Кожна така прогулянка одразу закінчується в одному й тому ж місці, і
  птах отримує винагороду \(g_i\) . Отже, середня винагорода за кожну
  граничну точку дорівнює \(g_i\) .
\item
  Тепер припустимо, що прогулянка починається з внутрішньої точки. Тоді
  стає зрозуміло, що середня винагорода за бал \(R(A)\) буде
  арифметичним середнім середніх нагород за чотири сусідні бали
\end{enumerate}

\[
\begin{cases} u(x, 0) = \sin(\pi x), & 0 \le x \le 1, \\ u_t(x, 0) = 0. & 0 \le x \le 1, \end{cases}
\]

Давайте знову запитаємо себе: чому середнє значення винагороди \(R(A)\)
наближається до розв'язання задачі Діріхле в точці A? Ми бачили, що
значення \(R(A)\) задовольняє двох

уравнениям

(во внутренних точках), \(R(A)\) = \(g_i\) (в граничных точках).

Если \(g_i\) ---это значения функции \(g(x, y)\) из граничного условия в
граничных точках \(p_i\) , то два наших уравнения точно совпадают с
двумя уравнениями, которые были получены при решении задачи Дирихле
методом конечных разностей. То есть величина \(R(A)\) соответствует
величине \(u_{i,f}\) в разностных уравнениях

(i, j)---внутренняя точка, \(u_{ij} = g_{ij},\) \(g_{ij}\) ---значение
решения в граничной точке (i, j).

Отже, значення \(R(A)\) дійсно наближає розв'язок рівняння часткових у
точці A.

Тепер можна вважати, що гра «Блукаючий п'яниця» складається з трьох
кроків.

\section{Розв'язання рівняння Лапласа за методом
Монте-Карло}\label{ux440ux43eux437ux432ux44fux437ux430ux43dux43dux44f-ux440ux456ux432ux43dux44fux43dux43dux44f-ux43bux430ux43fux43bux430ux441ux430-ux437ux430-ux43cux435ux442ux43eux434ux43eux43c-ux43cux43eux43dux442ux435-ux43aux430ux440ux43bux43e}

\begin{itemize}
\tightlist
\item
  Наступні правила дозволяють отримати розв'язок в одній внутрішній
  точці квадрата.
\item
  КРОК 1. Давайте зробимо кілька випадкових прогулянок, починаючи з
  точки A і закінчуючи однією з прикордонних точок. Давайте
  відстежувати, скільки разів прогулянки закінчувалися на кожній межі.
\item
  КРОК 2. Після завершення всіх прогулянок для кожної граничної точки ми
  обчислюємо відносну кількість прогулянок, які закінчилися в цій точці.
  Позначимо ці значення \(P_A(p_i)\) .

  \begin{itemize}
  \tightlist
  \item
    КРОК 3. Обчислюємо наближений розв'язок \(u(A)\) за формулою
    \(u(A) = g_1 P_A(p_1) + g_2 P_A(p_2) + \dots + g_N P_A(p_N)\) ,
  \end{itemize}
\end{itemize}

де \(g_i\) --- значення функції g у точці \(p_i\) , а N --- кількість
граничних точок.

Блукаючий п'яниця можна покращити, щоб розв'язати складнішу задачу.
Нижче наведено приклад такого оновлення.

Розв'язок задачі Діріхле з змінними коефіцієнтами.

Рассмотрим эллиптическую краевую вадачу в квадрате

Дайте физическую интерпретацию этой задачи.

\begin{enumerate}
\def\labelenumi{\arabic{enumi}.}
\setcounter{enumi}{3}
\tightlist
\item
  Для уравнения \(u_{tt} = u_{xx}\)найдите все решения вида
\end{enumerate}

\(u_{xx} + (\sin x) u_{yy} = 0\) , \(0 < x < \pi\) , \(0 < y < \pi\) ,
\((\Gamma Y)\) \(u(x, y) = g(x, y)\) на границе квадрата.

Щоб розв'язати цю задачу, замінимо \(u_{xx}\) , \(u_{yy}\) та \(\sin x\)
наступним чином:

Будет ли решением сумма двух решений такого вида?

\section{ФОРМУЛА ДАЛАМБЕРА МЕТА ЛЕКЦІЇ: Построить решение задачи Коши
для волнового
уравнения}\label{ux444ux43eux440ux43cux443ux43bux430-ux434ux430ux43bux430ux43cux431ux435ux440ux430-ux43cux435ux442ux430-ux43bux435ux43aux446ux456ux457-ux43fux43eux441ux442ux440ux43eux438ux442ux44c-ux440ux435ux448ux435ux43dux438ux435-ux437ux430ux434ux430ux447ux438-ux43aux43eux448ux438-ux434ux43bux44f-ux432ux43eux43bux43dux43eux432ux43eux433ux43e-ux443ux440ux430ux432ux43dux435ux43dux438ux44f}

и \(_{yy}=u_{l+1,\ j}-2u_{l,\ j}+u_{l-1,\ j},\) (центральные разностные
производные),

і підставляють їх у рівняння похідних похідних. Відповідна сітка
показана на рис. 43.2. Розв'язавши отримане рівняння за \(u_{i,l}\) ,
отримуємо

Посмотрим внимательно на это уравнение. Коэффициенты при
\(u_{i+1,\ j}\),~\(u_i\),~j+i,~\(u_l\),~j-1 н \(u_{i-1}\),~j
положительны, и их сумма равна единице. Другими словами, решение
\(u_i\),~j является взвешенным средним решений в четырех соседних
точках. Следовательно, можно Р И С

\pandocbounded{\includegraphics[keepaspectratio]{images/_page_324_Picture_9.jpeg}}

. 43.2. Узлы сетки для случайного маршруга.

модифицировать «Блуждающего пьяницу» так, чтобы вероятности перехода в
соседние точки были не 0,25, а равнялись бы коэффициенту в
соответствующем члене. Другими словами, если пьяница находится в точке
(i, j), то он переходит в точку

Ця задача описує рух нескінченної струни за заданих початкових умов.
Його розв'язав у 1750 році французький математик д'Аламбер. Формула
Д'Аламбера

\begin{verbatim}
с вероятностью  \frac{1}{2(1+\sin x_j)} ,  (i, j-1)  с вероятностью  \frac{1}{2(1+\sin x_j)} ,  (i+1, j)  с вероятностью  \frac{\sin x_j}{2(1+\sin x_j)} ,  (i-1, j)  с вероятностью  \frac{\sin x_j}{2(1+\sin x_j)} .
\end{verbatim}

В остальном игра не меняется. Модификация для других задач может
оказаться более хитроумной, но основная идея остается той же. Читатель
может заняться разработкой своих игр для решения других задач.
Параболический случай рассмотрен в разделе «Задачи».

\section{ЗАУВАЖЕННЯ}\label{ux437ux430ux443ux432ux430ux436ux435ux43dux43dux44f-7}

\begin{verbatim}
1. Отметим, что если относительные частоты попадания  P_A(p_l)  известны, то можно легко найти величину  u\left(A\right)  при других граничных условиях  g_l . Для этого достаточно подставить  g_l  в формулу
\end{verbatim}

полегшує пошук розв'язку, якщо початкові умови відомі. Крім того, це
дозволяє дати цікаву фізичну інтерпретацію розв'язку мовою хвиль, що
рухаються.

Як читач, ймовірно, пам'ятає, у параболічному випадку ми спочатку
розв'язали задачі дифузії на скінченному відрізку (методом розділення
змінних), а лише потім переходили до необмежених інтервалів (-
\(\infty < x < \infty\) ), де застосовували перетворення Фур'є. У
гіперболічному випадку ми зробимо навпаки. Почнемо з розв'язання
одномірного хвильового рівняння на всій дійсній осі, тобто з задачі Коші

(17.1)

а не разыгрывать вновь случайные прогулки.

\begin{enumerate}
\def\labelenumi{\arabic{enumi}.}
\setcounter{enumi}{1}
\tightlist
\item
  Во многих случаях исследователя интересует решение уравнения с
  частными производными в одной точке. Если граница области сложна и
  уравнение зависит от трех или четырех независимых переменных, метод
  Монте-Карло может оказаться удобным. Методы Монте-Карло возникли при
  решении очень трудных задач диффузии нейтронов. Эти задачи невозможно
  решать аналитически.
\end{enumerate}

\section{ЗАДАЧІ}\label{ux437ux430ux434ux430ux447ux456}

\begin{enumerate}
\def\labelenumi{\arabic{enumi}.}
\tightlist
\item
  Построить блок-схему игры «Блуждающий пьяница» для решения задачи
\end{enumerate}

(УЧП)

, \(0 < x < 1\) , \(0 < y < 1\) , (ГУ) \(u(x, y) = g(x, y)\) на границе,

\begin{itemize}
\tightlist
\item
  у вузлах внутрішньої сітки, якщо кількість вузлів сітки довільна. 2.
  Напишіть програму для розв'язання задачі відповідно до блок-схеми з
  завдання 1.
\item
  \begin{enumerate}
  \def\labelenumi{\arabic{enumi}.}
  \setcounter{enumi}{2}
  \tightlist
  \item
    Як модифікувати гру для вирішення проблеми
  \end{enumerate}
\end{itemize}

(УЧП)

содержащей только начальные условия. Эту задачу можно было бы решить
либо с помощью преобразования Фурье (по переменной t), либо с помощью
преобразования Лапласа (по переменной t). Однако мы воспользуемся другим
методом (методом канонических координат), который, как мы надеемся,
заинтересует читателя. Этот метод базируется на тех же идеях, что и
метод перехода к движущейся системе отсчета, с которым мы познакомились
в лекции 15. Итак, приступим к решению задачи (17.1).

\section{Решение одномерного волнового уравнения. Формула
Даламбера}\label{ux440ux435ux448ux435ux43dux438ux435-ux43eux434ux43dux43eux43cux435ux440ux43dux43eux433ux43e-ux432ux43eux43bux43dux43eux432ux43eux433ux43e-ux443ux440ux430ux432ux43dux435ux43dux438ux44f.-ux444ux43eux440ux43cux443ux43bux430-ux434ux430ux43bux430ux43cux431ux435ux440ux430}

Решение задачи (17.1) разобьем на несколько шагов:

\begin{verbatim}
ШАГ 1 (Замена координат (x, t) новыми каноническими координатами   (\xi, \eta)  ).
\end{verbatim}

Для решения задачи (17.1) воспользуемся тем, что если заменить две
независимые переменные x и t новыми пространственновременными
координатами

, \(0 < x < 1\) , \(0 < y < 1\) , (ГУ) \(u(x, y) = g(x, y)\) на границе?

Как выглядит в этой игре случайная прогулка?

\begin{enumerate}
\def\labelenumi{\arabic{enumi}.}
\setcounter{enumi}{3}
\tightlist
\item
  Можете ли вы так модифицировать игру, чтобы решить задачу
\end{enumerate}

\[
\begin{array}{l} u_{xx} + u_{yy} = 0, \quad 0 < x < 1, \quad 0 < y < 1, \\ u(x, 1) = 0, \quad 0 < x < 1, \quad 0 < y < 1, \\ u(x, 0) = 0, \quad 0 < x < 1, \quad 0 < y < 1, \\ u(0, y) = 1, \quad \frac{\partial u}{\partial x}(1, y) = 0 \end{array}
\]

то уравнение

\begin{verbatim}
\begin{aligned}
(\text{УЧП})\quad & u_t = \alpha^2 u_{xx} - \beta u_x - \gamma u, && 0 < x < 1,\ 0 < t < \infty, \\
(\text{ГУ})\quad & u(0,t) = f(t),\quad u(1,t) = g(t), && 0 < t < \infty, \\
(\text{НУ})\quad & u(x,0) = \varphi(x), && 0 \leqslant x \leqslant 1.
\end{aligned}
\end{verbatim}

преобразуется в уравнение

ПРИМІТКА. Замінимо це рівняння на скінченну різницю Кренка--Ніколсона,
виразимо \(u_{i+1,j}\) через п'ять значень \(u_{i+1,j-1}\) ,
\(u_{i+1,j+1}\) , \(u_{i,j+1}\) , \(u_{i,j+1}\) , \(u_{i,j+1}\) у
суміжних точках.

\section{КАЛЬКУЛЮС ВАРІАЦІЙ (РІВНЯННЯ
ЕЙЛЕРА-ЛАГРАНЖА)}\label{ux43aux430ux43bux44cux43aux443ux43bux44eux441-ux432ux430ux440ux456ux430ux446ux456ux439-ux440ux456ux432ux43dux44fux43dux43dux44f-ux435ux439ux43bux435ux440ux430-ux43bux430ux433ux440ux430ux43dux436ux430}

МЕТА ЛЕКЦІЇ: Ввести поняття функціоналу, тобто функції від функції, і
пояснити, як функціонали природно виникають у фізиці. Типовим
функціоналом є інтеграл

\[
J[y] = \int_a^b F(x, y, y')\,dx,
\]

який є функцією від функції \(y\) (при цьому вважається, що
підінтегральна функція \(F(x, y, y')\) задана). Приклад функціоналу:

\[
J[y] = \int_0^1 \left[y^2(x) + y'^2(x)\right] dx.
\]

Ми покажемо, як знайти функцію \(\overline{y}(x)\), що мінімізує
функціонал \(J[y]\). Виявляється, що мінімізуюча функція
\(\overline{y}\) повинна задовольняти так зване рівняння
Ейлера-Лагранжа. Воно відіграє ту саму роль, що й необхідна умова
мінімуму

\[
\frac{df(x)}{dx} = 0
\]

для функції \(f(x)\) у точці \(x\) в диференціальному численні.

Варіаційне числення тісно пов'язане з диференціальними рівняннями, але,
на жаль, більшість студентів його не вивчає. У цій лекції (і наступній)
подається вступ до варіаційного числення та показано, як можна
розв'язувати диференціальні рівняння в частинних похідних на основі
варіаційних принципів.

Варіаційне числення виникло одночасно з математичним аналізом у зв'язку
з розв'язанням задач максимізації та мінімізації функцій від функцій,
тобто функціоналів. Першою задачею варіаційного числення була задача про
брахістохрону, сформульована Йоганном Бернуллі у 1696 році. У цій задачі
потрібно знайти криву \(y(x)\) так, щоб мінімізувати час спуску без
тертя з однієї точки в іншу (див. рис. 44.1).

Бернуллі показав, що час спуску записується у вигляді

\[
T = \int_0^T dt = \int_a^L \frac{dt}{ds}\,ds = \int_0^L \frac{ds}{v}
= \frac{1}{\sqrt{2mg}} \int_0^L \frac{ds}{\sqrt{y}}
= \frac{1}{\sqrt{2mg}} \int_a^b \sqrt{\frac{1+y'^2}{y}}\,dy
\]

і, отже, є функціоналом \(T\{y\}\) від функції. Оскільки багато
функціоналів мають подібну будову, ми зосередимося на вивченні
функціоналу загального вигляду

(44.1) \[
J[y] = \int_a^b F(x, y, y')\,dx.
\]

Тепер сформулюємо головну мету лекції: знайти функцію, яка доставляє
мінімум (або максимум) функціоналу (44.1). Стратегія пошуку буде такою
ж, як і при пошуку мінімуму функції у звичайному диференціальному
численні. Там ми шукали критичні точки з умови \(f'(x)=0\). У
варіаційному численні все дещо складніше, бо аргументом є не число, а
функція. Проте загальний підхід той самий: ми обчислимо функціональну
похідну за функцією \(y(x)\) і прирівняємо її до нуля. Нове рівняння
буде аналогом умови \(df(x)/dx = 0\), але тепер це буде звичайне
диференціальне рівняння, відоме як рівняння Ейлера-Лагранжа.

Мінімізація функціоналу

\[
J[y] = \int_a^b F(x, y, y')\,dx
\]

Розглянемо задачу знаходження функції \(y(x)\), яка мінімізує цей
функціонал у класі гладких функцій, що задовольняють граничним умовам

\[
y(a) = A, \qquad y(b) = B.
\]

(див. рис. 44.2).

Нехай шукана функція \(y = \overline{y}(x)\) існує, і розглянемо малу
варіацію цієї функції, тобто функцію
\(\overline{y} + \varepsilon\eta(x)\), де \(\varepsilon\) --- мале
число, а \(\eta(x)\) --- гладка функція, яка задовольняє граничним
умовам \(\eta(a)=\eta(b)=0\).

Тоді, якщо обчислити інтеграл \(J\) для близької функції
\(\overline{y} + \varepsilon\eta\), функціонал зросте, тобто

\[
J[\overline{y}] \leqslant J[\overline{y} + \varepsilon\eta]
\]

для всіх \(\varepsilon\). Іншими словами, графік
\(\varphi(\varepsilon)=J[\overline{y}+\varepsilon\eta]\) як функції
\(\varepsilon\) має вигляд, показаний на рис. 44.3.

\pandocbounded{\includegraphics[keepaspectratio]{images/_page_329_Figure_5.jpeg}}

Рис. 44.1. Задача про брахістохрону (з цієї задачі почалося варіаційне
числення).

\pandocbounded{\includegraphics[keepaspectratio]{images/_page_329_Figure_7.jpeg}}

Рис. 44.2. Варіація функції: a --- варіація функції
\(\overline{y}(x)+\varepsilon\eta(x)\); б --- мінімізуюча функція
\(\overline{y}(x)\).

З рис. 44.3 видно, що тут потрібно обчислити похідну від

\[
\varphi(\varepsilon) = J[\overline{y} + \varepsilon\eta]
\]

за \(\varepsilon\), покласти \(\varepsilon = 0\) і прирівняти отриманий
вираз до нуля, тобто

\[
\frac{d\varphi(\varepsilon)}{d\varepsilon}
= \frac{d}{d\varepsilon} J[\overline{y} + \varepsilon\eta]\bigg|_{\varepsilon=0}
= \int_a^b \left[
\frac{\partial F}{\partial \overline{y}}\eta(x)
+ \frac{\partial F}{\partial \overline{y'}}\eta'(x)
\right] dx = 0.
\]

(Читач має виконати це обчислення самостійно.) Інтегруючи частинами,
дістаємо

\[
\frac{d\varphi(\varepsilon)}{d\varepsilon}
\equiv \int_a^b \left\{
\frac{\partial F}{\partial \overline{y}} - \frac{d}{dx}\left[\frac{\partial F}{\partial \overline{y'}}\right]
\right\} \eta(x)\,dx = 0.
\]

Оскільки цей інтеграл дорівнює нулю для будь-якої функції \(\eta(x)\),
яка задовольняє граничним умовам \(\eta(a)=\eta(b)=0\), то
підінтегральний вираз має дорівнювати нулю, тобто

\pandocbounded{\includegraphics[keepaspectratio]{images/_page_330_Figure_6.jpeg}}

Рис. 44.3. Графік функції \(J[y+\varepsilon\eta]\) в околі
\(\varepsilon=0\).

(44.2) \[
\frac{\partial F}{\partial \overline{y}} - \frac{d}{dx}\left[\frac{\partial F}{\partial \overline{y'}}\right] = 0
\]

(рівняння Ейлера-Лагранжа).

Рівняння (44.2) називається рівнянням Ейлера-Лагранжа. Хоча в загальному
вигляді воно здається складним, після підстановки конкретної функції
\(F(x,y,y')\) воно перетворюється на звичайне диференціальне рівняння
другого порядку відносно невідомої функції \(\overline{y}(x)\). Отже,
щоб визначити мінімізуючу функцію \(\overline{y}\), треба розв'язати
рівняння Ейлера-Лагранжа.

Отже, ми показали, що якщо функція \(y(x)\) мінімізує функціонал
\(J[y] = \int_a^b F(x,y,y')\,dx\) у класі гладких функцій із граничними
умовами \(y(a)=A\) і \(y(b)=B\), то вона повинна задовольняти рівнянню

\[
\frac{\partial F}{\partial y} - \frac{d}{dx}\left[\frac{\partial F}{\partial y'}\right] = 0.
\]

(Для спрощення позначень ми прибрали риску над \(\overline{y}\).) Для
ілюстрації теорії розглянемо приклад.

Пошук мінімальної функціональності
\(J[y] = \int_{0}^{1} [y^2 + y'^2] dx\) .

Спробуємо знайти функцію \(y(x)\), яка проходить через точки (0,0) і
(1,1) і мінімізує функціональність

\[
P_{0}
\]

З контексту зрозуміло, що шукана функція має бути диференційовною,
оскільки інтегративний вираз залежить від y'. Щоб визначити
\(\bar{y}(x)\) , запишемо рівняння Ейлера--Лагранжа (разом із граничною
умовою \(y(0)\)=0 та \(y(1)\)=1).

\begin{verbatim}
,  (x,  y, z) \in R^3  ,  u = 0   u_t =
\end{verbatim}

\(y(0) = 0,\) \(y(1) = 1.\) Але \(F(x, y, y') = y^2 + y'^2\) і, отже,

\begin{enumerate}
\def\labelenumi{\arabic{enumi}.}
\setcounter{enumi}{3}
\tightlist
\item
  Решите аналогичную (см. задачу 3) задачу для двумерного уравнения
\end{enumerate}

(УЧП)

(Продифференцировали \(F\) по \(y\) и \(y'.\) )

Уравнение Эйлера --- Лагранжа принимает вид

Зверніть увагу, що в перших двох задачах коефіцієнти є сталими. Читач
повинен розуміти, що параметр b має бути достатньо малим, а необережний
ряд може відрізнятися

\section{ЗАДАЧИ}\label{ux437ux430ux434ux430ux447ux438}

\begin{verbatim}
- 1. Подставьте разложение (46.8) в задачу (46.7) и получите последовательность задач P_0,\ P_1,\ P_2,\ \dots . 2. Покажите, что нелинейную задачу
\end{verbatim}

откуда

(НУ)

Решая это простое дифференциальное уравнение с граничными условиями
\(y(0)\)=0 и \(y(1)\)=1, получаем

\[
u(r, \theta) = u_0(r, \theta) + \frac{1}{4}u_1(r, \theta)
\]

.

График этой функции изображен на рис. 44.4. Ясно, что любая другая
гладкая функция, удовлетворяющая тем же граничным условиям, даст большее
значение функционала \(J(y)\).

\pandocbounded{\includegraphics[keepaspectratio]{images/_page_332_Figure_5.jpeg}}

Рис. 44 4. Все допустимые гладкие кривые удовлетворяют граничным
условиям \(y(0)\) = 0 и \(y(1)\) = 1: a---допустимая кривая; \(\delta\)
---минимизирующая функция \(y(x) = 0.42e^x - 0.42e^{-x}\) . \#
ЗАУВАЖЕННЯ 1. Уравнение Эйлера---Лагранжа аналогично равенству нулю
производной в дифференциальном исчислении. Если читатель помнит, то на
самом деле равенство нулю производной не является достаточным условием
экстремума. Например, производная функции \(f(x) = x^3\) при x = 0 равна
нулю, но эта точка не является нп точкой минимума, ни точкой максимума.
Так же и с уравнением Эйлера---Лагранжа: оно является только необходимым
условием экстремума, но не достаточным. Значит, экстремум может
достигаться на других функциях. Очень часто, однако, решение уравнения
Эйлера---Лагранжа по существу задачи дает локальный или даже глобальный
минимум.

\begin{enumerate}
\def\labelenumi{\arabic{enumi}.}
\setcounter{enumi}{1}
\tightlist
\item
  Если у (х) минимизирует функционал
\end{enumerate}

\[
\Delta u = 0,
\]

то площадь поверхности, образованной вращением кривой \(y(x)\) вокруг
оси x, будет минимальной (рис. 44.5). Решением уравнения Эйлера ---
Лагранжа в этом случае будет

цепная линия

\[
q_{xx} + \varphi_{uu} = 0
\]

где константы \(\alpha\) и \(\beta\) определяются из условий \(y(a)\) =
A и \(y(b)\) = B (т. е. цепная линия должна проходить через заданные
точки). Решить уравнение Эйлера---Лагранжа в данном случае достаточно
трудно, поэтому мы рекомендуем читателю только проверить, что цепная
линия ему удовлетворяет.

\begin{enumerate}
\def\labelenumi{\arabic{enumi}.}
\setcounter{enumi}{2}
\tightlist
\item
  В этой лекции мы нашли, что функция
\end{enumerate}

\[
q_{uu} + q_{vv} = 0
\]

доставляет минимум функционалу

\pandocbounded{\includegraphics[keepaspectratio]{images/_page_353_Figure_9.jpeg}}

При этом оказывается, что \(J[\overline{y}] = 0.46\) . Если подставить в
функционал любую другую гладкую функцию, график которой проходит через
точки (0, 0) и (1, 1), то значение функционала J{[}y{]} станет больше.

\pandocbounded{\includegraphics[keepaspectratio]{images/_page_333_Picture_8.jpeg}}

Рис. 44 5. Минимальная поверхность вращения.

\begin{enumerate}
\def\labelenumi{\arabic{enumi}.}
\setcounter{enumi}{3}
\tightlist
\item
  Основные законы физики чаще всего формулируются на языке вариационных
  принципов, а не дифференциальных уравнений. В качестве примеров можно
  привести принцип Ферма (свет при распространении из одной точки в
  другую выбирает путь, которому соответствует наименьшее время
  распространения) или принцип Гамильтона (в консервативном поле частица
  движется так, что интеграл действия
\end{enumerate}

\[
w = f(z)
\]

(кинетическая энергия--- потенциальная энергия) \(dt\)

будет минимальным). Таким образом, физические явления развиваются только
так, что эти функционалы принимают минимальные значения.

Основные идеи вариационного исчисления можно распространить на
функционалы, зависящие от функции нескольких независимых переменных,
вида

\pandocbounded{\includegraphics[keepaspectratio]{images/_page_355_Figure_5.jpeg}}

Соответствующее уравнение Эйлера --- Лагранжа для такого функционала
будет уже уравнением с частными производными

\pandocbounded{\includegraphics[keepaspectratio]{images/_page_355_Figure_7.jpeg}}

Такого рода функционалы мы рассмотрим в следующей лекции. Однако общая
философия метода будет иной. Некоторым новым способом мы найдем функцию
\(u(x, y)\), которая минимизирует функционал, а поскольку она является
решением уравнения Эйлера --- Лагранжа, то тем самым мы найдем решение
уравнения с частными производными. Другими словами, мы будем искать
решение уравнения с частными производными методом минимизации
функционала. В этой лекции все было наоборот: задачу минимизации
функционала, зависящего от функции одной переменной, мы сводили к
решению уравнения Эйлера --- Лагранжа. Методы решения дифференциальных
уравнений путем минимизации соответствующих функционалов принято
называть прямыми методами вариационного исчисления. Методы минимизации
функционалов путем решения соответствующих уравнений Эйлера --- Лагранжа
называются непрямыми методами вариационного исчисления. Следующая лекция
посвящена хорошо известному прямому методу Ритца. \# ЗАДАЧІ 1. Среди
всех кривых, удовлетворяющих условиям \(y(0)\) = 0, и \(y(1)\) = 1,
найдите ту, которая минимизирует функционал

\[
u + iv = (x + iy)^2 = x^2 - y^2 + 2ixy
\]

Как вы интерпретируете получелный результат? Чему равна величина
\(J[\overline{y}]\) ? Каков смысл величины \(J[\overline{y}]\) ?

\begin{enumerate}
\def\labelenumi{\arabic{enumi}.}
\setcounter{enumi}{1}
\tightlist
\item
  Кинетическая энергия колеблющейся материальной точки определяется
  выражением \(KE = \frac{1}{2} m\dot{y}^2\) , где y = dy/dt, а
  потенциальная равна \(PE = \frac{1}{2} ky^2\) .
\end{enumerate}

\pandocbounded{\includegraphics[keepaspectratio]{images/_page_335_Picture_2.jpeg}}

Принцип Гамильтона утверждает, что материальная точка будет двигаться
так, что интеграл

\begin{verbatim}
\left\{
\end{verbatim}

будет минимален. Если этот закоп справедлив, то получите
дифференциальное уравнение движения материальной точки.

\begin{enumerate}
\def\labelenumi{\arabic{enumi}.}
\setcounter{enumi}{2}
\tightlist
\item
  Покажите, что в классе гладких функций, удовлетворяющих граничным
  условиям \(y(0)\)=0 и \(y(\pi/2)=1\) , минимум функционала
\end{enumerate}

Чи знаєте ви, як виглядає графік цієї функції з сегментом {[}0, 1{]}?
Подумай про це. Зверніть увагу, що синусоїдальне перетворення функції
\(f(x)\) визначається лише для додатних цілих чисел n.~Іншими словами,
скінченний синус і косинус перетворюють функції у чисельні
послідовності.

\pandocbounded{\includegraphics[keepaspectratio]{images/_page_185_Figure_3.jpeg}}

РИС. 25.1. Графики функции \(f(x)\) = 1 и ее преобразования.

\section{Свойства
преобразований}\label{ux441ux432ux43eux439ux441ux442ux432ux430-ux43fux440ux435ux43eux431ux440ux430ux437ux43eux432ux430ux43dux438ux439}

Прежде чем приступить к решению задач, мы должны получить некоторые
свойства этих преобразований.

\begin{verbatim}
Якщо you(x, t) — функція  \partial \theta yx  змінних, і ми виконуємо перетворення змінної x, то
\end{verbatim}

достигается на функции \(\overline{y}(x) = \sin x\) . Вычислите
\(J(\sin x)\) . 4. Исходя из функционала

(Зверніть увагу, що перетворення було виконано на змінній x, і отримана
послідовність залежить лише від часу t.)

А как быть с производными? Ниже приводятся несколько полезных формул

получите уравнение Эйлера --- Лагранжа

\section{ВАРИАЦИОННЫЕ МЕТОДЫ РЕШЕНИЯ УРАВНЕНИЙ С ЧАСТНЫМИ ПРОИЗВОДНЫМИ
МЕТА ЛЕКЦІЇ: Показать, как можно решать дифференциальное уравнение,
рассматривая его как уравнение Эйлера---Лагранжа для некоторого
функционала и находя некоторым новым методом функцию, минимизирующую
этот функционал. Тогда минимизирующая функция будет решением уравнения с
частными производными. Задача, конечно, состоит в том, чтобы найти такой
функционал, для которого исходное уравнение является уравнением
Эйлера---Лагранжа. Хорошо известен результат (теорема о минимуме
энергии), который гласит, что нахождение решения и некоторой
эллиптической краевой
задачи}\label{ux432ux430ux440ux438ux430ux446ux438ux43eux43dux43dux44bux435-ux43cux435ux442ux43eux434ux44b-ux440ux435ux448ux435ux43dux438ux44f-ux443ux440ux430ux432ux43dux435ux43dux438ux439-ux441-ux447ux430ux441ux442ux43dux44bux43cux438-ux43fux440ux43eux438ux437ux432ux43eux434ux43dux44bux43cux438-ux43cux435ux442ux430-ux43bux435ux43aux446ux456ux457-ux43fux43eux43aux430ux437ux430ux442ux44c-ux43aux430ux43a-ux43cux43eux436ux43dux43e-ux440ux435ux448ux430ux442ux44c-ux434ux438ux444ux444ux435ux440ux435ux43dux446ux438ux430ux43bux44cux43dux43eux435-ux443ux440ux430ux432ux43dux435ux43dux438ux435-ux440ux430ux441ux441ux43cux430ux442ux440ux438ux432ux430ux44f-ux435ux433ux43e-ux43aux430ux43a-ux443ux440ux430ux432ux43dux435ux43dux438ux435-ux44dux439ux43bux435ux440ux430ux43bux430ux433ux440ux430ux43dux436ux430-ux434ux43bux44f-ux43dux435ux43aux43eux442ux43eux440ux43eux433ux43e-ux444ux443ux43dux43aux446ux438ux43eux43dux430ux43bux430-ux438-ux43dux430ux445ux43eux434ux44f-ux43dux435ux43aux43eux442ux43eux440ux44bux43c-ux43dux43eux432ux44bux43c-ux43cux435ux442ux43eux434ux43eux43c-ux444ux443ux43dux43aux446ux438ux44e-ux43cux438ux43dux438ux43cux438ux437ux438ux440ux443ux44eux449ux443ux44e-ux44dux442ux43eux442-ux444ux443ux43dux43aux446ux438ux43eux43dux430ux43b.-ux442ux43eux433ux434ux430-ux43cux438ux43dux438ux43cux438ux437ux438ux440ux443ux44eux449ux430ux44f-ux444ux443ux43dux43aux446ux438ux44f-ux431ux443ux434ux435ux442-ux440ux435ux448ux435ux43dux438ux435ux43c-ux443ux440ux430ux432ux43dux435ux43dux438ux44f-ux441-ux447ux430ux441ux442ux43dux44bux43cux438-ux43fux440ux43eux438ux437ux432ux43eux434ux43dux44bux43cux438.-ux437ux430ux434ux430ux447ux430-ux43aux43eux43dux435ux447ux43dux43e-ux441ux43eux441ux442ux43eux438ux442-ux432-ux442ux43eux43c-ux447ux442ux43eux431ux44b-ux43dux430ux439ux442ux438-ux442ux430ux43aux43eux439-ux444ux443ux43dux43aux446ux438ux43eux43dux430ux43b-ux434ux43bux44f-ux43aux43eux442ux43eux440ux43eux433ux43e-ux438ux441ux445ux43eux434ux43dux43eux435-ux443ux440ux430ux432ux43dux435ux43dux438ux435-ux44fux432ux43bux44fux435ux442ux441ux44f-ux443ux440ux430ux432ux43dux435ux43dux438ux435ux43c-ux44dux439ux43bux435ux440ux430ux43bux430ux433ux440ux430ux43dux436ux430.-ux445ux43eux440ux43eux448ux43e-ux438ux437ux432ux435ux441ux442ux435ux43d-ux440ux435ux437ux443ux43bux44cux442ux430ux442-ux442ux435ux43eux440ux435ux43cux430-ux43e-ux43cux438ux43dux438ux43cux443ux43cux435-ux44dux43dux435ux440ux433ux438ux438-ux43aux43eux442ux43eux440ux44bux439-ux433ux43bux430ux441ux438ux442-ux447ux442ux43e-ux43dux430ux445ux43eux436ux434ux435ux43dux438ux435-ux440ux435ux448ux435ux43dux438ux44f-ux438-ux43dux435ux43aux43eux442ux43eux440ux43eux439-ux44dux43bux43bux438ux43fux442ux438ux447ux435ux441ux43aux43eux439-ux43aux440ux430ux435ux432ux43eux439-ux437ux430ux434ux430ux447ux438}

\[
w = \ln\left\{\frac{z-1}{z+1}\right\}
\]

, в области D , u = 0 , на границе D ,

эквивалентно нахождению такой функции u (тоже обращающейся в нуль на
границе области D), которая минимизирует функционал потенинальной
энергии

То есть, \(\Delta u = f\) ---уравнение Эйлера---Лагранжа для J{[}u{]}.
Функция, приближенно минимизирующая функционал J{[}u{]}, ищется методом
Ритца; тем самым мы получаем решение (приближениое) уравнения с частными
производными. В этой лекции мы познакомимся с методом Ритца и посмотрим,
как с помощью этого метода можно и и и и и и и уравнения.

Очень хорошо решать граничные задачи (как, например, задачу о
закрепленной мембране) путем отыскания гладкой поверхности, которая
минимизирует потенциальную энергию мембраны. То есть, если мы
рассматриваем исходное уравнение как уравнение Эйлера --- Лагранжа для
искоторого функционала J{[}u{]}, то решение дифференциального уравнения
можно получить путем минимизации функционала (поскольку минимизирующая
функция функционала является решением соответствующего уравнения

Эйлера --- Лагранжа). В лекции 44 мы рассматривали такие функционалы, у
которых уравнение Эйлера --- Лагранжа было обыкновенным дифференциальным
уравнением. В этой лекции рассматриваются функционалы, для которых
уравнение Эйлера --- Лагранжа относится к классу уравнений с частными
производными. Например для функционала

уравнение Эйлера---Лагранжа (подобно случаю обыкновенного
дифференциального уравнения в последней лекции) имеет вид

и, следовательно, для решения задачи Дирихле в единичном квадрате

(УЧП)

\[
(\text{УЧП}) \quad u_{tt} = u_{xx} + \sin(\pi x), \quad 0 < x < 1, \quad 0 < t < \infty,
\]

, \(0 < x < 1\) , \(0 < y < 1\) , (ГУ) \(u = g\) , на границе квадрата,

можно, наоборот, найти такую функцию \(u(x, y)\), которая минимизирует
J{[}u{]} и равна g на границе. Вероятно, не стоит удивляться тому, что
функционал

\[
(\text{ГУ}) \quad
\]

представляет собой потенциальную эпергию мембраны и фактически мы
находим поверхность с минимальной потенциальной энерешей. Вопрос,
конечно, состоит в следующем: если дано дифференциальное уравнение, как
мы находим функционал, выражающий потенциальную энергию решения? Ответ
на этот вопрос дает хорошо известная теорема (теорема о минимуме
энергии), которая утверждает:

Решение и задачи Дирихле

\begin{verbatim}
^{\Lambda}u = f   в области   ^{D}     (ГУ)     ^{U}=0   на границе   ^{D}
\end{verbatim}

является той же функциен, которая минимизирует (среди функций,
удовлетворяющих граничному условию u=0) энергегический функционал

\begin{verbatim}
\begin{cases} u_{tx}, \quad 0 = \begin{cases} 1, \quad -1 < x < 1, \\ 0 \quad \text{в остальных точках,} \end{cases}
\end{verbatim}

\[
= \ln\left[\frac{z-1}{z+1}\right] + i \arg\left[\frac{z-1}{z+1}\right]
\]

и, следовательно, для решения задачи (45.1) достаточно найти функцию
\(\overline{u}\) (среди функций, обращающихся в нуль на граннце),
минимизирующую \(J\{u\}\) . Этим замечанием завершается первая часть
лекции. В оставшейся части мы покажем, как находить минимизирующую
функцию \(\overline{u}\) для J{[}u{]}, используя метод Ритца.

\section{Метод Ритца минимизации
функционалов}\label{ux43cux435ux442ux43eux434-ux440ux438ux442ux446ux430-ux43cux438ux43dux438ux43cux438ux437ux430ux446ux438ux438-ux444ux443ux43dux43aux446ux438ux43eux43dux430ux43bux43eux432}

Этот метод --- один из многих вариационных методов, описанных в
литературе по вариационному исчислению. Идея, предложенная математиком
В. Ритцем, очень проста. Метод состоит из следующих шагов.

ШАГ 1. Выбираем n и функцию u, минимизирующую функционал

\[
= \arg\left[\frac{x^2+y^2-1+i2y}{(x+1)^2+y^2}\right] =
\]

ищем в виде

\[
\begin{cases} u(x, 0) = 1, \\ u_t(x, 0) = 0, \end{cases}
\]

где функции \(\phi_1, \phi_2, \phi_3, \ldots, \phi_n\) принадлежат к
классу достаточно хороших функций и все обращаются в нуль на границе,
так что из них можно построить разумное приближение для решения задачи.
Эти функции принято называть пробными функциями. Типичный набор пробных
функций для задачи Дирихле в единичном квадрате приведен ниже:

\[
x^{2} + y^{2} = 1,
\]

--- обращаются в нуль на границе,
\(\varphi_2(x, y) = x\varphi_1(x, y),\)
\(\varphi_3(x, y) = y\varphi_1(x, y),\)
\(\varphi_4(x, y) = x^2\varphi_1(x, y),\)
\(\varphi_5(x, y) = xy\varphi_1(x, y),\)
\(\varphi_6(x, y) = y^2\varphi_1(x, y),\) \(\vdots\) \(\vdots\)

Иными словами, первые четыре приближения имеют вид

Решение задачи выполним за три следующих шага:

ШАГ І (Определение подходящего преобразования.)

Поскольку переменная x меняется от 0 до 1, мы воспользуемся конечным
преобразованием, а именно синус-преобразованием; позже станет ясно
почему. Мы могли бы решить эту задачу с помощью преобразования Лапласа
по переменной t (по трудоемкости этот метод не отличается от конечного
синус-преобразования).

ШАГ 2 (Выполнение преобразования).

Для удобства будем использовать обозначение \(S_n\)(t) = S{[}u{]} .
Применим синус-преобразование к исходному уравнению

\begin{verbatim}
\begin{aligned} & u_{ti} = c^{2}u_{xx}, & -\infty < x < +\infty, & 0 < t < \infty, \\ & (\text{НУ}) & \begin{cases} u_{t}(x, 0) = 0, \\ u_{t}(x, 0) = \begin{cases} 1, & -1 < x < 1, \\ 0, & \text{в остальных случаях.} \end{cases} \end{aligned}
\end{verbatim}

и воспользуемся тождеством для синус-преобразования, получим

в координатах (x, t).

\pandocbounded{\includegraphics[keepaspectratio]{images/_page_134_Figure_7.jpeg}}

Рис. 18.4. Решение задачи (18.4) в xt-плоскости.

Задача (18.4) описывает поведение струны, которой сообщается рачальная
единичная скорость на отрезке -1 \textless{} x \textless{} 1. Смещение .
.

описывается формулой Даламбера

\[
\begin{aligned} & (\text{УЧП}) & \phi_{xx} + \phi_{yy} = 0, & \text{внутри } D, \\ & (\text{ГУ}) & \begin{cases} \varphi(x, y) = 1 & \text{на } x^2 + y^2 = 1, \\ \varphi(x, y) = 2 & \text{на } (x - 1)^2 + y^2 = 9. \end{cases} \end{aligned}
\]

является теперь функцией коэффициентов \(a_1a_2\dots a_n\) .
Следовательно, для того чтобы найти минимум функционала J, приравняем
нулю частные производные

\[
Aa = b
b_i = -\int_0^1 \int_0^1 f(x, y) \varphi_i(x, y) dx dy,
\]

где

Это решение в различные моменты времени изображено на рис. 18.5.

\pandocbounded{\includegraphics[keepaspectratio]{images/_page_135_Figure_5.jpeg}}

Рис. 18.5. Решение задачи (18.4) в различные моменты времени.

Этим завершается наша интерпретация формулы Даламбера в плоскости
переменных \(\boldsymbol{x}\) и t. В оставшейся части лекции мы займемся
решением смешанной задачи для полубесконечной струны:

(18.6)

\[
u = \gamma [(x-s)(x-t)-\gamma^2],
\]

Единственное отличие в том, что аппроксимирующие функции
\(\phi_1, \phi_2, \phi_3, \ldots, \phi_n\) должны удовлетворять
граничным условиям:

\[
\gamma = 2t/[(x-t)^2 + y^2].
\]

является иравнение Лапласа, но это доказательство проводится аналогично
тому, как это делалось в лекции 44 для функционала, зависящего от
функции одной переменной.

\begin{enumerate}
\def\labelenumi{\arabic{enumi}.}
\setcounter{enumi}{2}
\item
  В книгах по вариационному исчислению показано, как строить
  энергетический функционал не только для уравнения Пуассона с граничным
  условием Дирихле u=0, но и для многих других дифференциальных
  уравнений и многих других типов граничных условий. Следовательно,
  большое число краевых задач можно решить, минимизируя соответствующие
  энергетические функционалы.
\item
  Чем большее значение п мы выбираем, тем меньшее значение
  \(J[\bar{u}_n]\) получается. Значит, с увеличением n растет точность л
  приближенного решения дифференциального уравнения. Один из методов
  определения и, основан на последовательном вычислении функционала
  \(J[u_n]\) для возрастающих значений n.~Вычисления прекращаются, если
  с ростом п функционал перестает практически уменьшаться.
\item
  Если и велико, то проведение вычислений по методу Ритца требует
  привлечения ЭВМ. На рис. 45.2 (с. 345) приведена блок-
\end{enumerate}

схема решения краевой задачи методом Ритца.

Пользователь должен позаботиться о подпрограмме, которая должна
вычислять значения пробных функций. Такая подпрограмма может иметь вид:

SUBROUTINE BC (X.Y.PHI)

\begin{itemize}
\tightlist
\item
  SUBROUTINE PROVIDED BY THE USER TO C - EVALUATE THE - FUNCTIONS
  PHI(1), PHI(2), \ldots, PHI(N) DIMENSION PHI(20)
\end{itemize}

\[
PHI(1) = X\*Y\*(1-X)\*(1-Y)
\]

\(PHI(2) = X^*PHI(1)\) (используемые функции)

\[
PHI(3) = Y\*PHI(1)
\]

PHI(N) = (что получится)RETURN END

\section{\texorpdfstring{\textbf{ЗА}ДАЧИ}{ЗАДАЧИ}}\label{ux437ux430ux434ux430ux447ux438-1}

Найдите энергетический функционал для задачи

(УЧП)

Если теперь преобразовать начальные условия краевой задачи, то получим
начальные условия для обыкновенного дифференциального уравнения

, \(0 < x < 1\) , \(0 < y < 1\) , (ГУ) \(u = 0\) на границе.

\begin{enumerate}
\def\labelenumi{\arabic{enumi}.}
\setcounter{enumi}{1}
\tightlist
\item
  Как минимизировать функционал
\end{enumerate}

методом Ритца?

УКАЗАНИЕ. Если ввести новую функцию \(z(x)\):

то можно заметить, что она удовлетворяет граничным условиям \(z(0)\) = 0
и \(z(1)\) = 0.1

\begin{itemize}
\tightlist
\item
  \begin{enumerate}
  \def\labelenumi{\arabic{enumi}.}
  \setcounter{enumi}{2}
  \tightlist
  \item
    Напишите программу для проведения вычислений по блоксхеме,
    изображенной на рис. 45 2.
  \end{enumerate}
\item
  \begin{enumerate}
  \def\labelenumi{\arabic{enumi}.}
  \setcounter{enumi}{3}
  \tightlist
  \item
    Покажите, что уравнение
  \end{enumerate}
\end{itemize}

\[
\varphi(r) = a \ln r + b, \quad r = u^2 + v^2.
\]

является уравнением Эйлера --- Лагранжа для функционала

\[
\varphi(x, y) = 0.57 \ln(u^2 + v^2) + 1,
\]

\begin{enumerate}
\def\labelenumi{\arabic{enumi}.}
\setcounter{enumi}{4}
\tightlist
\item
  Задачу Дирихле
\end{enumerate}

Итак, решим теперь новое семейство задач Кощи:

(ОДУ)

где A=0.06 и B=0.04. Как найти потенциальную энергию этого решения? Мы
рекомендуем читателю вспомнить синуспреобразование и самостоятельно
получить решение \(u\left(x,y\right)\) .

*) При такой замене возникает особенность в т. x=1. Лучше ввести функцию
\(z(x)\) по формуле \(z(x)\)=\(y(x)\)=x. Прим. ред.

\pandocbounded{\includegraphics[keepaspectratio]{images/_page_343_Figure_2.jpeg}}

Далсе. Приближенное решение
\(u_n(x,y) = a_1 \varphi_1 + a_2 \varphi_2 + \dots + a_n \varphi_n\) . В
это время можно численно найти эначение функционала \(J[u_n]\) . Если
необходимо найти значение решения \(u_n(x,y)\) в некоторой точке, можно
написать небольшую программу эля быполнения отой операции. При
необходимости можно повторить вычисления по этой программе с другими
значениями п Рис. 452. Блок-схема метода Ритца.

\section{РЕШЕНИЕ УРАВНЕНИЙ С ЧАСТНЫМИ ПРОИЗВОДНЫМИ МЕТОДАМИ ТЕОРИИ
ВОЗМУЩЕНИЙ МЕТА ЛЕКЦІЇ: Показать, как различные сложные задачи
(нелинейные уравнения с переменными коэффициентами, области неправильной
формы и т. д.) можно решить методом возмущения более простых задач.
Иными словами, мы хотим показать, как следует изменить решение простой
задачи, чтобы оно давало приближенное решение трудной
задачи.}\label{ux440ux435ux448ux435ux43dux438ux435-ux443ux440ux430ux432ux43dux435ux43dux438ux439-ux441-ux447ux430ux441ux442ux43dux44bux43cux438-ux43fux440ux43eux438ux437ux432ux43eux434ux43dux44bux43cux438-ux43cux435ux442ux43eux434ux430ux43cux438-ux442ux435ux43eux440ux438ux438-ux432ux43eux437ux43cux443ux449ux435ux43dux438ux439-ux43cux435ux442ux430-ux43bux435ux43aux446ux456ux457-ux43fux43eux43aux430ux437ux430ux442ux44c-ux43aux430ux43a-ux440ux430ux437ux43bux438ux447ux43dux44bux435-ux441ux43bux43eux436ux43dux44bux435-ux437ux430ux434ux430ux447ux438-ux43dux435ux43bux438ux43dux435ux439ux43dux44bux435-ux443ux440ux430ux432ux43dux435ux43dux438ux44f-ux441-ux43fux435ux440ux435ux43cux435ux43dux43dux44bux43cux438-ux43aux43eux44dux444ux444ux438ux446ux438ux435ux43dux442ux430ux43cux438-ux43eux431ux43bux430ux441ux442ux438-ux43dux435ux43fux440ux430ux432ux438ux43bux44cux43dux43eux439-ux444ux43eux440ux43cux44b-ux438-ux442.-ux434.-ux43cux43eux436ux43dux43e-ux440ux435ux448ux438ux442ux44c-ux43cux435ux442ux43eux434ux43eux43c-ux432ux43eux437ux43cux443ux449ux435ux43dux438ux44f-ux431ux43eux43bux435ux435-ux43fux440ux43eux441ux442ux44bux445-ux437ux430ux434ux430ux447.-ux438ux43dux44bux43cux438-ux441ux43bux43eux432ux430ux43cux438-ux43cux44b-ux445ux43eux442ux438ux43c-ux43fux43eux43aux430ux437ux430ux442ux44c-ux43aux430ux43a-ux441ux43bux435ux434ux443ux435ux442-ux438ux437ux43cux435ux43dux438ux442ux44c-ux440ux435ux448ux435ux43dux438ux435-ux43fux440ux43eux441ux442ux43eux439-ux437ux430ux434ux430ux447ux438-ux447ux442ux43eux431ux44b-ux43eux43dux43e-ux434ux430ux432ux430ux43bux43e-ux43fux440ux438ux431ux43bux438ux436ux435ux43dux43dux43eux435-ux440ux435ux448ux435ux43dux438ux435-ux442ux440ux443ux434ux43dux43eux439-ux437ux430ux434ux430ux447ux438.}

Например, мы покажем, как нелинейную задачу

(УЧП)

\[
w = \ln \left[ \frac{z-1}{z+1} \right].
\]

, \(0 < r < 1\) ,\\
(ГУ) \(u(1, \theta) = \cos \theta\) , \(0 \le \theta \le 2\pi\) .

Очень часто задачу, решение которой не известно, можно путем
непрерывного изменения каких-то параметров свести к близкой, но легко
решаемой задаче. Тогда и решение исходной задачи будет близко решению
модифицирующей задачи.

Например, уравнение Лапласа

\[
\begin{cases} 1, & n = 1, \\ 0, & n = 2, 3, \dots, \end{cases}
\]

в помощью семейства уравнений

\[
,   0 \le \varepsilon \le 1  ,
\right.
,   0 < r < 1  ,   0 \le \theta < 2\pi  ,   u(1, \theta) = g(\theta)  ,   0 \le \theta < 2\pi  ,
\]

\begin{verbatim}
при одних граничных условиях (например,   g(\theta) = \cos \theta  ) имеет два различных вещественных решения, а при других (например,   g(\theta) = A \cos \theta   с A > 20,65) эта задача не имеет ни одного вещественного решения. — Прим. ред.
\end{verbatim}

4) Краевая задача

\begin{verbatim}
можно непрерывно модифицировать в нелинейное уравнение   \Delta u = u^2 = 0   1).
\end{verbatim}

На рис. 46.1 изображена схема решения нелинейной задачи путем сведения
ее к решению уравнения Лапласа2)

\pandocbounded{\includegraphics[keepaspectratio]{images/_page_362_Figure_1.jpeg}}

Рис. 46.1. Схема метода возмущении.

Чтобы найти решение возмущенного уравнения Лапласа

\[
u(0, t) = f(t),
u_{x}(0, t) = 0.
\]

\pandocbounded{\includegraphics[keepaspectratio]{images/_page_362_Figure_2.jpeg}}

\[
x + ct = \text{const},
\]

\pandocbounded{\includegraphics[keepaspectratio]{images/_page_362_Figure_5.jpeg}}

\[
0 < r < 1,     0 \le \theta < 2\pi,     u(1, 0) = 0,     0 \le \theta < 2\pi,
\]

нмеет два решения (одно --- нулсьое, другое --- нетривиальное), в то
время как задача

\pandocbounded{\includegraphics[keepaspectratio]{images/_page_362_Figure_7.jpeg}}

\[
u(1, \theta) = 0, \quad 0 \le \theta < 2\pi,
\]

\begin{verbatim}
\begin{array}{ccccccccccccccccccccccccccccccccccc
\end{verbatim}

\[
\begin{array}{ccc} (\text{УЧ}\Pi) & \Delta u = 0, & 0 < r < 1, \\ (\text{ГУ}) & u(1, \theta) = \cos \theta, & 0 \leqslant \theta \leqslant 2\pi, \end{array}
H(x-a) =
u_{tt} = c^2 u_{xx}, \quad 0 < x < L, \quad 0 < t < \infty,
\begin{cases} 0, x < a, \\ 1, x \geqslant a, \end{cases}
u(0, t) = g_1(t),
\]

В результате получаем

\[
u_x(0, t) = g_1(t),
H(x-a) =
u_x(0, t) - \gamma_1 u(0, t) = g_1(t),
\begin{cases} 0, & x < a, \\ 1, & x \geqslant a, \end{cases}
u(0, t) = g_1(t),
u_x(0, t) = g_1(t),
f(x) = \mathcal{F}^{-1}[F] = \frac{1}{\sqrt{2\pi}} \int_{-\infty}^{+\infty} F(\omega) e^{-i\omega x} d\omega \qquad F(\omega) = \mathcal{F}[f] = \frac{1}{\sqrt{2\pi}} \int_{-\infty}^{+\infty} f(x) e^{-i\omega x} dx
u_x(0, t) - \gamma_1 u(0, t) = g_1(t),
\begin{cases} \Delta u_{1} = -r^{2} \cos^{2} \theta = -\frac{r^{2}}{2} \left[1 + \cos(2\theta)\right] = -\frac{r^{2}}{2} - \frac{r^{2}}{2} \cos(2\theta), \\ u_{1}(1, \theta) = 0. \end{cases}
\begin{cases} 1, |x| < a \\ 0, |x| > a \end{cases}
\Delta u = -\frac{r^2}{2} - \frac{r^3}{2} \cos{(2\theta)}
\]

\section{ЗАМЕЧ
АНИЯ}\label{ux437ux430ux43cux435ux447-ux430ux43dux438ux44f}

\begin{enumerate}
\def\labelenumi{\arabic{enumi}.}
\tightlist
\item
  Для решения задачи методом конечного синус- или косинуспреобразования
  краевые условия при x=0 и x=L должны иметь вид
\end{enumerate}

\[
u_p(r, \theta) = Ar^4 + Br^4 \cos(2\theta).
\begin{cases} 1, |x| < 1 \\ 0, |x| > 1 \end{cases}
u_p(r, \theta) = -\frac{r^4}{32} - \frac{r^4}{24} \cos(2\theta).
\sum_{n=0}^{\infty} [a_n \cos(n\theta) + b_n \sin(n\theta)] - \frac{1}{32} - \frac{1}{24} \cos(2\theta) = 0.
\]

Отметим, что если величина \(u(0, t)\) положительна, то величина
\(u_x\)(0, t) также положительна, если \(u(L, t)\) положительна, то
\(u_x\)(L, t)

\pandocbounded{\includegraphics[keepaspectratio]{images/_page_144_Picture_4.jpeg}}

Рис. 19.6. Возникновение неоднородных граничных условий упругого
закрепления: h---коэффициент упругости пружины.

отрицательна. Однородные граничные условия (19.2) можно записать в виде

(19.3)

\[
P_{2}
\]

\(f(x)\) * \(g(x)\)

\[
\Delta u = 0, \quad 0 < r < 1 + \frac{1}{4} \sin \theta,
(1+x^2)^{-1}
\]

досягається в кінцевій точці сімейства граничних умов виду (див. рисунок
46 3)

\[
Tu_x(0, t) = -\beta u_t(0, t).
Tu_{x}(0, t) = \varphi[u(0, t)],
\frac{a}{x^2+a^2}
Tu_{x}(0, t) = -hu^{3}(0, t)
\frac{2ax}{(x^2+a^2)^2}
mu_{tt}(L, t) = -ku_x(L, t) + mg.
P_{0}
\]

\pandocbounded{\includegraphics[keepaspectratio]{images/_page_189_Picture_10.jpeg}}

Отже, розв'язуючи кожну з задач Дірікла всередині кола для функцій
\(u_0\) , \(u_1\) , \(u_2\) , отримуємо розв'язок

(46.9)

\begin{enumerate}
\def\labelenumi{\arabic{enumi}.}
\setcounter{enumi}{17}
\tightlist
\item
\end{enumerate}

Задачі Діріхле в деформованій області (зверніть увагу, що тут

\begin{verbatim}
Ми вже взяли  \epsilon = 1/4  :)
\end{verbatim}

Як вправу, читача просять знайти пертурбацію першого порядку \(u_1\) і
перевірити, наскільки добре є приблизне розв'язання

\[
\begin{cases} 1-|x|, |x| < 1 \\ 0, |x| > 1 \end{cases}
\]

якщо розглянути рівняння з параметром

(46.11)

і побачити його розв'язок у вигляді

Підставляючи цю декомпозицію на (46.11), отримуємо таку послідовність
завдань

\[
u_{x}(0, t) = \frac{h}{T} [u(0, t) - \theta_{1}(t)]
\frac{2}{\pi} \frac{\sin a\omega}{\omega}
u_{tt} = u_{xx}, \quad 0 < x < 1, \quad 0 < t < \infty,
\frac{2}{\pi} \frac{\sin a\omega}{(\omega^2+a^2)^2}
\Delta u = 0,
q_{uu} + q_{vv} = 0
w = f(z)
\]

\begin{verbatim}
\left\{
\end{verbatim}

\[
\Rightarrow.
\frac{\pi}{2} e^{-a+\omega}
\]

\begin{verbatim}
\begin{aligned} (Y & \Pi) & \varphi_{xx} + \varphi_{yy} = 0, & -\infty < x < \infty, & 0 < y < \infty, \\ (47.1) & \varphi(x, 0) = \begin{cases} 0, & |x| > 1, \\ 1, & |x| \le 1. \end{cases}
\end{verbatim}

\[
\frac{\pi}{2} e^{-a+\omega}
w = \ln\left\{\frac{z-1}{z+1}\right\}
\frac{2ax}{(x^2+a^2)^2}
\]

\begin{verbatim}
\begin{cases} (\text{УЧП}) & \varphi_{uu} + \varphi_{vv} = 0, \quad -\infty < u < \infty, \quad 0 < v < \pi, \\ (\text{47.2}) & \begin{cases} \varphi(u, 0) = 0, \\ \varphi(u, \pi) = 1. \end{cases}
\end{verbatim}

\[
\frac{\cos(ax), |x| > \pi/2a}{0, |x| > \pi/2a}
u(x, t) = \sum_{n=1}^{\infty} c_n X_n(x) T_n(t),
u_{tt} = \alpha^2 u_{xx}, \quad 0 < x < L, \quad 0 < t < \infty,
(\Gamma Y)
\]

Щоб візуально зрозуміти, як працює ця функція, ми радимо читачеві
побудувати її за кількома лініями y=c.

У другому прикладі ми перетворюємо область між двома неконцентричними
колами у кільце.

Задача Діріхле в області між двома некопоцентричними колами

Припустимо, ви хочете знайти потенціал між двома колами

\[
\cos(ax)
\quad 0 < t < \infty,
\frac{\pi}{2} [\delta(\omega+a) + \delta(\omega-a)]
\]

(HY)

\[
u = \gamma [(x-s)(x-t)-\gamma^2],
\quad 0 \le x \le L,
u(x, t) = X(x) T(t).
F(\omega) = \frac{?}{\pi} \int_{0}^{\infty} f(x) \sin(\omega x) dx
T'' - \alpha^2 \lambda T = 0,
X'' - \lambda X = 0,
u(x, t) = [C \sin(\beta x) + D \cos(\beta x)][A \sin(\alpha \beta t) + B \cos(\alpha \beta t)]
\beta_n = \frac{n\pi}{L}, \quad n = 0, 1, 2, \ldots
u_n(x, t) = X_n(x) T_n(t) = \sin(n\pi x/L) [a_n \sin(n\pi at/L) + b_n \cos(n\pi at/L)],
\]

f'\,'(0)

\[
u_n(x, t) = R_n \sin\left(\frac{n\pi x}{L}\right) \cos\left(\frac{n\pi\alpha (t - \delta_n)}{L}\right).
\frac{1}{\alpha} F\left(\frac{\omega}{\alpha}\right)
H(x-a) =
\frac{2\omega}{\pi (\alpha^2 + \omega^2)}
\]

Функція Хевісайду,

\[
[2/\pi\omega]^{1/2}
\begin{cases} 0, & x < a, \\ 1, & x \geqslant a, \end{cases}
\frac{2}{\pi\omega} [1-\cos(\omega a)]
f(x) = \mathcal{F}^{-1}[F] = \frac{1}{\sqrt{2\pi}} \int_{-\infty}^{+\infty} F(\omega) e^{-i\omega x} d\omega \qquad F(\omega) = \mathcal{F}[f] = \frac{1}{\sqrt{2\pi}} \int_{-\infty}^{+\infty} f(x) e^{-i\omega x} dx
e^{-1}a + \omega
\begin{cases} 1, |x| < a \\ 0, |x| > a \end{cases}
\frac{1}{2}e^{-\alpha}\sin\omega
\delta(x-a)
a_{n} = \frac{2}{n\pi\alpha} \int_{0}^{L} g(x) \sin(n\pi x/L) dx, \\ b_{n} = \frac{2}{L} \int_{0}^{L} f(x) \sin(n\pi x/L) dx.
u(x, t) = \sum_{n=1}^{\infty} \sin(n\pi x/L) \left[ a_n \sin(n\pi \alpha t/L) + b_n \cos(n\pi \alpha t/L) \right],
u(x, t) = \sum_{n=1}^{\infty} b_n \sin(n\pi x/L) \cos(n\pi\alpha t/L)
u(x, \theta) = f(x)
f(x) = \int_{0}^{\infty} F(\omega) \cos(\omega x) d\omega \qquad \qquad F(x) = \frac{2}{\pi} \int_{0}^{\infty} f(x) \cos(\omega x) dx
f(x) = \sum_{n=1}^{\infty} b_n \sin(n\pi x/L).
u_n(x, t) = b_n \sin(n\pi x/L) \cos(n\pi\alpha t/L).
u(x, t) = \sin(\pi x/L)\cos(\pi \alpha t/L) + 0.5\sin(3\pi x/L)\cos(3\pi \alpha t/L) + 0.25\sin(5\pi x/L)\cos(5\pi \alpha t/L).
R_n \sin(n\pi x/L) \cos[n\pi\alpha(t-\delta_n)/L],
\omega_n = \frac{n\pi\alpha}{L} = \frac{n\pi}{L} \sqrt{\frac{T}{\rho}},
S_n = \frac{2}{\pi} \int_0^{\pi} f(x) \sin(nx) dx.
u(x, 0) = \sin(\pi x/L) + 0.5 \sin(3\pi x/L),
y = \frac{\pi}{b - a} (x - a).
u(x, 0) = 0,
f(x) = \sum_{n=1}^{\infty} S_n \sin(nx).
u(x, 0) = \sin(3\pi x/L),
f(x) = \sum_{n=1}^{\infty} S_n \sin(nx)
u(x, 0) =
S_n = \frac{2}{\pi} \int_{0}^{\pi} f(x) \sin(nx) dx
\begin{cases} 2hx, & 0 \le x \le 0.5, \\ 2h(1-x), & 0.5 \le x \le 1. \end{cases}
0 \le x \le \pi
\]

К а к

\pandocbounded{\includegraphics[keepaspectratio]{images/_page_153_Picture_8.jpeg}}

ие движения будет совершать струна, если ее отпустить? 6. Решите задачу
о затухающих колебаниях струны

\[
n = 1, 2, ...
\frac{2ax}{(x^2+a^2)^2}
\]

f'\,'(x)

Представляется ли вам разумным полученное решение? Удовлетворяет ли оно
уравнению с частными производными, граничным и начальным условиям?

\begin{enumerate}
\def\labelenumi{\arabic{enumi}.}
\setcounter{enumi}{6}
\tightlist
\item
  Как вы будете решать неоднородное уравнение с заданными начальными и
  граничными условиям?
\end{enumerate}

\[
\frac{\pi}{2} e^{ax}
\cos(ax)
C_n = \frac{2}{\pi} \int_0^{\pi} f(x) \cos(nx) dx.
\frac{\pi}{2} [\delta(\omega+a) + \delta(\omega-a)]
y = \frac{\pi}{b-a}(x-a).
f(x) = \int_{0}^{\infty} F(\omega) \sin(\omega x) d\omega
f(x) = \frac{C_0}{2} + \sum_{n=1}^{\infty} C_n \cos(nx) \qquad C_n = \frac{2}{\pi} \int_{0}^{\pi} f(x) \cos nx \, dx
0 < x < \infty
\]

\begin{enumerate}
\def\labelenumi{\arabic{enumi}.}
\tightlist
\item
\end{enumerate}

\[
\frac{1}{\alpha} F\left(\frac{\omega}{\alpha}\right)
x^{-1/2}
\]

f'\,'(x)

\begin{enumerate}
\def\labelenumi{\arabic{enumi}.}
\setcounter{enumi}{4}
\tightlist
\item
  \(H(a-x)\)
\end{enumerate}

\[
\begin{array}{ll} u\left(x, \ 0\right) = \sin{(\pi x)}, \\ u_{t}\left(x, \ 0\right) = 0, \end{array}
\]

\begin{verbatim}
6.  x^{-1} 
\end{verbatim}

\begin{enumerate}
\def\labelenumi{\alph{enumi})}
\tightlist
\item
\end{enumerate}

\begin{verbatim}
7.  \frac{x}{x^2+a^2} 
\end{verbatim}

\begin{enumerate}
\def\labelenumi{\arabic{enumi}.}
\setcounter{enumi}{4}
\tightlist
\item
  \(\cos(mx), m=1, 2, \dots\)
\end{enumerate}

\begin{verbatim}
9.  \arctan \frac{a}{x} 
\end{verbatim}

\(\mathbf{\Pi}\) для яких рівнянь функція \(u_1 + u_2\) також буде
розв'язком? Які висновки ви можете зробити зі своїх відповідей?

\begin{enumerate}
\def\labelenumi{\arabic{enumi}.}
\setcounter{enumi}{3}
\tightlist
\item
  Найдите четыре смещанные задачи, сумма решений которых дает решение
  следующей задачи:
\end{enumerate}

Откуда получаем два уравнения

\[
0 < x < \infty \qquad \qquad 0 < \omega < \infty
\]

j'\,'(x)

\begin{enumerate}
\def\labelenumi{\arabic{enumi}.}
\setcounter{enumi}{4}
\tightlist
\item
  Решите задачу Коши
\end{enumerate}

(OAY)

\[
S_n = \frac{2}{\pi} \int_0^{\pi} f(x) \sin(nx) dx.
\]

(HY) \(U_n(0) = 0.\) Можете ли вы проверить полученное решение?

\begin{enumerate}
\def\labelenumi{\arabic{enumi}.}
\setcounter{enumi}{5}
\tightlist
\item
  Предположим, что функции \(u_1\) и \(u_2\) удовлетворяют линейным
  однородным граничным условиям
\end{enumerate}

\[
y = \frac{\pi}{b - a} (x - a).
\]

\begin{verbatim}
\begin{cases}  u_x(1, t) + h_2 u(1, t) = 0.  Будет ли функция   u_i + u_i   удовлетворять этим условиям? , .
\end{verbatim}

Лекция 27

\section{УРАВНЕНИЯ ПЕРВОГО ПОРЯДКА (МЕТОД
ХАРАКТЕРИСТИК)}\label{ux443ux440ux430ux432ux43dux435ux43dux438ux44f-ux43fux435ux440ux432ux43eux433ux43e-ux43fux43eux440ux44fux434ux43aux430-ux43cux435ux442ux43eux434-ux445ux430ux440ux430ux43aux442ux435ux440ux438ux441ux442ux438ux43a}

ЦЕЛЬ ЛЕКЦИИ: Ввести понятие уравнения с частными производными первого
порядка (до сих пор мы рассматривали уравнения второго порядка) и
познакомиться с важным методом решения задач с начальными
условиями---методом характеристик. Задача, которой мы теперь займемся,
имеет вид

(УЧП)

\[
f(x) = \sum_{n=1}^{\infty} S_n \sin(nx).
S_n = \frac{2}{\pi} \int_{0}^{\pi} f(x) \sin(nx) dx
\]

\begin{verbatim}
Сразу же отметим, что впервые мы будем решать задачу с переменными коэффициентами. Оказывается, если от координат (x, t) перейти к новым, характеристическим координатам   (s, \tau)  , то наше УЧП превратится в обыкновенное дифференциальное уравнение. Затем можно решить обыкновенное дифференциальное уравнение, т. е. найти   u(s, \tau)  , а на последнем шаге выразить s и t через x и t и получить u(x, t).
\end{verbatim}

Напомним читателю, что, когда мы решали уравнение дифф

Воспользовавшись ортогональностью семейства функции (sin (nnx)) на
отреже {[}0, 1{]}, находим формулы для коэффициентов

(21.6)

\begin{verbatim}
константа   \alpha^2   играла роль коэффициента диффузии, а v—была скоростью среды. Очевидно, что если   \alpha=0   (т. е. диффузии нет), то решение будет смещаться как целое вдоль оси x со скоростью v (у нас осталась только конвекция). Другими словами, если начальное условие   u(x,0)=\varphi(x)  , то соответствующее решение уравнения
\end{verbatim}

\[
\frac{\pi}{2} e^{ax}
\]

\begin{verbatim}
будет иметь вид   u(x, t) = \varphi(x-vt)  .
\end{verbatim}

Это дает нам основания думать, что решение уравнения первого порядка

\[
C_n = \frac{2}{\pi} \int_0^{\pi} f(x) \cos(nx) dx.
\]

переносится вдоль потока со скоростью

\[
y = \frac{\pi}{b-a}(x-a).
\]

Конечно, если a и b---константы, то решение представляет собой бегущую с
постоянной скоростью волну. Если же \(a(x, t)\) и \(b(x, t)\)

\[
f(x) = \frac{C_0}{2} + \sum_{n=1}^{\infty} C_n \cos(nx).
\]

\begin{verbatim}
РИС. 27.1. Начальное значение в точке x оказывает влияние на решение только вдоль характеристики в плоскости переменных (x,t). Параметр s изменяется вдоль характеристик от нуля до бесконечности, начиная с линии начального условия;   \tau   постоянна вдоль каждой характеристики;   \tau   изменяется вдоль линии начального условия.
\end{verbatim}

зависят от x и t, то скорость потока изменяется как вдоль потока, так и
во времени (читатель сможет увидеть, что начальная кривая сильно
искажается). Здесь так много аналогий с конвекцией!

Вернемся к нашей основной задаче:

(УЧП)

\[
0 \le x \le \pi \qquad n = 0, 1, 2, \dots
\]

\begin{verbatim}
(НУ)   u(x, 0) = \varphi(x), -\infty < x < \infty.
\end{verbatim}

Решение этого линейного уравнения первого порядка основывается на
следующем физическом факте: начальное условие в некоторой точке x
переносится в tx-плоскости вдоль линии, которая

называется характеристи-кой (см. рис. 27.1)

Этим наше уразнение отличается от других уравнений (таких, как уравнение
теплопроводности \(u_t = u_{xx} )\),для которых начальное значение в
точке x оказывает влияние на решение во всех точках пространства и во
все моменты времени. Если читатель помнит, то начальное сме- Р И С

f'\,'(x)

\begin{verbatim}
. 27.2. Начальное возмущение   u\left(x,0\right)   в точке порождает две волны. Начальное возмущение струны согласно волновому уравнению распространяется вдоль двух характеристик.
\end{verbatim}

\begin{verbatim}
щение скрипичной струны в точке x оказывает влияние на решение вдоль   \partial syx   линий в tx-плоскости (соответствующих двум бегущим волнам) (см. рис. 27.2).
\end{verbatim}

\begin{verbatim}
Естественно, напрашивается идея ввести две новые координаты s и   \tau   (вместо x и t) так, чтобы
\end{verbatim}

\begin{itemize}
\item
  s изменялась вдоль характеристических кривых,
\item
  т изменялась вдоль начальной кривой (желательно вдоль линии t=0).
\end{itemize}

Рассмотрим сначала новую координату s. Если s выбрать так, что она
удовлетворяет приведенному выше условию, то уравнение

\begin{enumerate}
\def\labelenumi{\arabic{enumi}.}
\setcounter{enumi}{1}
\tightlist
\item
  \(\frac{a_{0}}{2} + \sum_{n=1}^{\infty} a_{n} \cos(nx)\) \(a_{n}\) 3.
  \(f(n-x)\) \((-1)^{n} \frac{2}{\pi} C_{n}\) 4. 1
\end{enumerate}

превратится в обыкновенное дифференциальное уравнение К о н

\begin{enumerate}
\def\labelenumi{\arabic{enumi}.}
\setcounter{enumi}{4}
\tightlist
\item
  \(\cos(mx), m=1, 2, \dots\)
\end{enumerate}

\begin{verbatim}
ечно, остается открытым вопрос, как найти эти характеристики. Ответ прост: мы выберем характеристики   \{[x(s), t(s)]: 0 < s < \infty\}   так, чтобы они удовлетворяли системе
\end{verbatim}

\begin{verbatim}
\begin{cases}
\end{verbatim}

\end{document}
